\documentclass[UTF8, 12pt, fontset=fandol, oneside]{ctexart}
\usepackage{researchnotes}

% 保留分册独立编译时的章号前缀。
\renewcommand{\thesection}{14.\arabic{section}}

\title{第十四章\quad 多元微分学}
\author{}
\date{}

\begin{document}

\maketitle

在本章中, 我们将讨论多元函数的微分学. 将一元函数微分学的理论推广到多元函数 (向量) 函数是本章的中心内容. 除此之外, 由于多元函数与一元函数具有一些本质的差别, 在本章中我们还要研究一些只有在多元函数的情形才具有的问题, 如隐函数存在定理等。

\section{偏导数与全微分}

对于一个多元函数 $u=f(x_1,x_2,\cdots,x_n)$, 用一元微分学的工具来研究它是一个很自然的问题, 因此引进偏导数以及方向导数等概念是水到渠成的事情. 在本节中, 我们主要引进这些概念, 并讨论它们的简单性质。

\subsection{偏导数}

\noindent\textbf{定义 14.1.1}\quad 设函数 $u=f(x)=f(x_1,x_2,\cdots,x_n)$ 在区域 $D\subset\mathbb R^n$ 上有定义, $x_0=(x_1^0,x_2^0,\cdots,x_n^0)\in D$. 对于 $1\le i\le n$, 若一元函数
\[
  f(x_1^0,\cdots,x_{i-1}^0,x_i,x_{i+1}^0,\cdots,x_n^0)
\]
在 $x_i^0$ 处的导数, 即
\[
  \lim_{x_i\to x_i^0}
  \frac{f(x_1^0,\cdots,x_{i-1}^0,x_i,x_{i+1}^0,\cdots,x_n^0)
  -f(x_1^0,x_2^0,\cdots,x_n^0)}{x_i-x_i^0}
\]
存在, 则称 $f(x)$ 在 $x_0$ 处关于 $x_i$ 可偏导, 并称上述极限为 $f(x)$ 在 $x_0$ 处关于 $x_i$ 的偏导数, 记为 $\dfrac{\partial f(x_0)}{\partial x_i}, f'_{x_i}(x_0), \left.\dfrac{\partial u}{\partial x_i}\right|_{x_0}$ 等。

容易看出, 当一个 $n$ 元函数 $f(x)$ 关于每个分量都可偏导时, $f(x)$ 具有 $n$ 个偏导数. 若 $f(x)$ 在 $D$ 内的每一点关于 $x_i$ 都可偏导, 则将其偏导数记为 $\dfrac{\partial f(x)}{\partial x_i}, f'_{x_i}(x), \dfrac{\partial u}{\partial x_i}$ 等. 此时, 对每个固定的分量 $x_i$, $\dfrac{\partial f(x)}{\partial x_i}$ 仍然是 $D$ 中的一个 $n$ 元函数。

特别地, 若 $z=f(x,y)$ 是一个二元函数, 则我们用 $\dfrac{\partial f(x,y)}{\partial x}$ (或 $f'_x(x,y), \dfrac{\partial z}{\partial x}$) 与 $\dfrac{\partial f(x,y)}{\partial y}$ (或 $f'_y(x,y), \dfrac{\partial z}{\partial y}$) 来记它的两个偏导数. 同理, 对三元函数 $u=f(x,y,z)$, 可以类似地给出相应记号: $\dfrac{\partial f(x,y,z)}{\partial x}, \dfrac{\partial f(x,y,z)}{\partial y}, \dfrac{\partial f(x,y,z)}{\partial z}$, 或 $f'_x(x,y,z), f'_y(x,y,z), f'_z(x,y,z)$, 或 $\dfrac{\partial u}{\partial x}, \dfrac{\partial u}{\partial y}, \dfrac{\partial u}{\partial z}$。

对一个二元函数 $z=f(x,y)((x,y)\in D)$, 我们可以讨论其在 $(x_0,y_0)\in D$ 处偏导数的几何意义. 当函数具有较好的性质时, 该函数的图像是 $\mathbb R^3$ 中的一块曲面. $y=y_0$ 是一个过点 $(x_0,y_0,0)$ 且平行于 $Ozx$ 平面的平面, 它与曲面 $z=f(x,y)((x,y)\in D)$ 的交线为
\[
  L:\begin{cases}
    x=x,\\
    y=y_0,\\
    z=f(x,y_0).
  \end{cases}
\]
因此, $\dfrac{\partial f(x_0,y_0)}{\partial x}$ 是平面 $y=y_0$ 内的曲线 $L$ 在 $x=x_0$ 处的切线斜率 $k_1$. 同样, $\dfrac{\partial f(x_0,y_0)}{\partial y}$ 是平面 $x=x_0$ 与曲面 $z=f(x,y)((x,y)\in D)$ 的交线在 $y=y_0$ 处的切线斜率 $k_2$ (在图 14.1.1 中, $k_1=\tan\alpha, k_2=\tan\beta$)。

读者应该注意的是, 函数 $z=f(x,y)$ 在 $(x_0,y_0)$ 处的两个偏导数存在与否只与曲面与 $x=x_0$ 及 $y=y_0$ 的交线性质有关. 因此, 一个二元函数在点 $(x_0,y_0)$ 处即使两个偏导数都存在, 该函数在点 $(x_0,y_0)$ 处也未必连续. 例如, 设函数
\[
  f(x,y)=
  \begin{cases}
    0, & xy=0,\\
    1, & xy\ne0,
  \end{cases}
\]
容易看出 $f'_x(0,0)=f'_y(0,0)=0$, 但 $f(x,y)$ 在 $(0,0)$ 处不连续。

下面我们举两个例子:

\noindent\textbf{例 14.1.1}\quad 设函数 $f(x,y)=\tan\dfrac{x^2}{y}$, 求 $\dfrac{\partial f(0,1)}{\partial x}$ 及 $\dfrac{\partial f(0,1)}{\partial y}$。

\noindent\textbf{解}\quad 由偏导数的定义有
\[
  \frac{\partial f(0,1)}{\partial x}
  =(\tan x^2)'\big|_{x=0}
  =2x\sec^2 x^2\big|_{x=0}=0,
\]
\[
  \frac{\partial f(0,1)}{\partial y}
  =(0)'\big|_{y=1}=0.
\]

\noindent\textbf{例 14.1.2}\quad 设函数 $f(x,y)=\arcsin\dfrac{x}{\sqrt{x^2+y^2}}$, 求 $f'_x(x,y)$ 及 $f'_y(x,y)$。

\noindent\textbf{解}\quad 对这种计算题, 我们在求 $f'_x(x,y)$ 时将 $y$ 看成常数即可. 因此, 我们有
\[
\begin{aligned}
  f'_x(x,y)
  &=
  \frac{1}{\sqrt{1-\dfrac{x^2}{x^2+y^2}}}\cdot
  \frac{1}{x^2+y^2}
  \left(\sqrt{x^2+y^2}-\frac{x^2}{\sqrt{x^2+y^2}}\right)\\
  &=\frac{1}{\sqrt{y^2}}\cdot\frac{y^2}{x^2+y^2}
  =\frac{y\operatorname{sgn} y}{x^2+y^2};
\end{aligned}
\]
同理,
\[
  f'_y(x,y)
  =
  \frac{1}{\sqrt{1-\dfrac{x^2}{x^2+y^2}}}
  \left(-\frac{xy}{(x^2+y^2)^{3/2}}\right)
  =\frac{-x\operatorname{sgn} y}{x^2+y^2},
\]
其中 $\operatorname{sgn} y$ 是符号函数。

\subsection{方向导数}

一个多元函数的方向导数是用来刻画定义域内从一点出发的射线上函数的变化情况的. 从一点出发的一条射线可以看成是一个方向, 而一个方向可由以原点为心的单位球面上的点所确定. 设 $n$ 维单位球面为 $x_1^2+x_2^2+\cdots+x_n^2=1$, 则每一个以原点为起点, 终点在该球面上的有向线段表示一个单位向量 $v$, 它可以记为 $v=(\cos\theta_1,\cos\theta_2,\cdots,\cos\theta_n)$, 其中 $\theta_i(1\le i\le n)$ 是该向量与 $x_i$ 轴正向的夹角, $\cos\theta_i(1\le i\le n)$ 也称为是 $v$ 的方向余弦. 特别地, 对于 $\mathbb R^2$ 内任一个单位向量 $v$, 记它与正实轴的夹角为 $\theta$, 则有 $v=(\cos\theta,\sin\theta)$. 有了这些准备, 我们可以给出以下定义。

\noindent\textbf{定义 14.1.2}\quad 设函数 $u=f(x)$ 在区域 $D\subset\mathbb R^n$ 上有定义, $x_0\in D$,
\[
  v=(\cos\theta_1,\cos\theta_2,\cdots,\cos\theta_n)
\]
为一方向. 如果极限
\[
  \lim_{t\to0+0}\frac{f(x_0+tv)-f(x_0)}{t}
\]
存在, 则称该极限为 $f(x)$ 在 $x_0$ 处沿方向 $v$ 的方向导数, 记为 $\dfrac{\partial f(x_0)}{\partial v}$ 或 $\left.\dfrac{\partial u}{\partial v}\right|_{x_0}$。

\noindent\textbf{注 1}\quad 在上述定义中, 若记 $x_0+tv=x$, 则
\[
  \frac{\partial f(x_0)}{\partial v}
  =\lim_{x\to x_0}\frac{f(x)-f(x_0)}{(x-x_0)v}.
\]
从上述定义可以看出, 一个函数在某点处的方向导数不依赖于坐标系的选取。

\noindent\textbf{注 2}\quad 设 $f(x)$ 在 $x_0=(x_1^0,x_2^0,\cdots,x_n^0)$ 处关于 $x_i$ 的偏导数 $\dfrac{\partial f(x_0)}{\partial x_i}(1\le i\le n)$ 存在并等于 $A$, 若记 $v_i$ 是第 $i$ 个分量为 $1$ 的单位向量, 则 $\dfrac{\partial f(x_0)}{\partial v_i}=A$. 但读者应该注意的是 $\dfrac{\partial f(x_0)}{\partial(-v_i)}=-A$。

\noindent\textbf{例 14.1.3}\quad 设函数 $f(x,y)=\sqrt{x^2+y^2}$, 试求 $f(x,y)$ 在 $(0,0)$ 处各个方向的方向导数。

\noindent\textbf{解}\quad 设 $v=(\cos\theta,\sin\theta)(0\le\theta<2\pi)$, 则
\[
\begin{aligned}
  \frac{\partial f(0,0)}{\partial v}
  &=\lim_{t\to0+0}
  \frac{f(t\cos\theta,t\sin\theta)-f(0,0)}{t}\\
  &=\lim_{t\to0+0}
  \frac{\sqrt{t^2\cos^2\theta+t^2\sin^2\theta}}{t}\\
  &=\lim_{t\to0+0}\frac{t}{t}=1.
\end{aligned}
\]

从注 2 我们可知, 在上例中, $f(x,y)$ 在 $(0,0)$ 处两个偏导数都不存在. 当然, 我们也可从偏导数的定义直接证明这一结果。

\noindent\textbf{例 14.1.4}\quad 设函数
\[
  f(x,y)=
  \begin{cases}
    1, & y=x^2,\ \text{且 } x\ne0,\\
    0, & \text{其他},
  \end{cases}
\]
即 $f(x,y)$ 在抛物线 $y=x^2$ 上除去原点外取值为 $1$, 其余各处取值为 $0$. 证明: 对任意方向 $v$, 有 $\dfrac{\partial f(0,0)}{\partial v}=0$, 但 $f(x,y)$ 在 $(0,0)$ 处不连续。

\noindent\textbf{证明}\quad 记 $\mathbf 0=(0,0)$, 对任意方向 $v$, 当 $t$ 很小时, $\mathbf 0+tv$ 与抛物线不交, 从而 $f(tv)=0=f(0,0)$, 因此 $\dfrac{\partial f(0,0)}{\partial v}=0$. 由于
\[
  \lim_{k\to\infty}f\left(\frac1k,\frac1{k^2}\right)=1\ne f(0,0),
\]
所以 $f(x,y)$ 在 $(0,0)$ 处不连续。

\subsection{全微分}

一元函数 $y=f(x)$ 的微分是函数关于自变量的一阶线性的近似, 由其几何意义, 实际上是用切线来局部近似曲线 $y=f(x)$. 对于一个二元函数, 由于它的图像一般是一块曲面, 因此需用切平面来局部近似它, 而这些都与多元函数的全微分有关. 对一个 $n$ 元函数, 我们有以下定义。

\noindent\textbf{定义 14.1.3}\quad 设函数 $f(x)=f(x_1,x_2,\cdots,x_n)$ 在区域 $D\subset\mathbb R^n$ 上有定义, 且 $x_0=(x_1^0,x_2^0,\cdots,x_n^0)\in D$. 记 $\Delta x=(\Delta x_1,\Delta x_2,\cdots,\Delta x_n)$, 并称它为自变量的全增量. 再设
\[
  x=x_0+\Delta x=(x_1+\Delta x_1,x_2+\Delta x_2,\cdots,x_n+\Delta x_n)\in D.
\]
若存在仅依赖于 $x_0$ 的常数 $A_i(i=1,2,\cdots,n)$, 使得
\[
  \Delta f(x_0)=f(x_0+\Delta x)-f(x_0)
  =\sum_{i=1}^n A_i\Delta x_i+o(|\Delta x|),\quad |\Delta x|\to0,
\]
则称 $f(x)$ 在 $x_0$ 处可微, 并称 $\sum\limits_{i=1}^n A_i\Delta x_i$ 为 $f(x)$ 在 $x_0$ 处的全微分, 记为 $\mathrm df(x_0)$, 即
\[
  \mathrm df(x_0)=\sum_{i=1}^n A_i\Delta x_i.
\]
若 $f(x)$ 在 $D$ 内每一点处均可微, 则称 $f(x)$ 在 $D$ 内是可微函数。

当 $x_i(1\le i\le n)$ 是自变量时, 定义 $\mathrm dx_i=\Delta x_i$. 因此 $f(x)$ 在 $x_0$ 处的全微分可以记为 $\mathrm df(x_0)=\sum\limits_{i=1}^n A_i\mathrm dx_i$。

从全微分的定义可立即推出, 当 $f(x)$ 在 $x_0$ 处可微时, 它的全微分是函数改变量的线性主要部分. 本书中涉及一个多元函数的微分时, 一般情形下总是指上述的全微分, 因此今后我们将全微分简称为微分。

\noindent\textbf{定理 14.1.1}\quad 设函数 $f(x)$ 在区域 $D\subset\mathbb R^n$ 上有定义, 在 $x_0=(x_1^0,\cdots,x_n^0)\in D$ 处可微, 记其微分为 $\mathrm df(x_0)=\sum\limits_{i=1}^n A_i\mathrm dx_i$, 则
\[
\begin{gathered}
  (1)\quad f(x)\text{ 在 }x_0\text{ 处连续};\\
  (2)\quad \text{对于 } i(1\le i\le n), f(x)\text{ 关于 }x_i\text{ 可偏导, 并且有 }
  \frac{\partial f(x_0)}{\partial x_i}=A_i.
\end{gathered}
\]

\noindent\textbf{证明}\quad (1) 记 $\Delta x=x-x_0$, 则 $x\to x_0$ 等价于 $\Delta x\to0$, 即它等价于: 对于 $\forall i(1\le i\le n)$, 有 $\Delta x_i\to0$. 由此我们有
\[
  \lim_{\Delta x\to0}[f(x_0+\Delta x)-f(x_0)]
  =
  \lim_{\Delta x\to0}\left[\sum_{i=1}^n A_i\Delta x_i+o(|\Delta x|)\right]=0,
\]
即 $f(x)$ 在 $x_0$ 处连续。

(2) 对任何固定的 $i(1\le i\le n)$, 当 $j=1,2,\cdots,n$ 且 $j\ne i$ 时, 令 $x_j=x_j^0$, 即 $\Delta x_j=0$. 此时我们有 $\Delta x=(0,\cdots,\Delta x_i,\cdots,0)$ 和 $|\Delta x|=|\Delta x_i|$, 因此有
\[
\begin{aligned}
&\lim_{\Delta x_i\to0}
  \frac{f(x_1^0,x_2^0,\cdots,x_{i-1}^0,x_i^0+\Delta x_i,x_{i+1}^0,\cdots,x_n^0)-f(x_0)}{\Delta x_i}\\
&\quad =
  \lim_{\Delta x_i\to0}\left(A_i+\frac{o(|\Delta x_i|)}{\Delta x_i}\right)=A_i,
\end{aligned}
\]
从而 $f(x)$ 关于 $x_i$ 可偏导, 并且有 $\dfrac{\partial f(x_0)}{\partial x_i}=A_i$. 证毕。

上述定理说明, 当函数 $f(x)$ 在 $x_0$ 处可微时, 必有
\[
  \mathrm df(x_0)=\sum_{i=1}^n\frac{\partial f(x_0)}{\partial x_i}\mathrm dx_i.
\]
特别地, 当 $z=f(x,y)$ 在 $(x_0,y_0)$ 处可微时, 我们有
\[
  \mathrm df(x_0,y_0)=\frac{\partial f(x_0,y_0)}{\partial x}\mathrm dx+
  \frac{\partial f(x_0,y_0)}{\partial y}\mathrm dy;
\]
而当 $u=f(x,y,z)$ 在 $(x_0,y_0,z_0)$ 处可微时, 则有
\[
  \mathrm df(x_0,y_0,z_0)=
  \frac{\partial f(x_0,y_0,z_0)}{\partial x}\mathrm dx+
  \frac{\partial f(x_0,y_0,z_0)}{\partial y}\mathrm dy+
  \frac{\partial f(x_0,y_0,z_0)}{\partial z}\mathrm dz.
\]

我们已经有了例子表明存在可偏导但不连续的多元函数, 因此这样的函数不可微. 这说明, 可偏导与可微在多元函数的情形不是等价的. 但是我们有以下的定理。

\noindent\textbf{定理 14.1.2}\quad 设函数 $f(x)$ 在区域 $D\subset\mathbb R^n$ 上有定义, $x_0=(x_1^0,x_2^0,\cdots,x_n^0)\in D$, 再设 $f(x)$ 在 $x_0$ 的邻域 $U(x_0,\delta_0)(\delta_0>0)$ 内存在各个偏导数, 并且这些偏导数在 $x_0$ 处连续, 则 $f(x)$ 在 $x_0$ 处可微。

\noindent\textbf{证明}\quad 我们对 $n$ 作归纳法来证明定理。

当 $n=1$ 时, 定理结论显然成立。

假设当 $n=k$ 时定理成立, 即
\[
\begin{aligned}
&f(x_1^0+\Delta x_1,x_2^0+\Delta x_2,\cdots,x_k^0+\Delta x_k)
  -f(x_1^0,x_2^0,\cdots,x_k^0)\\
&\quad =
  \sum_{i=1}^k
  \frac{\partial f(x_1^0,x_2^0,\cdots,x_k^0)}{\partial x_i}\Delta x_i
  +o(1)\left(\sqrt{\sum_{i=1}^k(\Delta x_i)^2}\right)
\end{aligned}
\]
\[
  \left(\sqrt{\sum_{i=1}^k(\Delta x_i)^2}\to0\right).
\]
当 $n=k+1$ 时,
\[
\begin{aligned}
&\Delta f(x_1^0,\cdots,x_k^0,x_{k+1}^0)\\
&=f(x_1^0+\Delta x_1,\cdots,x_k^0+\Delta x_k,x_{k+1}^0+\Delta x_{k+1})
  -f(x_1^0,\cdots,x_k^0,x_{k+1}^0)\\
&=\big[f(x_1^0+\Delta x_1,\cdots,x_k^0+\Delta x_k,x_{k+1}^0+\Delta x_{k+1})
  -f(x_1^0+\Delta x_1,\cdots,x_k^0+\Delta x_k,x_{k+1}^0)\big]\\
&\quad+\big[f(x_1^0+\Delta x_1,\cdots,x_k^0+\Delta x_k,x_{k+1}^0)
  -f(x_1^0,\cdots,x_k^0,x_{k+1}^0)\big].
\end{aligned}
\]
由一元函数的拉格朗日微分中值定理及 $\dfrac{\partial f(x)}{\partial x_{k+1}}$ 在 $x_0$ 处的连续性知
\[
\begin{aligned}
&f(x_1^0+\Delta x_1,\cdots,x_k^0+\Delta x_k,x_{k+1}^0+\Delta x_{k+1})
  -f(x_1^0+\Delta x_1,\cdots,x_k^0+\Delta x_k,x_{k+1}^0)\\
&=f'_{x_{k+1}}(x_1^0+\Delta x_1,\cdots,x_k^0+\Delta x_k,
  x_{k+1}^0+\theta\Delta x_{k+1})\Delta x_{k+1}\\
&=f'_{x_{k+1}}(x_1^0,\cdots,x_k^0,x_{k+1}^0)\Delta x_{k+1}
  +o(1)(|\Delta x_{k+1}|),
\end{aligned}
\tag{14.1.1}
\]
其中 $0<\theta<1$. 由归纳法假设有
\[
\begin{aligned}
&f(x_1^0+\Delta x_1,\cdots,x_k^0+\Delta x_k,x_{k+1}^0)
  -f(x_1^0,\cdots,x_k^0,x_{k+1}^0)\\
&=
\sum_{i=1}^k
\frac{\partial f(x_1^0,\cdots,x_k^0,x_{k+1}^0)}{\partial x_i}\Delta x_i
+o(1)\left(\sqrt{\sum_{i=1}^k(\Delta x_i)^2}\right)
\end{aligned}
\tag{14.1.2}
\]
\[
  \left(\sqrt{\sum_{i=1}^k(\Delta x_i)^2}\to0\right).
\]
注意到当 $\sqrt{\sum\limits_{i=1}^{k+1}(\Delta x_i)^2}\to0$ 时, 有
\[
  \frac{
  o(1)(|\Delta x_{k+1}|)
  +o(1)\left(\sqrt{\sum\limits_{i=1}^k(\Delta x_i)^2}\right)}
  {\sqrt{\sum\limits_{i=1}^{k+1}(\Delta x_i)^2}}\to0,
\]
因此由式 (14.1.1), (14.1.2) 即得
\[
\begin{aligned}
&\Delta f(x_1^0,\cdots,x_k^0,x_{k+1}^0)\\
&=
  f(x_1^0+\Delta x_1,\cdots,x_k^0+\Delta x_k,x_{k+1}^0+\Delta x_{k+1})
  -f(x_1^0,\cdots,x_k^0,x_{k+1}^0)\\
&=\sum_{i=1}^{k+1}
  \frac{\partial f(x_1^0,x_2^0,\cdots,x_k^0,x_{k+1}^0)}{\partial x_i}\Delta x_i
  +o(1)\left(\sqrt{\sum_{i=1}^{k+1}(\Delta x_i)^2}\right),
\end{aligned}
\]
即 $f(x_1,x_2,\cdots,x_{k+1})$ 在 $(x_1^0,x_2^0,\cdots,x_{k+1}^0)$ 处可微, 且
\[
  \mathrm df(x_1^0,x_2^0,\cdots,x_{k+1}^0)
  =\sum_{i=1}^{k+1}
  \frac{\partial f(x_1^0,x_2^0,\cdots,x_{k+1}^0)}{\partial x_i}\mathrm dx_i.
\]
这就证明了定理的结论对一切 $n\in\mathbb N$ 成立. 证毕。

\noindent\textbf{注}\quad 若函数 $f(x)$ 在区域 $D\subset\mathbb R^n$ 上关于自变量的各个分量都具有连续偏导数, 通常我们称 $f(x)$ 在 $D$ 内是 $C^1$ 的, 记为 $f(x)\in C^1(D)$, 此时我们也称 $f(x)$ 在 $D$ 内连续可微。

\noindent\textbf{例 14.1.5}\quad 证明函数
\[
  f(x,y)=
  \begin{cases}
    (x^2+y^2)\sin\dfrac1{x^2+y^2}, & x^2+y^2\ne0,\\
    0, & x^2+y^2=0
  \end{cases}
\]
在 $(0,0)$ 处可微, 但 $f'_x(x,y)$ 及 $f'_y(x,y)$ 在 $(0,0)$ 处不连续。

\noindent\textbf{证明}\quad 显然 $f'_x(0,0)=f'_y(0,0)=0$. 由于
\[
  \Delta f(0,0)=f(\Delta x,\Delta y)
  =[(\Delta x)^2+(\Delta y)^2]\sin\frac1{(\Delta x)^2+(\Delta y)^2}
  =o(1)\sqrt{(\Delta x)^2+(\Delta y)^2},
\]
\[
  \left(\sqrt{(\Delta x)^2+(\Delta y)^2}\to0\right),
\]
因此 $f(x,y)$ 在 $(0,0)$ 处可微. 当 $(x,y)\ne(0,0)$ 时, 有
\[
  f'_x(x,y)=2x\sin\frac1{x^2+y^2}-\frac{2x}{x^2+y^2}\cos\frac1{x^2+y^2}.
\]
由于 $(x_k,y_k)=\left(\dfrac1{\sqrt{2k\pi}},0\right)\to(0,0)$ 时, $f'_x(x_k,y_k)\to\infty$, 因此 $f'_x(x,y)$ 在 $(0,0)$ 处不连续。

由函数的对称性, 同理 $f'_y(x,y)$ 在 $(0,0)$ 处也不连续。

对于可微函数, 它的方向导数与偏导数密切相关. 对此我们有下面的结论。

\noindent\textbf{定理 14.1.3}\quad 设函数 $f(x)$ 在区域 $D\subset\mathbb R^n$ 上有定义, 且在 $x_0\in D$ 处可微, 则 $f(x)$ 在 $x_0$ 处沿方向 $v=(\cos\theta_1,\cos\theta_2,\cdots,\cos\theta_n)$ 的方向导数为
\[
  \frac{\partial f(x_0)}{\partial v}
  =\sum_{i=1}^n\frac{\partial f(x_0)}{\partial x_i}\cos\theta_i.
\]

\noindent\textbf{证明}\quad 由 $f(x)$ 的可微性及方向导数的定义, 我们有
\[
\begin{aligned}
  \lim_{t\to0+0}\frac{f(x_0+tv)-f(x_0)}{t}
  &=\lim_{t\to0+0}
  \frac{\sum\limits_{i=1}^n\dfrac{\partial f(x_0)}{\partial x_i}t\cos\theta_i+o(|t|)}{t}\\
  &=\sum_{i=1}^n\frac{\partial f(x_0)}{\partial x_i}\cos\theta_i.
\end{aligned}
\]
证毕。

在本小节的最后, 我们再举一个例子来说明多元函数微分的在近似计算中的应用。

\noindent\textbf{例 14.1.6}\quad 求 $2.99^2\times1.02^3$ 的近似值。

\noindent\textbf{解}\quad 取函数 $f(x,y)=x^2y^3$, 显然 $f(x,y)$ 的两个偏导数在 $\mathbb R^2$ 上连续, 从而它在 $\mathbb R^2$ 处处可微. 取 $(x_0,y_0)=(3,1), \Delta x=-0.01, \Delta y=0.02$, 则
\[
\begin{aligned}
  2.99^2\times1.02^3
  &=f(2.99,1.02)=f(3-0.01,1+0.02)\\
  &\approx f(3,1)+f'_x(3,1)\Delta x+f'_y(3,1)\Delta y\\
  &=3^2\times1^3+f'_x(3,1)\cdot(-0.01)+f'_y(3,1)\cdot(0.02)\\
  &=9+6\cdot(-0.01)+27\cdot(0.02)=9.48.
\end{aligned}
\]

\subsection{梯度}

与方向导数密切相关的概念是梯度. 在本节我们先介绍梯度的概念与基本的性质, 在本书后面的章节中我们还会多次讨论它的进一步的性质。

在上一小节中, 我们知道当 $f(x)$ 在 $x_0\in\mathbb R^n$ 处可微时, $f(x)$ 沿任意方向 $v=(\cos\theta_1,\cos\theta_2,\cdots,\cos\theta_n)$ 的方向导数为
\[
  \frac{\partial f(x_0)}{\partial v}=
  \sum_{i=1}^n\frac{\partial f(x_0)}{\partial x_i}\cos\theta_i.
\]
当 $f(x)$ 在 $x_0$ 处的 $n$ 个偏导数不全为零时, $f(x)$ 沿方向
\[
  v_0=
  \frac1{\sqrt{\sum\limits_{i=1}^n\left(\dfrac{\partial f(x_0)}{\partial x_i}\right)^2}}
  \left(\frac{\partial f(x_0)}{\partial x_1},\frac{\partial f(x_0)}{\partial x_2},\cdots,\frac{\partial f(x_0)}{\partial x_n}\right)
\]
的方向导数达到最大值. 因此向量
\[
  \left(\frac{\partial f(x_0)}{\partial x_1},\frac{\partial f(x_0)}{\partial x_2},\cdots,\frac{\partial f(x_0)}{\partial x_n}\right)
\]
是 $f(x)$ 在 $x_0$ 处方向导数达到最大的方向, 同时它的模就是该方向的方向导数. 由此我们引进下面的定义。

\noindent\textbf{定义 14.1.4}\quad 设函数 $f(x)$ 在 $x_0$ 处可微, 则称向量
\[
  \left(\frac{\partial f(x_0)}{\partial x_1},\frac{\partial f(x_0)}{\partial x_2},\cdots,\frac{\partial f(x_0)}{\partial x_n}\right)
\]
为 $f(x)$ 在 $x_0$ 处的梯度, 记为 $\operatorname{grad}f(x_0)$, 即
\[
  \operatorname{grad}f(x_0)=
  \left(\frac{\partial f(x_0)}{\partial x_1},\frac{\partial f(x_0)}{\partial x_2},\cdots,\frac{\partial f(x_0)}{\partial x_n}\right).
\]

设 $f,g$ 均是可微函数, 则从定义容易推出梯度有以下简单性质:
\[
\begin{aligned}
  &(1)\quad \operatorname{grad}C=0,\ C\text{ 为常数};\\
  &(2)\quad \text{对于 } \forall\alpha,\beta\in\mathbb R,\ \text{有 }
  \operatorname{grad}(\alpha f+\beta g)=\alpha\operatorname{grad}f+\beta\operatorname{grad}g;\\
  &(3)\quad \operatorname{grad}(f\cdot g)=f\cdot\operatorname{grad}g+g\cdot\operatorname{grad}f;\\
  &(4)\quad \operatorname{grad}\frac fg=\frac1{g^2}(g\cdot\operatorname{grad}f-f\cdot\operatorname{grad}g)\quad(g\ne0).
\end{aligned}
\]
从梯度的定义可以看出, 当 $f(x)$ 在 $x_0$ 处可微时, $f(x)$ 沿方向 $v$ 的方向导数可以简单地记成
\[
  \frac{\partial f(x_0)}{\partial v}=\operatorname{grad}f(x_0)\cdot v.
\]
读者易于给出方向导数与梯度的关系的几何意义。

设函数 $u=f(x,y,z)$ 在 $(x_0,y_0,z_0)$ 处可微, $\operatorname{grad}f(x_0,y_0,z_0)$ 不是零向量. 记 $H=|\operatorname{grad}f(x_0,y_0,z_0)|$, $l_v$ 为从 $(x_0,y_0,z_0)$ 出发由方向 $v$ 确定的射线, 则对于 $\forall h\in(-H,H)$, 所有使得 $\dfrac{\partial f}{\partial v}=h$ 的 $l_v$ 的集合构成了顶点在 $(x_0,y_0,z_0)$, 母线是 $l_v$ 的锥面. 当 $h$ 从 $-H$ 连续变为 $H$ 时, 其母线 $l_v$ 与 $\operatorname{grad}f(x_0,y_0,z_0)$ 的夹角从 $\pi$ 变到 $0$. 特别地, 当它们的夹角为 $\dfrac\pi2$ 时, $l_v$ 的全体构成了过 $(x_0,y_0,z_0)$ 且以 $\operatorname{grad}f(x_0,y_0,z_0)$ 为法线的平面。

\noindent\textbf{例 14.1.7}\quad 设 $k\ge2$ 为正整数, 函数 $f(x,y)$ 在极坐标 $(r,\theta)$ 下有表示式
\[
  f(x,y)=
  \begin{cases}
    r\sin k\theta, &(x,y)\ne(0,0),\\
    0, &(x,y)=(0,0).
  \end{cases}
\]
(1) 求 $f(x,y)$ 在 $(0,0)$ 处的方向导数; (2) $f(x,y)$ 在 $(0,0)$ 处是否连续? (3) $f(x,y)$ 在 $(0,0)$ 处是否可微?

\noindent\textbf{解}\quad (1) 对固定的 $\theta_0\in[0,2\pi)$, $\theta=\theta_0$ 即确定一个从原点出发的方向 $v_0=(\cos\theta_0,\sin\theta_0)$. 我们有
\[
  \frac{\partial f(0,0)}{\partial v_0}
  =\lim_{r\to0}\frac{r\sin k\theta_0-0}{r}
  =\sin k\theta_0.
\]
因此 $f(x,y)$ 在 $(0,0)$ 处沿方向 $(\cos\theta,\sin\theta)$ 的方向导数即为 $\sin k\theta$。

(2) 显然 $\lim\limits_{(x,y)\to(0,0)}f(x,y)=0$, 所以 $f(x,y)$ 在 $(0,0)$ 处连续。

(3) 当 $k\ge2$ 时, 由于 $\sin k\theta$ 至少在 $\theta=\dfrac{\pi}{2k}, \dfrac{5\pi}{2k}\in[0,2\pi)$ 处取最大值 $1$, 我们断言 $f(x,y)$ 在 $(0,0)$ 处不可微. 如若断言不真, 假设 $f(x,y)$ 在 $(0,0)$ 处可微, 则 $\operatorname{grad}f(0,0)$ 是非零向量, 因此它的方向导数只能在一个方向即梯度的方向达到最大值. 此矛盾证明了断言。

\noindent\textbf{注}\quad 当 $k=2$ 且 $r\ne0$ 时,
\[
  r\sin2\theta=\frac{2r^2\sin\theta\cos\theta}{r}
  =\frac{2xy}{\sqrt{x^2+y^2}}.
\]
由此我们可将 $f(x,y)$ 写成如下形式:
\[
  f(x,y)=
  \begin{cases}
    \dfrac{2xy}{\sqrt{x^2+y^2}}, &x^2+y^2\ne0,\\
    0, &x^2+y^2=0.
  \end{cases}
\]
同理, 当 $k=3$ 且 $r\ne0$ 时,
\[
  f(x,y)=
  \begin{cases}
    \dfrac{3x^2y-y^3}{x^2+y^2}, &x^2+y^2\ne0,\\
    0, &x^2+y^2=0.
  \end{cases}
\]
请读者自己对上述函数在 $(x,y)$ 坐标下重新解例 14.1.7。

\subsection{向量函数的导数与全微分}

在本小节中, 我们将定义向量函数的导数与全微分. 由于涉及向量函数, 今后我们有时用行向量, 但有时也用列向量来记同一个向量, 不过从上下文中, 读者可以清楚地知道它何时是行向量, 何时是列向量。

\noindent\textbf{定义 14.1.5}\quad 设向量函数 $u=f(x)=(f_1(x),f_2(x),\cdots,f_m(x))^\mathrm T$ 在区域 $D\subset\mathbb R^n$ 上有定义, $x_0\in D$, $\Delta x=(\Delta x_1,\cdots,\Delta x_n)^\mathrm T$ 为 $x$ 在 $x_0$ 处的全增量. 如果存在 $m\times n$ 矩阵
\[
  A=
  \begin{pmatrix}
    A_{11}&A_{12}&\cdots&A_{1n}\\
    A_{21}&A_{22}&\cdots&A_{2n}\\
    \cdots&\cdots&\cdots&\cdots\\
    A_{m1}&A_{m2}&\cdots&A_{mn}
  \end{pmatrix}
\]
使得当 $|\Delta x|\to0$ 时, 下式成立:
\begin{equation}
\begin{aligned}
  \Delta f(x_0)
  &=(\Delta f_1(x_0),\cdots,\Delta f_m(x_0))^\mathrm T\\
  &=A(\Delta x_1,\cdots,\Delta x_n)^\mathrm T
  +(\alpha_1(|\Delta x|),\cdots,\alpha_m(|\Delta x|))^\mathrm T,
\end{aligned}
\tag{14.1.3}
\end{equation}
其中 $A$ 中的元素仅依赖于 $x_0$ 而不依赖于 $\Delta x$, 对于 $\forall j(1\le j\le m)$, $\alpha_j(|\Delta x|)$ 依赖于 $\Delta x$, 并且满足 $\lim\limits_{|\Delta x|\to0}\dfrac{\alpha_j(|\Delta x|)}{|\Delta x|}=0$, 则称 $f(x)$ 在 $x_0$ 处可微或可导, 矩阵 $A$ 称为 $f(x)$ 在 $x_0$ 处的 Fréchet 导数 (简称导数), 记做 $f'(x_0)$ 或 $Df(x_0)$; $A\Delta x$ 称为 $f(x)$ 在 $x_0$ 处的全微分 (简称微分), 记做 $\mathrm df(x_0)$, 即
\[
  \mathrm df(x_0)=A\Delta x=f'(x_0)\Delta x=Df(x_0)\Delta x.
\]
在上述定义中, 若规定 $\mathrm dx=\Delta x$, 则有 $\mathrm df(x_0)=f'(x_0)\mathrm dx$。

式 (14.1.3) 用矩阵可以简写成
\[
  \Delta f(x_0)=f(x_0+\Delta x)-f(x_0)=A\Delta x+\alpha(|\Delta x|),
\]
其中 $\alpha(|\Delta x|)=(\alpha_1(|\Delta x|),\cdots,\alpha_m(|\Delta x|))^\mathrm T$ 满足
\[
  \lim_{|\Delta x|\to0}\frac{\alpha(|\Delta x|)}{|\Delta x|}=0.
\]

下面我们考虑如何判断一个向量函数的可微性, 以及当它可微时, 如何求它的导数. 下面的定理对此给出了回答。

\noindent\textbf{定理 14.1.4}\quad 设 $D$ 是 $\mathbb R^n$ 中的区域, $x_0\in D$, 向量函数 $f(x)=(f_1(x),\cdots,f_m(x))^\mathrm T$ 在 $D$ 上有定义, 则 $f(x)$ 在 $x_0$ 处可微的充分必要条件是对于 $\forall j(1\le j\le m)$, $f_j(x)$ 在 $x_0$ 处可微. 记
\[
  A=
  \begin{pmatrix}
    \dfrac{\partial f_1(x_0)}{\partial x_1}&\dfrac{\partial f_1(x_0)}{\partial x_2}&\cdots&\dfrac{\partial f_1(x_0)}{\partial x_n}\\
    \dfrac{\partial f_2(x_0)}{\partial x_1}&\dfrac{\partial f_2(x_0)}{\partial x_2}&\cdots&\dfrac{\partial f_2(x_0)}{\partial x_n}\\
    \cdots&\cdots&\cdots&\cdots\\
    \dfrac{\partial f_m(x_0)}{\partial x_1}&\dfrac{\partial f_m(x_0)}{\partial x_2}&\cdots&\dfrac{\partial f_m(x_0)}{\partial x_n}
  \end{pmatrix},
\]
则当 $f(x)$ 在 $x_0$ 处可微时, 有
\[
  \mathrm df(x_0)=A\,\mathrm dx\quad\text{或}\quad f'(x_0)=A.
\]

\noindent\textbf{证明}\quad \textbf{必要性}\quad 设 $f(x)$ 在 $x_0$ 处可微, 从而存在 $m\times n$ 矩阵 $A=(A_{ji})$ 使得当 $|\Delta x|\to0$ 时, 有
\[
\begin{aligned}
  \Delta f(x_0)
  &=(\Delta f_1(x_0),\cdots,\Delta f_m(x_0))^\mathrm T\\
  &=\left(\sum_{i=1}^n A_{1i}\Delta x_i,\cdots,\sum_{i=1}^n A_{mi}\Delta x_i\right)^\mathrm T
  +(\alpha_1(|\Delta x|),\cdots,\alpha_m(|\Delta x|))^\mathrm T,
\end{aligned}
\]
其中 $\alpha_j(|\Delta x|)$ 依赖于 $\Delta x$, 且 $\lim\limits_{|\Delta x|\to0}\dfrac{\alpha_j(|\Delta x|)}{|\Delta x|}=0(j=1,\cdots,m)$. 比较上式两边向量的分量, 对 $j=1,\cdots,m$, 当 $|\Delta x|\to0$ 时, 有
\[
  \Delta f_j(x_0)=\sum_{i=1}^n A_{ji}\Delta x_i+\alpha_j(|\Delta x|).
\]
由多元函数可微的定义知, $f_j(x)$ 在 $x_0$ 处可微, 并且
\[
  A_{ji}=\frac{\partial f_j(x_0)}{\partial x_i}
  \quad (i=1,\cdots,n;\ j=1,\cdots,m).
\]

\textbf{充分性}\quad 设对于 $\forall j(1\le j\le m)$, $f_j(x)$ 在 $x_0$ 处可微, 则有
\[
  \Delta f_j(x_0)=\sum_{i=1}^n\frac{\partial f_j(x_0)}{\partial x_i}\Delta x_i+\alpha_j(|\Delta x|),
\]
其中 $\alpha_j(|\Delta x|)$ 依赖于 $\Delta x$, 且 $\lim\limits_{|\Delta x|\to0}\dfrac{\alpha_j(|\Delta x|)}{|\Delta x|}=0$. 因此
\[
\begin{aligned}
  \Delta f(x_0)
  &=(\Delta f_1(x_0),\cdots,\Delta f_m(x_0))^\mathrm T\\
  &=
  \begin{pmatrix}
    \dfrac{\partial f_1(x_0)}{\partial x_1}&\cdots&\dfrac{\partial f_1(x_0)}{\partial x_n}\\
    \dfrac{\partial f_2(x_0)}{\partial x_1}&\cdots&\dfrac{\partial f_2(x_0)}{\partial x_n}\\
    \cdots&\cdots&\cdots\\
    \dfrac{\partial f_m(x_0)}{\partial x_1}&\cdots&\dfrac{\partial f_m(x_0)}{\partial x_n}
  \end{pmatrix}
  \begin{pmatrix}
    \Delta x_1\\
    \Delta x_2\\
    \vdots\\
    \Delta x_n
  \end{pmatrix}
  +
  \begin{pmatrix}
    \alpha_1(|\Delta x|)\\
    \alpha_2(|\Delta x|)\\
    \vdots\\
    \alpha_m(|\Delta x|)
  \end{pmatrix}\\
  &=A\Delta x+\alpha(|\Delta x|),
\end{aligned}
\]
其中 $\alpha(|\Delta x|)$ 满足 $\lim\limits_{|\Delta x|\to0}\dfrac{\alpha(|\Delta x|)}{|\Delta x|}=0$. 由定义知 $f(x)$ 在 $x_0$ 处可微, 且 $f'(x_0)=A$. 证毕。

当向量函数 $f(x)=(f_1,\cdots,f_m)^\mathrm T$ 中的每个分量函数在 $x_0$ 处均可微时, 矩阵 $A=f'(x_0)$ 称为 $f(x)$ 在 $x_0$ 处的雅可比 (Jacobi) 矩阵, 记为 $J_f(x_0)$. 特别地, 当 $f(x)$ 是 $n$ 维向量函数时, $A$ 是 $n\times n$ 矩阵, 此时 $J_f(x_0)$ 的行列式称为 $f(x)$ 在 $x_0$ 处的雅可比行列式, 记为
\[
  |f'(x_0)|\quad\text{或}\quad
  \left.\frac{\partial(f_1,f_2,\cdots,f_n)}{\partial(x_1,x_2,\cdots,x_n)}\right|_{x_0}.
\]
显然, 对于一个向量函数 $f(x)$, 若其分量函数在 $x_0$ 处具有各个偏导数时, 我们就可以形式地定义 $J_f(x_0)$, 但当 $f(x)$ 在 $x_0$ 处不可微, 仅仅存在所有的偏导数时, $J_f(x_0)$ 不能刻画 $f(x)$ 在 $x_0$ 附近的变化情况。

当 $f_j(x)(1\le j\le m)$ 的各个偏导数都在 $x_0$ 处连续时, $f_j(x)$ 在 $x_0$ 处可微, 因此 $f(x)$ 在 $x_0$ 处可微. 另外, 若对于 $\forall j(1\le j\le m)$, $f_j(x)$ 的各个偏导数在区域 $D$ 上连续, 我们称 $f(x)$ 在 $D$ 上是 $C^1$ 的, 记为 $f(x)\in C^1(D)$. 特别地, 我们称 $\mathbb R^n$ 中区域 $D$ 到 $\Omega$ 的变换 $y=f(x)$ 是 $C^1$ 的, 如果 $f(x)\in C^1(D)$, 并且 $f^{-1}(y)\in C^1(\Omega)$。

\noindent\textbf{注}\quad 利用向量函数导数的记号, 对一个多元函数 $f(x)$, 当它可微时, 我们有
\[
  f'(x)=\left(\frac{\partial f}{\partial x_1},\cdots,\frac{\partial f}{\partial x_n}\right)=\operatorname{grad}f(x).
\]

\noindent\textbf{例 14.1.8}\quad 设向量 $w=(r,\varphi,\theta)$, 向量函数
\[
  f(w)=(x(w),y(w),z(w))^\mathrm T
  =(r\sin\varphi\cos\theta,r\sin\varphi\sin\theta,r\cos\varphi)^\mathrm T,
\]
其中 $D=\{(r,\varphi,\theta):0<r<+\infty,0<\varphi<\pi,0<\theta<2\pi\}$; 试求 $f'(w)$ 以及 $f(w)$ 在 $w$ 处的雅可比行列式。

\noindent\textbf{解}\quad 由于 $x=r\sin\varphi\cos\theta$, $y=r\sin\varphi\sin\theta$, $z=r\cos\varphi$ 均在 $D$ 上具有连续偏导数, 因此 $f'(w)$ 存在. 由计算得
\[
\begin{aligned}
  f'(w)
  &=
  \begin{pmatrix}
    \dfrac{\partial x}{\partial r}&\dfrac{\partial x}{\partial\varphi}&\dfrac{\partial x}{\partial\theta}\\
    \dfrac{\partial y}{\partial r}&\dfrac{\partial y}{\partial\varphi}&\dfrac{\partial y}{\partial\theta}\\
    \dfrac{\partial z}{\partial r}&\dfrac{\partial z}{\partial\varphi}&\dfrac{\partial z}{\partial\theta}
  \end{pmatrix}\\
  &=
  \begin{pmatrix}
    \sin\varphi\cos\theta&r\cos\varphi\cos\theta&-r\sin\varphi\sin\theta\\
    \sin\varphi\sin\theta&r\cos\varphi\sin\theta&r\sin\varphi\cos\theta\\
    \cos\varphi&-r\sin\varphi&0
  \end{pmatrix}.
\end{aligned}
\]
对上述矩阵取行列式, 即得 $f(w)$ 在 $w=(r,\varphi,\theta)$ 处的雅可比行列式
\[
  \frac{\partial(x,y,z)}{\partial(r,\varphi,\theta)}=r^2\sin\varphi.
\]

\section{多元函数求导法}

\subsection{导数的四则运算}

在本小节中, 我们总假定涉及的多元 (向量) 函数存在各个偏导数, 然后来讨论多元函数求偏导数的一些公式. 读者应该注意的是, 一般情况下一个函数的导数是一个向量 (或矩阵)。

\noindent\textbf{定理 14.2.1}\quad 设函数 $f(x),g(x)$ 在区域 $D\subset\mathbb R^n$ 上可导, 则对于 $\forall x\in D$, 有
\[
\begin{aligned}
  &(1)\quad (f(x)\pm g(x))'=f'(x)\pm g'(x);\\
  &(2)\quad (f(x)g(x))'=f(x)g'(x)+g(x)f'(x);\\
  &(3)\quad \left(\frac{f(x)}{g(x)}\right)'=
  \frac{g(x)f'(x)-f(x)g'(x)}{g^2(x)}\quad (g(x)\ne0).
\end{aligned}
\]
如果 $f(x):\mathbb R^n\to\mathbb R$ 是一个 $n$ 元可微函数, $g(x):\mathbb R^n\to\mathbb R^m$ 是一个可微的 $m$ 维列向量函数, 则有
\[
  (4)\quad (f(x)g(x))'=f(x)g'(x)+g(x)f'(x).
\]

\noindent\textbf{证明}\quad 我们只证 (3). (1), (2) 和 (4) 的证明留给读者. 由定义有
\[
  \left(\frac{f(x)}{g(x)}\right)'=
  \left(
    \frac{\partial\left(\dfrac{f(x)}{g(x)}\right)}{\partial x_1},
    \frac{\partial\left(\dfrac{f(x)}{g(x)}\right)}{\partial x_2},
    \cdots,
    \frac{\partial\left(\dfrac{f(x)}{g(x)}\right)}{\partial x_n}
  \right).
\]
由于
\[
  \frac{\partial\left(\dfrac{f(x)}{g(x)}\right)}{\partial x_i}
  =
  \frac{g(x)\dfrac{\partial f(x)}{\partial x_i}-f(x)\dfrac{\partial g(x)}{\partial x_i}}{g^2(x)}
  \quad (i=1,2,\cdots,n),
\]
将其代入上式即知 (3) 成立. 证毕。

\noindent\textbf{注}\quad 注意到 (4) 中等式左边 $(f(x)g(x))'$ 是一个 $m\times n$ 矩阵, $g(x)$ 是一个 $m\times1$ 矩阵, 而 $f'(x)$ 是一个 $1\times n$ 矩阵, 因此 $g(x)f'(x)$ 是一个 $m\times n$ 矩阵. 显然, $f(x)g'(x)$ 为一个 $m\times n$ 矩阵. 至于 (4) 的证明, 读者只要将 (4) 中等式左边矩阵中的每个分量与其右边矩阵中相对应的分量比较即可。

\subsection{复合函数的求导法}

下面我们主要来考虑复合函数的求导法。

\noindent\textbf{定理 14.2.2}\quad 设函数 $f(u)=f(u_1,u_2,\cdots,u_m)$ 在区域 $\Omega\subset\mathbb R^m$ 上有定义, 并且在 $u_0=(u_1^0,u_2^0,\cdots,u_m^0)^\mathrm T\in\Omega$ 处可微, 再设
\[
  u=u(x)=(u_1(x),u_2(x),\cdots,u_m(x))^\mathrm T
\]
在区域 $D\subset\mathbb R^n$ 上有定义, 在 $x_0=(x_1^0,x_2^0,\cdots,x_n^0)\in D$ 处可微, 并且 $u_0=u(x_0)$, 则 $f(u(x))$ 在 $x_0$ 处可微, 并且
\[
  \mathrm df(u(x_0))=f'(u(x_0))u'(x_0)\mathrm dx.
\]

\noindent\textbf{证明}\quad 由 $u$ 在 $x_0$ 处可微, 我们有
\[
  \Delta u(x_0)=u(x_0+\Delta x)-u(x_0)
  =u'(x_0)\Delta x+\alpha(|\Delta x|),
\]
其中 $\alpha(|\Delta x|)$ 依赖于 $\Delta x$, 且满足
\[
  \frac{\alpha(|\Delta x|)}{|\Delta x|}\to0\quad(|\Delta x|\to0).
\]
再由 $f(u)$ 在 $u_0$ 处可微, 有
\[
  \Delta f(u_0)=f(u_0+\Delta u)-f(u_0)=f'(u_0)\Delta u+\beta(|\Delta u|),
\]
其中 $\beta(|\Delta u|)$ 依赖于 $\Delta u$, 且满足
\[
  \frac{\beta(|\Delta u|)}{|\Delta u|}\to0\quad(|\Delta u|\to0).
\]
我们规定 $\beta(0)=0$. 在这样的规定下, $\beta(|\Delta u|)$ 在 $\Delta u=0$ 时连续。

我们有
\[
\begin{aligned}
  \Delta f(u(x_0))
  &=f(u(x_0+\Delta x))-f(u(x_0))\\
  &=f'(u(x_0))(u'(x_0)\Delta x+\alpha(|\Delta x|))
  +\beta(|u'(x_0)\Delta x+\alpha(|\Delta x|)|)\\
  &=f'(u(x_0))u'(x_0)\Delta x+\gamma(|\Delta x|),
\end{aligned}
\]
其中
\[
\begin{aligned}
  \gamma(|\Delta x|)
  &=f'(u(x_0))\alpha(|\Delta x|)
  +\beta(|\Delta u(x_0)|)\\
  &=f'(u(x_0))\alpha(|\Delta x|)
  +\beta(|u'(x_0)\Delta x+\alpha(|\Delta x|)|).
\end{aligned}
\]
下面我们证明
\[
  \frac{\gamma(|\Delta x|)}{|\Delta x|}\to0\quad(|\Delta x|\to0).
\]
首先, 显然有
\begin{equation}
  \lim_{|\Delta x|\to0}
  \frac{|f'(u(x_0))\alpha(|\Delta x|)|}{|\Delta x|}
  \le
  \lim_{|\Delta x|\to0}
  \frac{|f'(u(x_0))|\cdot|\alpha(|\Delta x|)|}{|\Delta x|}=0.
\tag{14.2.1}
\end{equation}
其次, 注意到当 $|\Delta x|$ 很小时, 有
\[
\begin{aligned}
  \frac{|\Delta u(x_0)|}{|\Delta x|}
  &=\frac{|u'(x_0)\Delta x+\alpha(|\Delta x|)|}{|\Delta x|}\\
  &\le \frac1{|\Delta x|}\bigl(\|u'(x_0)\|\,|\Delta x|+|\alpha(|\Delta x|)|\bigr)
  \le \|u'(x_0)\|+1,
\end{aligned}
\]
其中 $\|u'(x_0)\|$ 是矩阵 $u'(x_0)$ 的范数. 因此当 $|\Delta x|\to0$ 时, 必有 $|\Delta u(x_0)|\to0$. 我们有
\[
  \frac{|\beta(|\Delta u(x_0)|)|}{|\Delta x|}
  =
  \begin{cases}
    \dfrac{|\beta(|\Delta u(x_0)|)|}{|\Delta u(x_0)|}
    \cdot\dfrac{|\Delta u(x_0)|}{|\Delta x|}, &|\Delta u(x_0)|\ne0,\\[1.2em]
    0, &|\Delta u(x_0)|=0,
  \end{cases}
\]
从而有
\begin{equation}
  \frac{|\beta(|\Delta u(x_0)|)|}{|\Delta x|}\to0
  \quad(|\Delta x|\to0).
\tag{14.2.2}
\end{equation}
结合式 (14.2.1), (14.2.2) 得
\[
  \frac{\gamma(|\Delta x|)}{|\Delta x|}\to0\quad(|\Delta x|\to0).
\]
我们最后有
\[
  \Delta f(u(x_0))=f'(u(x_0))u'(x_0)\Delta x+\gamma(|\Delta x|),
\]
其中 $\gamma(|\Delta x|)$ 依赖于 $\Delta x$ 且 $\dfrac{\gamma(|\Delta x|)}{|\Delta x|}\to0(|\Delta x|\to0)$. 由微分的定义知, $f(u(x))$ 在 $x_0$ 处可微, 并且
\[
  \mathrm df(u(x_0))=f'(u(x_0))u'(x_0)\mathrm dx.
\]
证毕。

从定理 14.2.2 我们可以推出以下结论。

\noindent\textbf{推论 1}\quad 设向量函数 $f(u)=(f_1(u),f_2(u),\cdots,f_l(u))^\mathrm T$ 在区域 $\Omega\subset\mathbb R^m$ 上可微, $u(x)=(u_1(x),u_2(x),\cdots,u_m(x))^\mathrm T$ 在区域 $D\subset\mathbb R^n$ 上可微, 且 $u(D)\subset\Omega$, 则 $g=f(u(x))$ 在 $D$ 上可微, 并且
\[
  \mathrm df(u(x))=f'(u(x))u'(x)\mathrm dx.
\]
顺便我们有
\[
  [f(u(x))]'=f'(u(x))u'(x).
\]

\noindent\textbf{证明}\quad 由定理 14.2.2 知 $\mathrm df_j(u(x))=f'_j(u(x))u'(x)\mathrm dx\ (j=1,2,\cdots,l)$, 从而
\[
  \mathrm df(u(x))=
  \begin{pmatrix}
    f'_1(u(x))u'(x)\mathrm dx\\
    f'_2(u(x))u'(x)\mathrm dx\\
    \vdots\\
    f'_l(u(x))u'(x)\mathrm dx
  \end{pmatrix}
  =f'(u(x))u'(x)\mathrm dx.
\]
证毕。

\noindent\textbf{注}\quad 读者应该注意到 $f'(u(x))u'(x)$ 是一个 $l\times n$ 矩阵。

\noindent\textbf{推论 2}\quad 设 $D$ 与 $\Omega$ 为 $\mathbb R^n$ 中的区域, $y=f(x)$ 是 $D$ 到 $\Omega$ 的 $C^1$ 变换, 则对于 $\forall x\in D$, 有
\begin{equation}
  (f^{-1})'(y)\cdot f'(x)=E,
\tag{14.2.3}
\end{equation}
其中 $y=f(x)$; 对于 $\forall y\in\Omega$, 有
\begin{equation}
  f'(x)\cdot(f^{-1})'(y)=E,
\tag{14.2.4}
\end{equation}
其中 $x=f^{-1}(y)$. 特别地, 当 $y=f(x)$ 时, 有
\begin{equation}
  (f^{-1})'(y)=[f'(x)]^{-1},
\tag{14.2.5}
\end{equation}
其中 $E$ 是 $n\times n$ 单位矩阵, $[f'(x)]^{-1}$ 为 $f'(x)$ 的逆矩阵。

\noindent\textbf{证明}\quad 由 $C^1$ 变换的定义知 $f$ 与 $f^{-1}$ 都是可微向量函数, 因此 $(f^{-1}\circ f)(x)=x$ 两边对 $x$ 求导数即得式 (14.2.3), 而 $(f\circ f^{-1})(y)=y$ 两边对 $y$ 求导数即得式 (14.2.4). 在式 (14.2.3) 的两边同时右乘 $[f'(x)]^{-1}$ 即得式 (14.2.5). 证毕。

\noindent\textbf{注}\quad 在推论 2 中, 如果我们在式 (14.2.4) 两边取行列式, 则有
\begin{equation}
  \frac{\partial(y_1,\cdots,y_n)}{\partial(x_1,\cdots,x_n)}
  \cdot
  \frac{\partial(x_1,\cdots,x_n)}{\partial(y_1,\cdots,y_n)}
  =1.
\tag{14.2.6}
\end{equation}
下面的推论 3 是多元复合函数求偏导数的基础。

\noindent\textbf{推论 3}\quad 设函数 $f(u)=f(u_1,u_2,\cdots,u_m)$ 在区域 $\Omega\subset\mathbb R^m$ 上有定义, 并且在 $u_0=(u_1^0,u_2^0,\cdots,u_m^0)^\mathrm T\in\Omega$ 处可微, 又设向量函数
\[
  u=u(x)=(u_1(x),u_2(x),\cdots,u_m(x))^\mathrm T
\]
在区域 $D\subset\mathbb R^n$ 上有定义, 在 $x_0=(x_1^0,x_2^0,\cdots,x_n^0)\in D$ 处可微, 并且 $u_0=u(x_0)$, 则对于 $\forall i(1\le i\le n)$, $f(u(x))$ 在 $x_0$ 处关于 $x_i$ 可偏导, 并且
\[
  \frac{\partial f(u(x_0))}{\partial x_i}
  =\sum_{j=1}^m\left(
  \frac{\partial f(u_0)}{\partial u_j}\cdot
  \frac{\partial u_j(x_0)}{\partial x_i}\right).
\]

\noindent\textbf{证明}\quad 由定理 14.1.4 的注有
\[
  \left.(f(u(x)))'\right|_{x=x_0}
  =\left(
  \frac{\partial f(u(x_0))}{\partial x_1},
  \frac{\partial f(u(x_0))}{\partial x_2},
  \cdots,
  \frac{\partial f(u(x_0))}{\partial x_n}
  \right).
\]
用矩阵形式表示 $f'(u(x_0))u'(x_0)$, 有
\[
\begin{aligned}
  f'(u(x_0))u'(x_0)
  &=
  \left(
  \frac{\partial f(u(x_0))}{\partial u_1},
  \frac{\partial f(u(x_0))}{\partial u_2},
  \cdots,
  \frac{\partial f(u(x_0))}{\partial u_m}
  \right)\\
  &\quad\cdot
  \begin{pmatrix}
    \dfrac{\partial u_1(x_0)}{\partial x_1}&\dfrac{\partial u_1(x_0)}{\partial x_2}&\cdots&\dfrac{\partial u_1(x_0)}{\partial x_n}\\
    \dfrac{\partial u_2(x_0)}{\partial x_1}&\dfrac{\partial u_2(x_0)}{\partial x_2}&\cdots&\dfrac{\partial u_2(x_0)}{\partial x_n}\\
    \vdots&\vdots&\ddots&\vdots\\
    \dfrac{\partial u_m(x_0)}{\partial x_1}&\dfrac{\partial u_m(x_0)}{\partial x_2}&\cdots&\dfrac{\partial u_m(x_0)}{\partial x_n}
  \end{pmatrix}\\
  &=
  \left(
  \sum_{j=1}^m\frac{\partial f(u_0)}{\partial u_j}\frac{\partial u_j(x_0)}{\partial x_1},
  \sum_{j=1}^m\frac{\partial f(u_0)}{\partial u_j}\frac{\partial u_j(x_0)}{\partial x_2},
  \cdots,
  \sum_{j=1}^m\frac{\partial f(u_0)}{\partial u_j}\frac{\partial u_j(x_0)}{\partial x_n}
  \right).
\end{aligned}
\]
由定理 14.2.2 有 $\left.(f(u(x)))'\right|_{x=x_0}=f'(u(x_0))u'(x_0)$, 再比较它们的分量即得推论 3. 证毕。

\noindent\textbf{注}\quad 利用定理 14.2.2 类似的证明方法可以证明, 若将推论 3 中的条件 ``$u(x)$ 在 $x_0$ 处可微'' 减弱成 ``$u(x)$ 在 $x_0$ 处存在各个偏导数'', 则推论 3 的结论仍然成立. 请读者自己对此加以证明 (见本章习题). 但是, 在推论 3 中, 条件 ``$f(u)$ 在 $u_0$ 处可微'' 不能减弱成 ``$f(u)$ 在 $u_0$ 处存在各个偏导数''. 很容易找到例子来说明这一点, 例如取
\[
  f(u,v)=
  \begin{cases}
    1, &uv\ne0,\\
    0, &uv=0,
  \end{cases}
  \quad (u,v)\in\mathbb R^2,
\]
而令
\[
  \begin{cases}
    u=t,\\
    v=t
  \end{cases}
  \quad(-\infty<t<+\infty),
\]
则
\[
  g(t)=f(t,t)=
  \begin{cases}
    1, &t\ne0,\\
    0, &t=0
  \end{cases}
\]
在 $t=0$ 处不连续, 从而不可导。

如同一元函数, 我们将推论 3 中给出的复合函数求导公式称为链锁法则. 为了今后应用方便, 下面我们列出常见的二元复合函数的求导公式. 设 $z=f(u,v), u=u(x,y), v=v(x,y)$ 都为可微函数, 则
\[
  \frac{\partial z}{\partial x}
  =\frac{\partial f}{\partial u}\cdot\frac{\partial u}{\partial x}
  +\frac{\partial f}{\partial v}\cdot\frac{\partial v}{\partial x},
\]
\[
  \frac{\partial z}{\partial y}
  =\frac{\partial f}{\partial u}\cdot\frac{\partial u}{\partial y}
  +\frac{\partial f}{\partial v}\cdot\frac{\partial v}{\partial y}.
\]
有时, 若给出的函数具有形式 $z=f(u(x,y),v(x,y))$, 我们可以用以下记号:
\[
  \frac{\partial z}{\partial x}=f'_1\frac{\partial u}{\partial x}+f'_2\frac{\partial v}{\partial x},
  \qquad
  \frac{\partial z}{\partial y}=f'_1\frac{\partial u}{\partial y}+f'_2\frac{\partial v}{\partial y},
\]
其中 $f'_i(i=1,2)$ 指的是 $f$ 对其第 $i$ 个变量求偏导数. 对于具有多个中间变量的复合函数, 我们可以引进类似的记号, 这对抽象形式给出的函数的求导运算提供了简洁的记号。

\noindent\textbf{例 14.2.1}\quad 设函数 $u=e^{x^2+y^2+z^2}; z=x^2\sin y$, 求 $\dfrac{\partial u}{\partial x}, \dfrac{\partial u}{\partial y}$。

\noindent\textbf{解}\quad
\[
\begin{aligned}
  \frac{\partial u}{\partial x}
  &=\frac{\partial e^{x^2+y^2+z^2}}{\partial x}\cdot1
  +\frac{\partial e^{x^2+y^2+z^2}}{\partial z}\cdot\frac{\partial z}{\partial x}\\
  &=2xe^{x^2+y^2+z^2}+e^{x^2+y^2+z^2}(2z)(2x\sin y)\\
  &=2x(1+2x^2\sin^2 y)e^{x^2+y^2+z^2},
\end{aligned}
\]
\[
\begin{aligned}
  \frac{\partial u}{\partial y}
  &=\frac{\partial e^{x^2+y^2+z^2}}{\partial y}\cdot1
  +\frac{\partial e^{x^2+y^2+z^2}}{\partial z}\cdot\frac{\partial z}{\partial y}\\
  &=2ye^{x^2+y^2+z^2}+e^{x^2+y^2+z^2}(2z)(x^2\cos y)\\
  &=(2y+x^4\sin2y)\cdot e^{x^2+y^2+z^2}.
\end{aligned}
\]

\noindent\textbf{例 14.2.2}\quad 设函数 $u=f\left(xy,\dfrac yx,yz\right)$, 并设 $f$ 是可微函数, 试求 $\dfrac{\partial u}{\partial x}, \dfrac{\partial u}{\partial y}$ 和 $\dfrac{\partial u}{\partial z}$。

\noindent\textbf{解}\quad 对于这种抽象形式给出的函数, 我们采用上述约定的记号来求偏导数:
\[
  \frac{\partial u}{\partial x}
  =f'_1\frac{\partial(xy)}{\partial x}+f'_2\frac{\partial(y/x)}{\partial x}+f'_3\cdot0
  =yf'_1-\frac{y}{x^2}f'_2,
\]
\[
  \frac{\partial u}{\partial y}
  =f'_1\frac{\partial(xy)}{\partial y}
  +f'_2\frac{\partial(y/x)}{\partial y}
  +f'_3\frac{\partial(yz)}{\partial y}
  =xf'_1+\frac1x f'_2+zf'_3,
\]
\[
  \frac{\partial u}{\partial z}
  =f'_1\cdot0+f'_2\cdot0+f'_3\frac{\partial(yz)}{\partial z}
  =yf'_3.
\]

\noindent\textbf{例 14.2.3}\quad 记向量 $v=(r,\theta_1,\cdots,\theta_{n-1})$, 并设向量函数
\[
  f(v)=
  \begin{pmatrix}
    x_1(v)\\
    x_2(v)\\
    \vdots\\
    x_{n-1}(v)\\
    x_n(v)
  \end{pmatrix}
  =
  \begin{pmatrix}
    r\cos\theta_1\\
    r\sin\theta_1\cos\theta_2\\
    \vdots\\
    r\sin\theta_1\sin\theta_2\cdots\cos\theta_{n-1}\\
    r\sin\theta_1\sin\theta_2\cdots\sin\theta_{n-1}
  \end{pmatrix},
\]
试求 $f'(v)$ 及其雅可比行列式。

\noindent\textbf{解}\quad 由定义有

\[
  f'(v)=
  \begin{pmatrix}
    \cos\theta_1&-r\sin\theta_1&\cdots&0\\
    \sin\theta_1\cos\theta_2&r\cos\theta_1\cos\theta_2&\cdots&0\\
    \vdots&\vdots&\ddots&\vdots\\
    \sin\theta_1\cdots\cos\theta_{n-1}&r\cos\theta_1\cdots\cos\theta_{n-1}&\cdots&-r\sin\theta_1\cdots\sin\theta_{n-1}\\
    \sin\theta_1\cdots\sin\theta_{n-1}&r\cos\theta_1\cdots\sin\theta_{n-1}&\cdots&r\sin\theta_1\cdots\cos\theta_{n-1}
  \end{pmatrix}.
\]
下面用归纳法来求 $f'(v)$ 的行列式 $|f'(v)|$. 当 $n=2$ 时, 有 $|f'(v)|=r$; 当 $n=3$ 时, 有 $|f'(v)|=r^2\sin\theta_1$. 我们猜测对于 $n\ge3$, 有
\[
  |f'(v)|=r^{n-1}\sin^{n-2}\theta_1\sin^{n-3}\theta_2\cdots\sin\theta_{n-2}.
\]
事实上, 已知 $n=3$ 时上式成立. 假设 $n=k-1$ 时上式成立, 现在我们来证明当 $n=k$ 时上式成立. 为此将变换 $f$ 拆成下列两个变换的复合:
\[
  f_1:\begin{cases}
    x_1=y_1,\\
    x_2=y_2\cos y_3,\\
    x_3=y_2\sin y_3\cos y_4,\\
    \cdots\\
    x_{n-1}=y_2\sin y_3\sin y_4\cdots\cos y_n,\\
    x_n=y_2\sin y_3\sin y_4\cdots\sin y_n,
  \end{cases}
  \qquad
  f_2:\begin{cases}
    y_1=r\cos\theta_1,\\
    y_2=r\sin\theta_1,\\
    y_3=\theta_2,\\
    \cdots\\
    y_{n-1}=\theta_{n-2},\\
    y_n=\theta_{n-1}.
  \end{cases}
\]
则有 $f=f_1\circ f_2$. 当 $n=k$ 时, 由归纳法假设得
\[
\begin{aligned}
  |f'(v)|
  &=
  \frac{\partial(x_1,x_2,\cdots,x_k)}
  {\partial(r,\theta_1,\cdots,\theta_{k-1})}\\
  &=
  \left.
  \frac{\partial(x_1,x_2,\cdots,x_k)}
  {\partial(y_1,y_2,\cdots,y_k)}
  \right|_{f_2(v)}
  \frac{\partial(y_1,y_2,\cdots,y_k)}
  {\partial(r,\theta_1,\cdots,\theta_{k-1})}\\
  &=
  \left.(y_2^{k-2}\sin^{k-3}y_3\sin^{k-4}y_4\cdots\sin y_{k-1})\right|_{f_2(v)}
  \cdot r\\
  &=r^{k-1}\sin^{k-2}\theta_1\sin^{k-3}\theta_2\cdots\sin\theta_{k-2}.
\end{aligned}
\]

\noindent\textbf{例 14.2.4}\quad 设 $f(x)=(f_1(x),f_2(x),\cdots,f_m(x))$ 和 $g(x)=(g_1(x),g_2(x),\cdots,g_m(x))$ 为 $D\subset\mathbb R^n$ 到 $\mathbb R^m$ 的可微向量函数, 证明:
\[
  [f(x)(g(x))^\mathrm T]'=f(x)[(g(x))^\mathrm T]'+g(x)[(f(x))^\mathrm T]'.
\]

\noindent\textbf{证明}\quad 令
\[
  F(x)=
  \begin{pmatrix}
    (f(x))^\mathrm T\\
    (g(x))^\mathrm T
  \end{pmatrix}
  =(f_1(x),\cdots,f_m(x),g_1(x),\cdots,g_m(x))^\mathrm T,
\]
则 $F(x)$ 是 $D\to\mathbb R^{2m}$ 的可微函数. 记 $y=(y_1,y_2,\cdots,y_{2m})^\mathrm T$, 并令
\[
  G(y)=\sum_{i=1}^m y_i y_{m+i},
\]
则 $G(y)$ 是 $\mathbb R^{2m}\to\mathbb R$ 的可微函数. 我们可以验证
\[
  f(x)(g(x))^\mathrm T=G(F(x)).
\]
由链锁法则则有
\[
\begin{aligned}
  [f(x)(g(x))^\mathrm T]'
  &=[G(F(x))]'\\
  &=(y_{m+1},\cdots,y_{2m},y_1,\cdots,y_m)\big|_{y=F(x)}
  \begin{pmatrix}
    [(f(x))^\mathrm T]'\\
    [(g(x))^\mathrm T]'
  \end{pmatrix}\\
  &=f(x)[(g(x))^\mathrm T]'+g(x)[(f(x))^\mathrm T]'.
\end{aligned}
\]

\noindent\textbf{例 14.2.5}\quad 设 $f(x,y)$ 在 $\mathbb R^2$ 上具有连续偏导数, 试用极坐标来表示
\[
  |f'(x,y)|^2=
  \left(\frac{\partial f}{\partial x}\right)^2+
  \left(\frac{\partial f}{\partial y}\right)^2
  =|\operatorname{grad}f(x,y)|^2.
\]

\noindent\textbf{解}\quad 由
\[
  \begin{cases}
    x=r\cos\theta,\\
    y=r\sin\theta
  \end{cases}
\]
得
\[
  \frac{\partial f}{\partial x}
  =\frac{\partial f}{\partial r}\cdot\frac{\partial r}{\partial x}
  +\frac{\partial f}{\partial\theta}\cdot\frac{\partial\theta}{\partial x},
  \quad
  \frac{\partial f}{\partial y}
  =\frac{\partial f}{\partial r}\cdot\frac{\partial r}{\partial y}
  +\frac{\partial f}{\partial\theta}\cdot\frac{\partial\theta}{\partial y}.
\]
由 $r=\sqrt{x^2+y^2}$ 及 $\theta=\arctan\dfrac yx$ (此处分象限讨论, 我们以第一象限为例求之) 得
\[
  \frac{\partial r}{\partial x}=\frac{x}{\sqrt{x^2+y^2}}=\cos\theta,\quad
  \frac{\partial r}{\partial y}=\frac{y}{\sqrt{x^2+y^2}}=\sin\theta,
\]
\[
  \frac{\partial\theta}{\partial x}=-\frac{y}{x^2+y^2}=-\frac{\sin\theta}{r},\quad
  \frac{\partial\theta}{\partial y}=\frac{x}{x^2+y^2}=\frac{\cos\theta}{r}.
\]
因此
\[
  \frac{\partial f}{\partial x}
  =\cos\theta\frac{\partial f}{\partial r}
  -\frac{\sin\theta}{r}\cdot\frac{\partial f}{\partial\theta},
  \quad
  \frac{\partial f}{\partial y}
  =\sin\theta\frac{\partial f}{\partial r}
  +\frac{\cos\theta}{r}\cdot\frac{\partial f}{\partial\theta};
\]
从而有
\[
  |\operatorname{grad}f(x,y)|^2
  =
  \left(\frac{\partial f}{\partial x}\right)^2+
  \left(\frac{\partial f}{\partial y}\right)^2
  =
  \left(\frac{\partial f}{\partial r}\right)^2+
  \frac1{r^2}\left(\frac{\partial f}{\partial\theta}\right)^2.
\]

\noindent\textbf{例 14.2.6}\quad 设函数 $u=f(x-\lambda t)$ ($\lambda$ 是常数), 其中 $f$ 是 $\mathbb R$ 中的一个可微函数 (即 $f(y)$ 是可微函数), 证明它满足微分方程
\begin{equation}
  \frac{\partial u}{\partial t}+\lambda\frac{\partial u}{\partial x}=0,
\tag{14.2.7}
\end{equation}
并证明方程 (14.2.7) 的解一定具有形式 $u=g(x-\lambda t)$, 其中 $g$ 是任一个 $\mathbb R$ 中的可微函数。

\noindent\textbf{证明}\quad 记 $y=x-\lambda t$, 则
\[
  \frac{\partial u}{\partial t}
  =\frac{\mathrm df}{\mathrm dy}\cdot\frac{\partial y}{\partial t}
  =f'(y)(-\lambda)=-\lambda f'(y),
  \quad
  \frac{\partial u}{\partial x}
  =\frac{\mathrm df}{\mathrm dy}\cdot\frac{\partial y}{\partial x}
  =\frac{\partial f}{\partial y}=f'(y),
\]
从而
\[
  \frac{\partial u}{\partial t}+\lambda\frac{\partial u}{\partial x}
  =-\lambda f'(y)+\lambda f'(y)=0.
\]
设可微函数 $u=u(x,t)$ 是方程 (14.2.7) 的解, 作变量替换
\[
  \begin{cases}
    x=y+\lambda z,\\
    z=t,
  \end{cases}
\]
则 $u(x,t)=u(y+\lambda z,z)\triangleq u_1(y,z)$. 由
\[
  \frac{\partial u}{\partial t}
  =\frac{\partial u_1}{\partial y}\cdot\frac{\partial y}{\partial t}
  +\frac{\partial u_1}{\partial z}\cdot\frac{\partial z}{\partial t}
  =-\lambda\frac{\partial u_1}{\partial y}
  +\frac{\partial u_1}{\partial z},
\]
\[
  \frac{\partial u}{\partial x}
  =\frac{\partial u_1}{\partial y}\cdot\frac{\partial y}{\partial x}
  +\frac{\partial u_1}{\partial z}\cdot\frac{\partial z}{\partial x}
  =\frac{\partial u_1}{\partial y},
\]
有
\[
  \frac{\partial u}{\partial t}+\lambda\frac{\partial u}{\partial x}
  =\frac{\partial u_1}{\partial z}=0.
\]
这说明 $u_1(y,z)$ 与 $z$ 无关, 即 $u_1=g(y)$, 亦即 $u=g(x-\lambda t)$, 其中 $g$ 为 $\mathbb R$ 中的任一个可微函数。

\subsection{高阶偏导数}

设 $f(x)=f(x_1,x_2,\cdots,x_n)$ 在区域 $D\subset\mathbb R^n$ 上具有各个偏导数. 由定义, 它的每个偏导数 $\dfrac{\partial f(x)}{\partial x_i}(i=1,2,\cdots,n)$ 是 $D$ 上的一个 $n$ 元函数. 若它们仍具有各个偏导数, 则称它们的偏导数为 $f(x)$ 的二阶偏导数. 类似地可以定义三阶以及更高阶的偏导数. 二阶及二阶以上的偏导数统称为高阶偏导数。

显然, 若一个函数 $f(x)$ 的各个偏导数都存在 $n$ 个偏导数, 则 $f(x)$ 具有 $n^2$ 个二阶偏导数. 当
\[
  \frac{\partial\left(\dfrac{\partial f(x)}{\partial x_i}\right)}{\partial x_k}
  \quad(1\le i,k\le n)
\]
存在时, 我们将其记为 $\dfrac{\partial^2 f(x)}{\partial x_k\partial x_i}$, $f''_{x_kx_i}(x)$ 或 $f''_{ki}(x)$。

\noindent\textbf{例 14.2.7}\quad 设函数
\[
  f(x,y)=
  \begin{cases}
    xy\dfrac{x^2-y^2}{x^2+y^2}, &x^2+y^2\ne0,\\
    0, &x^2+y^2=0,
  \end{cases}
\]
求 $f''_{xy}(0,0)$ 及 $f''_{yx}(0,0)$。

\noindent\textbf{解}\quad 由偏导数的定义得
\[
  f'_x(x,y)=
  \begin{cases}
    y\left[\dfrac{x^2-y^2}{x^2+y^2}+\dfrac{4x^2y^2}{(x^2+y^2)^2}\right], &x^2+y^2\ne0,\\
    0, &x^2+y^2=0,
  \end{cases}
\]
\[
  f'_y(x,y)=
  \begin{cases}
    x\left[\dfrac{x^2-y^2}{x^2+y^2}-\dfrac{4x^2y^2}{(x^2+y^2)^2}\right], &x^2+y^2\ne0,\\
    0, &x^2+y^2=0.
  \end{cases}
\]
特别地, 我们有
\[
  f'_x(0,y)=-y,\qquad f'_y(x,0)=x.
\]
由此得
\[
  f''_{yx}(0,0)=-1,\qquad f''_{xy}(0,0)=1.
\]

从上例可知, 一个函数的各个高阶偏导数都存在时, 若改变对自变量分量的求导顺序, 高阶偏导数的值也有可能改变. 但如果我们附加一些条件, 可使高阶偏导数的值对求导顺序无关. 事实上, 我们有下述结果。

\noindent\textbf{定理 14.2.3}\quad 设函数 $f(x)$ 在区域 $D\subset\mathbb R^n$ 上有定义, $x_0\in D$, 且对于 $1\le j<k\le n$, $f''_{kj}(x)$ 与 $f''_{jk}(x)$ 在 $U(x_0,\delta)(\delta>0)$ 内存在, 并且在 $x_0$ 处连续, 则有
\[
  f''_{kj}(x_0)=f''_{jk}(x_0).
\]

\noindent\textbf{证明}\quad 记 $f(x)=f(x_1,x_2,\cdots,x_n)$. 为了使记号简便, 我们不妨设 $j=1,k=2$, 并记 $x=(x,y,x')$, $x_0=(x_0,y_0,x'_0)$ 其中 $x'=(x_3,x_4,\cdots,x_n)$, $x'_0=(x_3^0,x_4^0,\cdots,x_n^0)$. 对于充分小的 $\Delta x,\Delta y$, 观察:
\[
\begin{aligned}
  I(\Delta x,\Delta y)
  &\triangleq
  \frac{f(x_0+\Delta x,y_0+\Delta y,x'_0)-f(x_0+\Delta x,y_0,x'_0)}{\Delta x\Delta y}\\
  &\quad-\frac{f(x_0,y_0+\Delta y,x'_0)-f(x_0)}{\Delta x\Delta y}.
\end{aligned}
\]
记
\[
  g(x)=f(x,y_0+\Delta y,x'_0)-f(x,y_0,x'_0),
  \qquad
  h(y)=f(x_0+\Delta x,y,x'_0)-f(x_0,y,x'_0).
\]
一方面, 由一元函数的微分中值定理得
\[
\begin{aligned}
  I(\Delta x,\Delta y)
  &=\frac1{\Delta x\Delta y}
  \{[f(x_0+\Delta x,y_0+\Delta y,x'_0)-f(x_0+\Delta x,y_0,x'_0)]\\
  &\quad-[f(x_0,y_0+\Delta y,x'_0)-f(x_0)]\}
\end{aligned}
\]
\[
\begin{aligned}
  &=\frac{g(x_0+\Delta x)-g(x_0)}{\Delta x\Delta y}
  =\frac{g'(x_0+\theta_1\Delta x)}{\Delta y}\\
  &=\frac{f'_x(x_0+\theta_1\Delta x,y_0+\Delta y,x'_0)
  -f'_x(x_0+\theta_1\Delta x,y_0,x'_0)}{\Delta y}\\
  &=f''_{yx}(x_0+\theta_1\Delta x,y_0+\theta_2\Delta y,x'_0),
\end{aligned}
\]
其中 $0<\theta_1,\theta_2<1$. 另一方面, 我们有
\[
\begin{aligned}
  I(\Delta x,\Delta y)
  &=\frac1{\Delta x\Delta y}
  \{[f(x_0+\Delta x,y_0+\Delta y,x'_0)-f(x_0,y_0+\Delta y,x'_0)]\\
  &\quad-[f(x_0+\Delta x,y_0,x'_0)-f(x_0)]\}\\
  &=\frac{h(y_0+\Delta y)-h(y_0)}{\Delta x\Delta y}
  =\frac{h'(y_0+\theta_3\Delta y)}{\Delta x}\\
  &=\frac{f'_y(x_0+\Delta x,y_0+\theta_3\Delta y,x'_0)
  -f'_y(x_0,y_0+\theta_3\Delta y,x'_0)}{\Delta x}\\
  &=f''_{xy}(x_0+\theta_4\Delta x,y_0+\theta_3\Delta y,x'_0),
\end{aligned}
\]
其中 $0<\theta_3,\theta_4<1$. 由此得
\[
  f''_{yx}(x_0+\theta_1\Delta x,y_0+\theta_2\Delta y,x'_0)
  =f''_{xy}(x_0+\theta_4\Delta x,y_0+\theta_3\Delta y,x'_0).
\]
在上式中令 $\Delta x\to0,\Delta y\to0$, 并利用 $f''_{yx}(x)$ 与 $f''_{xy}(x)$ 在 $x_0$ 处的连续性即得
\[
  f''_{yx}(x_0)=f''_{xy}(x_0).
\]
证毕。

在今后, 我们碰到的大部分多元函数都具有各阶连续偏导数, 因此, 我们可以不考虑对自变量分量的求导顺序. 例如, 当 $y=f(x)(x\in D\subset\mathbb R^n)$ 具有各个二阶连续偏导数时, $f(x)$ 在 $x_0\in D$ 处具有 $\dfrac{n^2+n}{2}$ 个不同的二阶偏导数。

\subsection{复合函数的高阶偏导数}

对于复合函数的高阶偏导数, 很难得出一般性的公式, 需要逐次求之. 例如, 设 $u=f(x,y),x=\varphi(s,t),y=\psi(s,t)$, 并假设所有涉及的函数均具有二阶连续偏导数, 则有
\[
  \frac{\partial u}{\partial s}
  =\frac{\partial f}{\partial x}\cdot\frac{\partial\varphi}{\partial s}
  +\frac{\partial f}{\partial y}\cdot\frac{\partial\psi}{\partial s},
\]
\[
\begin{aligned}
  \frac{\partial^2u}{\partial s^2}
  &=\frac{\partial^2 f}{\partial x^2}\left(\frac{\partial\varphi}{\partial s}\right)^2
  +\frac{\partial^2 f}{\partial x\partial y}\cdot
  \frac{\partial\varphi}{\partial s}\cdot\frac{\partial\psi}{\partial s}
  +\frac{\partial f}{\partial x}\cdot\frac{\partial^2\varphi}{\partial s^2}\\
  &\quad+\frac{\partial^2 f}{\partial x\partial y}\cdot
  \frac{\partial\varphi}{\partial s}\cdot\frac{\partial\psi}{\partial s}
  +\frac{\partial^2 f}{\partial y^2}\left(\frac{\partial\psi}{\partial s}\right)^2
  +\frac{\partial f}{\partial y}\cdot\frac{\partial^2\psi}{\partial s^2}\\
  &=\frac{\partial^2 f}{\partial x^2}\left(\frac{\partial\varphi}{\partial s}\right)^2
  +2\frac{\partial^2 f}{\partial x\partial y}\cdot
  \frac{\partial\varphi}{\partial s}\cdot\frac{\partial\psi}{\partial s}
  +\frac{\partial^2 f}{\partial y^2}\left(\frac{\partial\psi}{\partial s}\right)^2\\
  &\quad+\frac{\partial f}{\partial x}\cdot\frac{\partial^2\varphi}{\partial s^2}
  +\frac{\partial f}{\partial y}\cdot\frac{\partial^2\psi}{\partial s^2}.
\end{aligned}
\]

\noindent\textbf{例 14.2.8}\quad 设函数 $u=f\left(xy,\dfrac yx\right)$, 试求 $\dfrac{\partial^2u}{\partial x^2}$ 及 $\dfrac{\partial^2u}{\partial x\partial y}$ (假定其二阶偏导数均连续)。

\noindent\textbf{解}\quad 由
\[
  \frac{\partial u}{\partial x}=yf'_1-\frac y{x^2}f'_2
\]
得
\[
\begin{aligned}
  \frac{\partial^2u}{\partial x^2}
  &=y^2f''_{11}-\frac y{x^2}f''_{12}
  -\frac y{x^2}f''_{12}+\frac{y^2}{x^4}f''_{22}
  +\frac{2y}{x^3}f'_2\\
  &=y^2f''_{11}-\frac{2y}{x^2}f''_{12}
  +\frac{y^2}{x^4}f''_{22}+\frac{2y}{x^3}f'_2.
\end{aligned}
\]
\[
\begin{aligned}
  \frac{\partial^2u}{\partial x\partial y}
  &=f'_1+xyf''_{11}+\frac yx f''_{12}
  -\frac1{x^2}f'_2-\frac{xy}{x^2}f''_{12}
  -\frac y{x^3}f''_{22}\\
  &=xyf''_{11}-\frac y{x^3}f''_{22}+f'_1-\frac1{x^2}f'_2.
\end{aligned}
\]

\noindent\textbf{例 14.2.9}\quad 设函数 $u=f(r),r=\sqrt{x^2+y^2+z^2}$, 且假设 $u$ 满足拉普拉斯方程
\[
  \frac{\partial^2u}{\partial x^2}
  +\frac{\partial^2u}{\partial y^2}
  +\frac{\partial^2u}{\partial z^2}=0,
\]
求 $f(r)$ 的表达式。

\noindent\textbf{解}\quad 由
\[
  \frac{\partial u}{\partial x}=f'(r)\frac xr
\]
得
\[
  \frac{\partial^2u}{\partial x^2}
  =f''(r)\frac{x^2}{r^2}+f'(r)\frac{r-\dfrac{x^2}{r}}{r^2}.
\]
同理,
\[
  \frac{\partial^2u}{\partial y^2}
  =f''(r)\frac{y^2}{r^2}+f'(r)\frac{r-\dfrac{y^2}{r}}{r^2},
\]
\[
  \frac{\partial^2u}{\partial z^2}
  =f''(r)\frac{z^2}{r^2}+f'(r)\frac{r-\dfrac{z^2}{r}}{r^2}.
\]
由于 $u$ 满足拉普拉斯方程, 因此有
\[
\begin{aligned}
  \frac{\partial^2u}{\partial x^2}
  +\frac{\partial^2u}{\partial y^2}
  +\frac{\partial^2u}{\partial z^2}
  &=f''(r)\frac{x^2+y^2+z^2}{r^2}
  +f'(r)\frac{3r^2-(x^2+y^2+z^2)}{r^3}\\
  &=f''(r)+\frac2r f'(r)=0.
\end{aligned}
\]
由此推出
\[
  (r^2f'(r))'=r^2f''(r)+2rf'(r)
  =r^2\left(f''(r)+\frac2r f'(r)\right)=0.
\]
从上式我们进一步推出: 存在常数 $C'$, 使得
\[
  f'(r)=\frac{C'}{r^2}.
\]
对上式两边求不定积分得
\[
  f(r)=-\frac{C'}r+C_1=\frac C{\sqrt{x^2+y^2+z^2}}+C_1,
\]
其中 $C=-C'$, $C_1$ 为常数。

\subsection{一阶微分的形式不变性与高阶微分}

设 $D\subset\mathbb R^n$ 是区域, 函数 $f(x)$ 在 $x\in D$ 处可微, 即有
\[
  \mathrm df(x)=\sum_{i=1}^n\frac{\partial f(x)}{\partial x_i}\mathrm dx_i.
\]
若对于 $\forall i(1\le i\le n)$, $\dfrac{\partial f}{\partial x_i}$ 仍是可微函数, 则我们称
\[
  \sum_{i=1}^n\left(\sum_{k=1}^n
  \frac{\partial^2 f(x)}{\partial x_k\partial x_i}\mathrm dx_k\right)\mathrm dx_i
\]
为 $f(x)$ 的二阶微分, 并记为 $\mathrm d^2f(x)$, 即
\[
  \mathrm d^2f(x)=\sum_{i=1}^n\sum_{k=1}^n
  \frac{\partial^2 f(x)}{\partial x_k\partial x_i}\mathrm dx_k\mathrm dx_i.
\]
对于 $k\ge2(k\in\mathbb N)$, 我们可以归纳地给出多元函数 $f(x)$ 的 $k$ 阶微分
\[
  \mathrm d^kf(x)=\mathrm d(\mathrm d^{k-1}f(x)).
\]
若形式地记 $\mathrm df(x)=\left(\sum_{i=1}^n\mathrm dx_i\frac{\partial}{\partial x_i}\right)f(x)$, 当 $f(x)$ 的每个 $k$ 阶偏导数都连续时, 我们可以将 $\mathrm d^kf(x)$ 记为
\[
  \mathrm d^kf(x)=\left(\sum_{i=1}^n\mathrm dx_i\frac{\partial}{\partial x_i}\right)^kf(x).
\]
这样的记号对于高阶微分 (二阶及二阶以上的微分) 来说比较容易记忆. 特别地, 对于一个二元函数 $f(x,y)$, 若它各个 $k$ 阶偏导数都存在且连续, 则我们有
\[
\begin{aligned}
  \mathrm d^kf(x,y)
  &=\left(\mathrm dx\frac{\partial}{\partial x}
  +\mathrm dy\frac{\partial}{\partial y}\right)^kf(x,y)\\
  &=\sum_{j=0}^k C_k^j
  \frac{\partial^k f(x,y)}{\partial x^{k-j}\partial y^j}
  \mathrm dx^{k-j}\mathrm dy^j,
\end{aligned}
\]
其中 $C_k^j=\dfrac{k!}{j!(k-j)!}(j=1,2,\cdots,k),C_k^0=1$。

现设 $f(u)=f(u_1,u_2,\cdots,u_m)$ 在区域 $D\subset\mathbb R^m$ 上可微, 则 $f(u)$ 的微分 $\mathrm df(u)=f'(u)\mathrm du$. 注意到此时 $u$ 是自变量。

现在设 $u=(u_1(x),u_2(x),\cdots,u_m(x))^\mathrm T$ 是区域 $\Omega\subset\mathbb R^n$ 上的一个 $n$ 元可微向量函数, 并且 $u(\Omega)\subset D$, 则复合函数 $f(u(x))$ 在 $\Omega$ 处可微, 从而有
\[
  \mathrm df(u(x))=f'(u(x))u'(x)\mathrm dx.
\]
由于 $\mathrm du=(u'_1(x)\mathrm dx,u'_2(x)\mathrm dx,\cdots,u'_m(x)\mathrm dx)^\mathrm T=u'(x)\mathrm dx$, 因此, 当 $u$ 是中间变量时, 我们仍然有
\[
  \mathrm df(u)=f'(u(x))u'(x)\mathrm dx=f'(u)\mathrm du.
\]
以上讨论说明了对于多元函数的一阶微分仍具有形式不变性. 在一元微积分中, 我们已经知道了对于二阶以上的微分不再具有形式不变性, 因此多元函数情形也不可能具有高阶微分的形式不变性。

\noindent\textbf{例 14.2.10}\quad 设函数 $f\left(xy,\dfrac yx\right)$ 具有二阶连续偏导数, 求它的所有二阶偏导数。

\noindent\textbf{解}\quad 在例 14.2.8 中我们曾经直接对该函数求过二阶偏函数, 现在我们利用 $f$ 的二阶微分来求之. 由
\[
  \mathrm df=f'_1\mathrm d(xy)+f'_2\mathrm d\left(\frac yx\right)
\]
得
\[
\begin{aligned}
  \mathrm d^2f
  &=\left(f''_{11}\mathrm d(xy)+f''_{12}\mathrm d\left(\frac yx\right)\right)\mathrm d(xy)
  +f'_1\mathrm d^2(xy)\\
  &\quad+\left(f''_{21}\mathrm d(xy)+f''_{22}\mathrm d\left(\frac yx\right)\right)
  \mathrm d\left(\frac yx\right)+f'_2\mathrm d^2\left(\frac yx\right)\\
  &=\left(y^2f''_{11}-\frac{2y^2}{x^2}f''_{12}
  +\frac{y^2}{x^4}f''_{22}+\frac{2y}{x^3}f'_2\right)\mathrm dx^2\\
  &\quad+2\left(xyf''_{11}-\frac y{x^3}f''_{22}
  +f'_1-\frac1{x^2}f'_2\right)\mathrm dx\mathrm dy\\
  &\quad+\left(x^2f''_{11}+2f''_{12}+\frac1{x^2}f''_{22}\right)\mathrm dy^2.
\end{aligned}
\]
因此
\[
  \frac{\partial^2 f}{\partial x^2}
  =y^2f''_{11}-\frac{2y^2}{x^2}f''_{12}
  +\frac{y^2}{x^4}f''_{22}+\frac{2y}{x^3}f'_2,
\]
\[
  \frac{\partial^2 f}{\partial x\partial y}
  =xyf''_{11}-\frac y{x^3}f''_{22}
  +f'_1-\frac1{x^2}f'_2,
\]
\[
  \frac{\partial^2 f}{\partial y^2}
  =x^2f''_{11}+2f''_{12}+\frac1{x^2}f''_{22}.
\]

\section{泰勒公式}

在利用高阶导数研究函数时, 泰勒公式是一个强有力的工具, 对于多元函数我们有以下定理。

\noindent\textbf{定理 14.3.1(泰勒公式)}\quad 设函数 $f(x)$ 在 $x_0=(x_1^0,x_2^0,\cdots,x_n^0)\in\mathbb R^n$ 的邻域 $U(x_0,\delta_0)(\delta_0>0)$ 内具有 $K+1$ 阶连续偏导数, 则对于 $\forall x_0+h=(x_1^0+h_1,x_2^0+h_2,\cdots,x_n^0+h_n)\in U(x_0,\delta_0)$, 有
\begin{equation}
\begin{aligned}
  f(x_0+h)
  &=f(x_0)+\sum_{k=1}^K\frac1{k!}
  \left(\sum_{i=1}^n h_i\frac{\partial}{\partial x_i}\right)^k f(x_0)\\
  &\quad+\frac1{(K+1)!}
  \left(\sum_{i=1}^n h_i\frac{\partial}{\partial x_i}\right)^{K+1}
  f(x_0+\theta h),
\end{aligned}
\tag{14.3.1}
\end{equation}
其中 $0<\theta<1$。

\noindent\textbf{证明}\quad 设 $h=(h_1,h_2,\cdots,h_n)$ 满足 $x_0+h\in U(x_0,\delta_0)$. 构造一元函数
\[
  \varphi(t)=f(x_0+th),\qquad t\in[0,1].
\]
由复合函数求导公式知, $\varphi(t)$ 具有 $K+1$ 阶连续导数. 因此, 由一元函数的泰勒公式有
\begin{equation}
  \varphi(t)=\varphi(0)+\sum_{k=1}^K
  \frac{\varphi^{(k)}(0)}{k!}t^k
  +\frac{\varphi^{(K+1)}(\theta t)}{(K+1)!}t^{K+1},
  \tag{14.3.2}
\end{equation}
其中 $0<\theta<1$. 注意到
\[
  \varphi'(t)=\sum_{i=1}^n\frac{\partial f(x_0+th)}{\partial x_i}h_i
  =\left(\sum_{i=1}^n h_i\frac{\partial}{\partial x_i}\right)f(x_0+th),
\]
\[
  \varphi''(t)=\sum_{j=1}^n\sum_{i=1}^n
  \frac{\partial^2 f(x_0+th)}{\partial x_i\partial x_j}h_ih_j
  =\left(\sum_{i=1}^n h_i\frac{\partial}{\partial x_i}\right)^2 f(x_0+th),
\]
一般地, 对于 $k=1,2,\cdots,K+1$, 有
\[
  \varphi^{(k)}(t)
  =\left(\sum_{i=1}^n h_i\frac{\partial}{\partial x_i}\right)^k f(x_0+th).
\]
因此, 将 $\varphi^{(k)}(t)(k=1,2,\cdots,K+1)$ 的表达式代入式 (14.3.2), 然后令 $t=1$ 即得所证. 证毕。

\noindent\textbf{注}\quad 当 $f(x)$ 在区域 $D$ 内具有 $K+1$ 阶连续偏导数时, 定理 14.3.1 中的余项
\[
  R_{K+1}(h)=\frac1{(K+1)!}
  \left(\sum_{i=1}^n h_i\frac{\partial}{\partial x_i}\right)^{K+1}f(x_0+\theta h)
  \quad(0<\theta<1)
\]
称为拉格朗日余项, 并称式 (14.3.1) 为 $f(x)$ 带拉格朗日余项的泰勒公式。

\noindent\textbf{推论 1}\quad 设函数 $f(x)\in C^K(U(x_0,\delta_0))(\delta_0>0)$, 即 $f(x)$ 在邻域 $U(x_0,\delta_0)$ 内具有 $K(K\ge1)$ 阶连续偏导数, $x_0+h\in U(x_0,\delta_0)$, 则
\begin{equation}
  f(x_0+h)=f(x_0)+\sum_{k=1}^K\frac1{k!}
  \left(\sum_{i=1}^n h_i\frac{\partial}{\partial x_i}\right)^k f(x_0)
  +o(|h|^K)
  \quad(|h|\to0).
  \tag{14.3.3}
\end{equation}

\noindent\textbf{证明}\quad 由定理 14.3.1 我们有
\[
\begin{aligned}
  f(x_0+h)
  &=f(x_0)+\sum_{k=1}^K\frac1{k!}
  \left(\sum_{i=1}^n h_i\frac{\partial}{\partial x_i}\right)^kf(x_0)\\
  &\quad+\frac1{K!}
  \left(\sum_{i=1}^n h_i\frac{\partial}{\partial x_i}\right)^K
  (f(x_0+\theta h)-f(x_0)),
\end{aligned}
\]
其中 $0<\theta<1$。

记
\[
  R_K(h)=\frac1{K!}
  \left(\sum_{i=1}^n h_i\frac{\partial}{\partial x_i}\right)^K
  (f(x_0+\theta h)-f(x_0)).
\]
由 $f(x)$ 在 $U(x_0,\delta_0)$ 内具有 $K(K\ge1)$ 阶连续偏导数, 知
\[
  \lim_{|h|\to0}\frac{R_K(h)}{|h|^K}
  =\lim_{|h|\to0}\frac1{K!}
  \left(\sum_{i=1}^n\frac{h_i}{|h|}\frac{\partial}{\partial x_i}\right)^K
  (f(x_0+\theta h)-f(x_0))=0.
\]
上述极限为零是因为当 $|h|\to0$ 时, 和式中的每一项的系数都是有界的, 而 $f(x)$ 在 $x_0+\theta h$ 处的每个 $K$ 阶偏导数由于其连续性都趋于 $f(x)$ 在 $x_0$ 相应的 $K$ 阶偏导数. 证毕。

\noindent\textbf{注}\quad 我们将推论 1 中的余项 $o(|h|^K)$ 称为皮亚诺余项, 并称式 (14.3.3) 为 $f(x)$ 带皮亚诺余项的泰勒公式. 由推论 1 我们可知, 若 $f(x)$ 在 $x_0$ 处能展成
\[
  f(x_0+h)=P_K(h_1,\cdots,h_n)+o(|h|^K)
  \quad(|h|\to0),
\]
其中 $P_K$ 是 $n$ 元 $K$ 次多项式, 上式必定是 $f(x)$ 在 $U(x_0,\delta_0)$ 内的泰勒公式。

现设 $f(x)$ 在区域 $D$ 内具有 $K+1$ 阶连续偏导数, 从定理 14.3.1 的证明中可以看出, 若要求 $f(x)(x\in D)$ 在 $x_0\in D$ 处的泰勒公式, 我们不必要求 $x$ 在 $x_0$ 的某个邻域内, 而只要求连接 $x_0$ 和 $x$ 的线段 $x_0+t(x-x_0)(t\in[0,1])$ 落在 $D$ 内即可。

在定理 14.3.1 中取 $K=0$, 我们有以下多元函数的拉格朗日微分中值定理。

\noindent\textbf{推论 2}\quad 设 $f(x)$ 在区域 $D\subset\mathbb R^n$ 内具有连续偏导数, 且对于 $\forall t\in[0,1],x_0+t(x-x_0)\in D$, 则有
\[
\begin{aligned}
  f(x)-f(x_0)
  &=\sum_{i=1}^n
  \frac{\partial f(x_0+\theta(x-x_0))}{\partial x_i}(x_i-x_i^0)\\
  &=f'(x_0+\theta(x-x_0))\cdot(x-x_0)^\mathrm T,
\end{aligned}
\]
其中 $0<\theta<1$。

由上述推论 2, 我们容易推出下述结论。

\noindent\textbf{推论 3}\quad 设函数 $f(x)$ 在区域 $D\subset\mathbb R^n$ 内的各个偏导数均为 0, 则 $f(x)$ 在 $D$ 内为常数函数。

请读者给出上述推论的证明。

以后我们还常常要用到泰勒公式中 $K=1$ 的情形. 在定理 14.3.1 中, 令 $K=1$, 我们有
\[
  f(x_0+h)=f(x_0)+f'(x_0)\cdot h^\mathrm T
  +\frac12 h\cdot H_f(x_0)h^\mathrm T+o(|h|^2)
  \quad(|h|\to0),
\]
其中
\[
  H_f(x_0)=
  \begin{pmatrix}
    \dfrac{\partial^2 f(x_0)}{\partial x_1^2} &
    \dfrac{\partial^2 f(x_0)}{\partial x_1\partial x_2} &
    \cdots &
    \dfrac{\partial^2 f(x_0)}{\partial x_1\partial x_n}\\
    \dfrac{\partial^2 f(x_0)}{\partial x_2\partial x_1} &
    \dfrac{\partial^2 f(x_0)}{\partial x_2^2} &
    \cdots &
    \dfrac{\partial^2 f(x_0)}{\partial x_2\partial x_n}\\
    \vdots & \vdots & & \vdots\\
    \dfrac{\partial^2 f(x_0)}{\partial x_n\partial x_1} &
    \dfrac{\partial^2 f(x_0)}{\partial x_n\partial x_2} &
    \cdots &
    \dfrac{\partial^2 f(x_0)}{\partial x_n^2}
  \end{pmatrix}
\]
称为 $f(x)$ 在 $x_0$ 处的海色 (Hessi) 矩阵。

\noindent\textbf{例 14.3.1}\quad 求函数 $f(x,y)=\dfrac{x-y}{x+y}$ 在 $(1,0)$ 附近带皮亚诺余项的泰勒公式 (直到二次项)。

\noindent\textbf{解}\quad 对于这类函数, 一般可直接求偏导, 然后写出泰勒公式, 但我们下面利用一元函数已知的泰勒公式来求之. 由于
\[
\begin{aligned}
  f(x,y)
  &=\frac{(x-1)+1-y}{1+[(x-1)+y]}\\
  &=[(x-1)+1-y]\{1-(x-1)-y+[(x-1)+y]^2
  +o((x-1)^2+y^2)\}
\end{aligned}
\]
$((x-1)^2+y^2\to0)$, 将上式整理得
\[
  f(x,y)=1-2y+2(x-1)y+2y^2+o((x-1)^2+y^2),
  \quad ((x-1)^2+y^2\to0).
\]

\noindent\textbf{例 14.3.2}\quad 求函数 $f(x,y)=xe^{x+y}$ 在原点 $(0,0)$ 处的所有四阶偏导数。

\noindent\textbf{解}\quad 由 $e^{x+y}$ 在 $(0,0)$ 处的泰勒公式有
\[
  f(x,y)=xe^{x+y}
  =x\left[1+\sum_{k=1}^K\frac{(x+y)^k}{k!}
  +o\left((\sqrt{x^2+y^2})^K\right)\right]
  \quad(\sqrt{x^2+y^2}\to0).
\]
令 $K=3$, 得
\[
\begin{aligned}
  f(x,y)
  &=x+x^2+xy+\frac12(x^3+2x^2y+xy^2)\\
  &\quad+\frac16(x^4+3x^3y+3x^2y^2+xy^3)
  +o\left((\sqrt{x^2+y^2})^4\right)
  \quad(\sqrt{x^2+y^2}\to0),
\end{aligned}
\]
于是有
\[
  \frac{C_4^4}{4!}\cdot\frac{\partial^4 f(0,0)}{\partial x^4}=\frac16,\quad
  \frac{C_4^3}{4!}\cdot\frac{\partial^4 f(0,0)}{\partial x^3\partial y}=\frac12,\quad
  \frac{C_4^2}{4!}\cdot\frac{\partial^4 f(0,0)}{\partial x^2\partial y^2}=\frac12,
\]
\[
  \frac{C_4^1}{4!}\cdot\frac{\partial^4 f(0,0)}{\partial x\partial y^3}=\frac16,\quad
  \frac{C_4^0}{4!}\cdot\frac{\partial^4 f(0,0)}{\partial y^4}=0,
\]
从而得
\[
  \frac{\partial^4 f(0,0)}{\partial x^4}=4,\quad
  \frac{\partial^4 f(0,0)}{\partial x^3\partial y}=3,\quad
  \frac{\partial^4 f(0,0)}{\partial x^2\partial y^2}=2,
\]
\[
  \frac{\partial^4 f(0,0)}{\partial x\partial y^3}=1,\quad
  \frac{\partial^4 f(0,0)}{\partial y^4}=0.
\]

\section{隐函数存在定理}

\subsection{单个方程的情形}

在讨论一般的隐函数存在问题之前, 我们先来讨论一下由一个二元方程确定隐函数的简单情形. 设 $F(x,y)=0$ 是区域 $D\subset\mathbb R^2$ 上的一个二元方程, 并且满足 $F(x_0,y_0)=0$, 其中 $(x_0,y_0)\in D$. 我们要研究的问题是: 何时方程 $F(x,y)=0$ 在点 $(x_0,y_0)$ 附近能唯一确定一个函数 $y=f(x)$, 使得 $f(x_0)=y_0$, 且 $F(x,f(x))=0$ 对 $x\in(x_0-\delta,x_0+\delta)(\delta>0)$ 成立? 我们称这样定义的函数 $y=f(x)$ 为由方程 $F(x,y)=0$ 所确定的隐函数. 显然这个问题等价于: 何时 $z=F(x,y)$ 的零点集合在 $(x_0,y_0)$ 附近是 $Oxy$ 平面内的一条曲线, 而且这条曲线是 $y=f(x)$ 的图像?

我们先来考查一个例子. $\mathbb R^2$ 中的单位圆周满足的方程为 $x^2+y^2=1$, 我们可将其写成 $x^2+y^2-1=0$. 现考虑函数 $z=F(x,y)=x^2+y^2-1$. 先看一个简单情况, 令 $(x_0,y_0)=(0,1)$, 我们来考查一下为什么 $F(x,y)=0$ 在该点的附近能确定一个隐函数 $y=\sqrt{1-x^2}$. 首先 $(x_0,y_0)=(0,1)$ 是 $F(x,y)=x^2+y^2-1$ 的零点, 并且平面 $x=0$ 与曲面 $z=F(x,y)$ 的交线是
\[
  \begin{cases}
    x=0,\\
    z=y^2-1.
  \end{cases}
\]
对于该曲线容易看出, 对于 $0<\delta_0<1$, 当 $y$ 从 $1-\delta_0$ 连续变化到 $1+\delta_0$ 时, $z$ 严格单调上升并且连续地从负变正. 因此, 对于充分小的 $\delta>0$, 对 $(-\delta,\delta)$ 中的每个 $x'$, 曲线
\[
  \begin{cases}
    x=x',\\
    z=y^2-1+x'^2
  \end{cases}
\]
都有此性质. 显然, 当具有上述性质时, $F(x,y)$ 在点 $(x_0,y_0)$ 附近的零点集合刚好是一个函数 $y=f(x)$ 的图像. 而什么时候能保证一个函数 $z=F(x,y)$ 在固定 $x$ 时关于 $y$ 是一个严格单调函数呢? 从多元函数偏导数的几何意义容易看出, 一个充分条件是 $F'_y(x,y)\ne0$。

对于满足 $F'_y(x,y)=0$ 的点 $(x_0,y_0)$, 读者容易发现此时未必有隐函数 $y=y(x)$ 存在, 如 $F(x,y)=x^2+y^2-1$ 在 $(1,0)$ 处。

下面的定理给出了隐函数存在的充分条件。

\noindent\textbf{定理 14.4.1(隐函数存在定理)}\quad 设二元函数 $F(x,y)$ 在 $U((x_0,y_0),\delta)(\delta>0)$ 内满足以下条件:

(1) $F(x_0,y_0)=0$;

(2) $F(x,y),F'_y(x,y)$ 在 $U((x_0,y_0),\delta)$ 内连续;

(3) $F'_y(x_0,y_0)\ne0$,

则 $\exists\delta_0>0(0<\delta_0<\delta)$, 使得在 $U(x_0,\delta_0)$ 内存在唯一满足下述条件的连续函数 $y=f(x)$:

(a) $y_0=f(x_0)$;

(b) $F(x,f(x))=0,\forall x\in U(x_0,\delta_0)$;

(c) 如果 $F'_x(x,y)$ 在 $U((x_0,y_0),\delta)$ 内连续, 则 $f(x)$ 在 $U(x_0,\delta_0)$ 存在连续导数, 并且有
\[
  f'(x)=-\frac{F'_x(x,f(x))}{F'_y(x,f(x))}.
\]

\noindent\textbf{证明}\quad 不妨设 $F'_y(x_0,y_0)>0$. 由 $F'_y(x,y)$ 的连续性及极限的保号性可知, 存在 $0<\delta_1,\delta_2<\delta$, 使得对于 $\forall(x,y)\in U(x_0,\delta_1)\times U(y_0,\delta_2)$, 有
\[
  F'_y(x,y)>0.
\]
特别地, 若固定 $x=x_0$, 则对于 $\forall y\in U(y_0,\delta_2)$, 有 $F'_y(x_0,y)>0$. 因此 $F(x_0,y)$ 在 $U(y_0,\delta_2)$ 内是 $y$ 的严格上升函数. 注意到 $F(x_0,y_0)=0$, 我们可知 $F(x_0,y_0-\delta_2)<0,F(x_0,y_0+\delta_2)>0$. 再由 $F(x,y)$ 在 $(x_0,y_0-\delta_2)$ 和 $(x_0,y_0+\delta_2)$ 处的连续性, 并利用极限的保号性, 我们可以找到 $\delta_0(0<\delta_0<\delta_1)$, 使得当 $x\in U(x_0,\delta_0)$ 时, 有
\[
  F(x,y_0-\delta_2)<0,\qquad F(x,y_0+\delta_2)>0.
\]
对任意给定的 $\tilde x\in U(x_0,\delta_0)$, 由 $F'_y(\tilde x,y)>0$ 推知, 当 $y\in U(y_0,\delta_2)$ 时, $F(\tilde x,y)$ 连续地从负数 $F(\tilde x,y_0-\delta_2)$ 严格上升到正数 $F(\tilde x,y_0+\delta_2)$ (见图 14.4.1). 因此存在唯一的 $\tilde y\in U(y_0,\delta_2)$, 使得
\[
  F(\tilde x,\tilde y)=0.
\]
这说明了对于 $\forall\tilde x\in U(x_0,\delta_0)$, 有唯一确定的 $\tilde y$ 与之对应. 由函数的定义, 若令 $\tilde y=f(\tilde x)$, 则函数 $\tilde y=f(\tilde x)$ 在 $U(x_0,\delta_0)$ 内有定义, 并且满足 (a),(b). 该函数的唯一性由 $F(x,y)$ 在 $U(x_0,\delta_0)\times U(y_0,\delta_2)$ 内关于 $y$ 严格上升所保证。

下面证 $y=f(x)$ 在 $U(x_0,\delta_0)$ 内连续. 任意取定 $\bar x\in U(x_0,\delta_0)$, 记 $\bar y=f(\bar x)$. 对于充分小的 $\varepsilon>0$, 则有
\[
  F(\bar x,\bar y-\varepsilon)<0,\qquad F(\bar x,\bar y+\varepsilon)>0.
\]
由 $F(x,y)$ 在 $(\bar x,\bar y-\varepsilon)$ 及 $(\bar x,\bar y+\varepsilon)$ 处的连续性, 必存在充分小的 $\delta'(0<\delta'<\delta_0)$, 使得当 $x\in U(\bar x,\delta')\subset U(x_0,\delta_0)$ 时, 有
\[
  F(x,\bar y-\varepsilon)<0,\qquad F(x,\bar y+\varepsilon)>0.
\]
从上面关于隐函数存在性的证明过程知, 当 $x\in U(\bar x,\delta')$ 时, 必有 $f(x)\in(\bar y-\varepsilon,\bar y+\varepsilon)$, 即
\[
  |f(x)-f(\bar x)|<\varepsilon.
\]
这说明 $f(x)$ 在 $U(x_0,\delta_0)$ 内连续。

最后, 设 $F'_x(x,y)$ 在 $U((x_0,y_0),\delta)$ 内存在且连续, 我们证明 $y=f(x)$ 在 $U(x_0,\delta_0)$ 内具有连续导数. 任取 $\bar x\in U(x_0,\delta_0)$, 再取 $\Delta x$ 充分小, 使得 $\bar x+\Delta x\in U(x_0,\delta_0)$. 记
\[
  \bar y=f(\bar x),\qquad
  \Delta y=f(\bar x+\Delta x)-f(\bar x).
\]
由多元函数拉格朗日微分中值定理, $\exists0<\theta<1$, 使得下式成立:
\[
\begin{aligned}
  0&=F(\bar x+\Delta x,\bar y+\Delta y)-F(\bar x,\bar y)\\
  &=F'_x(\bar x+\theta\Delta x,\bar y+\theta\Delta y)\Delta x
  +F'_y(\bar x+\theta\Delta x,\bar y+\theta\Delta y)\Delta y.
\end{aligned}
\]
再由 $F'_y(x,y)\ne0((x,y)\in U((x_0,y_0),\delta))$, 得
\[
  \frac{\Delta y}{\Delta x}
  =-\frac{F'_x(\bar x+\theta\Delta x,\bar y+\theta\Delta y)}
  {F'_y(\bar x+\theta\Delta x,\bar y+\theta\Delta y)}.
\]
令 $\Delta x\to0$, 由 $F'_x(x,y)$ 及 $F'_y(x,y)$ 的连续性得
\[
  f'(\bar x)=-\frac{F'_x(\bar x,f(\bar x))}{F'_y(\bar x,f(\bar x))}.
\]
从上式可以看出 $f'(x)$ 在 $U(x_0,\delta_0)$ 内连续. 证毕。

关于隐函数存在定理, 我们要注意以下几点:

\noindent\textbf{注 1}\quad 隐函数存在定理是一个具有重要理论价值的定理, 这在许多后续课程中读者将会有所体验. 值得指出的是, 在定理的条件下, $F(x,y)=0$ 所确定的隐函数是局部存在的, 并且定理的结论只给出了存在性. 在很多情况下, 我们未必能找到解析式 $y=f(x)$ 来表示这个隐函数. 例如, 对于开普勒 (Kepler) 方程
\[
  y-x-\varepsilon\sin y=0\quad(0<\varepsilon<1),
\]
容易验证它在 $(-\infty,+\infty)$ 上能确定函数 $y=f(x)$, 但人们却无法写出 $y=f(x)$ 的解析式。

\noindent\textbf{注 2}\quad 定理 14.4.1 只是给了隐函数存在的一个充分条件, 若对定理证明过程加以分析, 读者容易举出例子来说明, 当 $F(x,y)$ 不满足定理的条件时, 有时也可以唯一确定一个隐函数 (见本章习题)。

在隐函数存在定理中, 将 $F(x,y)$ 的 $x\in U(x_0,\delta)$ 换成 $\mathbb R^n$ 中的点 $x\in U(x_0,\delta)(\delta>0)$ 时, 相应的隐函数存在和唯一性的结论不仅成立, 而且其证明与定理 14.4.1 的证明相似. 下面不加证明地给出此结论。

\noindent\textbf{定理 14.4.2(隐函数存在定理)}\quad 记 $x=(x_1,x_2,\cdots,x_n),x_0=(x_1^0,x_2^0,\cdots,x_n^0)\in\mathbb R^n$, 假设函数 $F(x,y)=F(x_1,x_2,\cdots,x_n,y)$ 在 $U(x_0,\delta)\times U(y_0,\delta)(\delta>0)$ 内有定义, 并且满足

(1) $F(x_0,y_0)=0$;

(2) $F(x,y)$ 和 $F'_y(x,y)$ 在 $U(x_0,\delta)\times U(y_0,\delta)$ 内连续;

(3) $F'_y(x_0,y_0)\ne0$,

则 $\exists\delta_0(0<\delta_0<\delta)$, 使得在 $U(x_0,\delta_0)$ 内存在唯一满足下述条件的连续函数 $y=f(x)$:

(a) $y_0=f(x_0)$;

(b) 对于 $\forall x\in U(x_0,\delta_0),F(x,f(x))=0$;

(c) 如果 $F(x,y)$ 在 $U(x_0,\delta)\times U(y_0,\delta)$ 内存在各个连续偏导数, 那么 $y=f(x)$ 在 $U(x_0,\delta_0)$ 具有各个连续偏导数, 并且对于 $i=1,2,\cdots,n$ 及 $\forall x\in U(x_0,\delta_0)$, 有
\[
  \frac{\partial f(x)}{\partial x_i}
  =-\frac{F'_{x_i}(x_1,x_2,\cdots,x_n,y)}
  {F'_y(x_1,x_2,\cdots,x_n,y)},\qquad
  \text{其中 }y=f(x).
\]

\noindent\textbf{注}\quad 在定理 14.4.1 和定理 14.4.2 中, 若 $F(x,y)$ (或 $F(x,y)$) 除了满足相应定理的条件外, 再假设它们具有各个高阶偏导数, 从隐函数的一阶导数或偏导数的表达式直接看出, 所确定的隐函数也具有相应的高阶导数或高阶偏导数。

\noindent\textbf{例 14.4.1}\quad 设 $x=x(y,z),y=y(x,z),z=z(x,y)$ 都是由方程 $F(x,y,z)=0$ 确定的隐函数, 并且它们具有连续偏导数. 证明:
\[
  (1)\quad \frac{\partial x}{\partial y}\cdot
  \frac{\partial y}{\partial z}\cdot
  \frac{\partial z}{\partial x}=-1;
\]
\[
  (2)\quad \frac{\partial x}{\partial y}\cdot
  \frac{\partial x}{\partial z}\cdot
  \frac{\partial y}{\partial x}\cdot
  \frac{\partial y}{\partial z}\cdot
  \frac{\partial z}{\partial x}\cdot
  \frac{\partial z}{\partial y}=1.
\]

\noindent\textbf{证明}\quad (1) 因为
\[
  \frac{\partial x}{\partial y}=-\frac{F'_y(x,y,z)}{F'_x(x,y,z)},\quad
  \frac{\partial y}{\partial z}=-\frac{F'_z(x,y,z)}{F'_y(x,y,z)},\quad
  \frac{\partial z}{\partial x}=-\frac{F'_x(x,y,z)}{F'_z(x,y,z)},
\]
所以
\[
  \frac{\partial x}{\partial y}\cdot
  \frac{\partial y}{\partial z}\cdot
  \frac{\partial z}{\partial x}=-1.
\]

(2) 由 (1), 又因为
\[
  \frac{\partial x}{\partial z}=-\frac{F'_z(x,y,z)}{F'_x(x,y,z)},\quad
  \frac{\partial y}{\partial x}=-\frac{F'_x(x,y,z)}{F'_y(x,y,z)},\quad
  \frac{\partial z}{\partial y}=-\frac{F'_y(x,y,z)}{F'_z(x,y,z)},
\]
所以
\[
  \frac{\partial x}{\partial y}\cdot
  \frac{\partial x}{\partial z}\cdot
  \frac{\partial y}{\partial x}\cdot
  \frac{\partial y}{\partial z}\cdot
  \frac{\partial z}{\partial x}\cdot
  \frac{\partial z}{\partial y}=1.
\]

\noindent\textbf{注}\quad 在一元微分学中, 当 $y$ 是 $x$ 的函数时, $\dfrac{\mathrm dy}{\mathrm dx}$ 可以看成是 $\mathrm dy$ 与 $\mathrm dx$ 的商. 上例说明, 在多元微分学中, 当 $z=f(x,y)$ 可偏导时, $\dfrac{\partial z}{\partial x}$ 则是一个整体记号。

\noindent\textbf{例 14.4.2}\quad 证明在 $(0,0,0)$ 的邻域内方程
\begin{equation}
  -2x+y-x^2+y^2+z+\sin z=0
  \tag{14.4.1}
\end{equation}
确定隐函数 $z=f(x,y)$, 并求 $f(x,y)$ 在 $(0,0)$ 处带皮亚诺余项的泰勒公式 (直到二次项)。

\noindent\textbf{解}\quad 令 $F(x,y,z)=-2x+y-x^2+y^2+z+\sin z$, 则 $F(x,y,z)$ 具有各阶连续偏导数且满足 $F(0,0,0)=0$ 和 $F'_z(0,0,0)=1+\cos0=2\ne0$. 因此 $F(x,y,z)=0$ 在 $(0,0,0)$ 的某一邻域内唯一确定隐函数 $z=f(x,y)$, 并且 $f(x,y)$ 在该邻域内具有各阶连续偏导数。

注意到 $f(0,0)=0$, 因此 $z=f(x,y)$ 的泰勒公式具有下述形式:
\begin{equation}
  f(x,y)=a_1x+a_2y+b_{11}x^2+b_{12}xy+b_{22}y^2+o(x^2+y^2)
  \tag{14.4.2}
\end{equation}
\[
  (\sqrt{x^2+y^2}\to0),
\]
其中 $a_i,b_{ij}(i,j=1,2)$ 为常数. 由于
\[
  \sin z=z+o(|z|^2)
  =a_1x+a_2y+b_{11}x^2+b_{12}xy+b_{22}y^2+o(x^2+y^2)
  \tag{14.4.3}
\]
\[
  (\sqrt{x^2+y^2}\to0),
\]
将式 (14.4.2),(14.4.3) 代入式 (14.4.1) 得
\[
  -2x+y-x^2+y^2+2a_1x+2a_2y+2b_{11}x^2+2b_{12}xy+2b_{22}y^2+o(x^2+y^2)=0.
\]
由此推出
\[
  (2a_1-2)x+(2a_2+1)y+(2b_{11}-1)x^2
  +2b_{12}xy+(2b_{22}+1)y^2+o(x^2+y^2)=0,
\]
从而有
\[
  a_1=1,\qquad a_2=-\frac12,\qquad b_{11}=\frac12,\qquad
  b_{12}=0,\qquad b_{22}=-\frac12.
\]
这样我们就求出了如下 $f(x,y)$ 的带皮亚诺余项的泰勒公式:
\[
  f(x,y)=x-\frac12y+\frac12x^2-\frac12y^2+o(x^2+y^2)
  \quad(\sqrt{x^2+y^2}\to0).
\]

\subsection{方程组的情形}

我们下面来考查两个方程所组成的方程组
\[
  \begin{cases}
    F(x,u,v)=0,\\
    G(x,u,v)=0
  \end{cases}
\]
何时能在 $(x_0,u_0,v_0)$ 的邻域内确定 $u,v$ 是 $x$ 的函数. 对一个方程 $F(x,y)=0$ 的情形, $F'_y(x_0,y_0)\ne0$ 起着重要作用. 因此, 当我们考虑向量函数
\[
  \begin{pmatrix}
    z_1\\
    z_2
  \end{pmatrix}
  =
  \begin{pmatrix}
    F(x,u,v)\\
    G(x,u,v)
  \end{pmatrix}
\]
时, 若将 $x$ 看做常量, 它在 $(x_0,u_0,v_0)$ 处关于 $(u,v)$ 的导数将起着隐函数存在定理中 $F'_y(x_0,y_0)$ 的作用. 注意到此时该导数是
\[
  \left.
  \begin{pmatrix}
    F'_u & F'_v\\
    G'_u & G'_v
  \end{pmatrix}
  \right|_{(x_0,u_0,v_0)}.
\]
由于要确定两个隐函数, 因此该雅可比矩阵应该是非退化的, 即
\[
  \left.
  \begin{vmatrix}
    F_u & F_v\\
    G_u & G_v
  \end{vmatrix}
  \right|_{(x_0,u_0,v_0)}\ne0.
\]
对于上述问题, 我们有以下一般性的定理。

\noindent\textbf{定理 14.4.3(隐函数组存在定理)}\quad 设向量函数
\[
  F(x,u)=(F_1(x,u),F_2(x,u),\cdots,F_m(x,u))
\]
在 $U(x_0,\delta)\times U(u_0,\delta)(\delta>0)$ 内有定义, 其中 $x=(x_1,x_2,\cdots,x_n)\in U(x_0,\delta)$, $u=(u_1,u_2,\cdots,u_m)\in U(u_0,\delta)$, 并且满足以下条件:

(1) $F_j(x_0,u_0)=0,\quad j=1,2,\cdots,m$;

(2) 对于 $j=1,2,\cdots,m$, $F_j(x,u)$ 以及它的各个偏导数在 $U(x_0,\delta)\times U(u_0,\delta)$ 内连续;

(3)
\[
  \left.
  \frac{\partial(F_1,F_2,\cdots,F_m)}
  {\partial(u_1,u_2,\cdots,u_m)}
  \right|_{(x_0,u_0)}
  \ne0,
\]

则 $\exists\delta_0(0<\delta_0<\delta)$, 使得在 $U(x_0,\delta_0)$ 内存在唯一的 $m$ 维 $n$ 元向量函数
\[
  f(x)=(f_1(x),f_2(x),\cdots,f_m(x)),
\]
满足

(a) $u_0=(f_1(x_0),f_2(x_0),\cdots,f_m(x_0))$;

(b) 对于 $\forall j(1\le j\le m)$ 及 $\forall x\in U(x_0,\delta_0)$, 有
\[
  F_j(x,f_1(x),\cdots,f_m(x))=0;
\]

(c) $f(x)$ 的每个分量函数 $f_j(x)(j=1,2,\cdots,m)$ 在 $U(x_0,\delta_0)$ 内存在连续偏导数, 记
\[
  A=
  \begin{pmatrix}
    \dfrac{\partial F_1(x,u)}{\partial x_1} &
    \dfrac{\partial F_1(x,u)}{\partial x_2} &
    \cdots &
    \dfrac{\partial F_1(x,u)}{\partial x_n}\\
    \dfrac{\partial F_2(x,u)}{\partial x_1} &
    \dfrac{\partial F_2(x,u)}{\partial x_2} &
    \cdots &
    \dfrac{\partial F_2(x,u)}{\partial x_n}\\
    \vdots & \vdots & & \vdots\\
    \dfrac{\partial F_m(x,u)}{\partial x_1} &
    \dfrac{\partial F_m(x,u)}{\partial x_2} &
    \cdots &
    \dfrac{\partial F_m(x,u)}{\partial x_n}
  \end{pmatrix},
\]
\[
  B=
  \begin{pmatrix}
    \dfrac{\partial F_1(x,u)}{\partial u_1} &
    \dfrac{\partial F_1(x,u)}{\partial u_2} &
    \cdots &
    \dfrac{\partial F_1(x,u)}{\partial u_m}\\
    \dfrac{\partial F_2(x,u)}{\partial u_1} &
    \dfrac{\partial F_2(x,u)}{\partial u_2} &
    \cdots &
    \dfrac{\partial F_2(x,u)}{\partial u_m}\\
    \vdots & \vdots & & \vdots\\
    \dfrac{\partial F_m(x,u)}{\partial u_1} &
    \dfrac{\partial F_m(x,u)}{\partial u_2} &
    \cdots &
    \dfrac{\partial F_m(x,u)}{\partial u_m}
  \end{pmatrix},
\]
那么
\[
  f'(x)=-B^{-1}\cdot A,
\]
其中矩阵 $B^{-1}$ 为 $B$ 的逆矩阵。

\noindent\textbf{证明}\quad 为了证明本定理, 可以对 $m$ 作数学归纳法. 对 $m=2$ 的证明是本质的, 而对 $m\ge3$ 的证明则只具有形式的复杂性. 在此我们只给出 $m=2$ 的证明。

由于
\[
  \left.
  \frac{\partial(F_1,F_2)}{\partial(u_1,u_2)}
  \right|_{(x_0,u_0)}
  =
  \left.
  \begin{vmatrix}
    \dfrac{\partial F_1}{\partial u_1} &
    \dfrac{\partial F_1}{\partial u_2}\\
    \dfrac{\partial F_2}{\partial u_1} &
    \dfrac{\partial F_2}{\partial u_2}
  \end{vmatrix}
  \right|_{(x_0,u_0)}
  \ne0,
\]
不失一般性, 我们假定 $\dfrac{\partial F_2(x_0,u_0)}{\partial u_2}\ne0$. 记 $u_0=(u_1^0,u_2^0)$. 由定理 14.4.2 知, $\exists\delta_1(0<\delta_1<\delta)$, 使得方程 $F_2(x,u_1,u_2)=0$ 在 $U(x_0,\delta_1)\times U(u_1^0,\delta_1)$ 内唯一确定一个函数
\[
  u_2=g(x,u_1),
\]
满足
\[
  F_2(x,u_1,g(x,u_1))=0,\qquad u_2^0=g(x_0,u_1^0),
\]
且 $g(x,u_1)$ 在该邻域内具有各个连续偏导数. 将 $u_2=g(x,u_1)$ 代入方程 $F_1(x,u_1,u_2)=0$, 得到一个关于 $(x,u_1)$ 的方程
\[
  H(x,u_1)\triangleq F_1(x,u_1,g(x,u_1))=0.
\]
由于
\[
\begin{aligned}
  \frac{\partial H}{\partial u_1}
  &=\frac{\partial F_1}{\partial u_1}
  +\frac{\partial F_1}{\partial u_2}\cdot
  \frac{\partial g(x,u_1)}{\partial u_1}\\
  &=\frac{\partial F_1}{\partial u_1}
  +\frac{\partial F_1}{\partial u_2}
  \left(-\frac{\partial F_2}{\partial u_1}\bigg/
  \frac{\partial F_2}{\partial u_2}\right)\\
  &=\frac1{\dfrac{\partial F_2}{\partial u_2}}\cdot
  \frac{\partial(F_1,F_2)}{\partial(u_1,u_2)},
\end{aligned}
\]
我们有
\[
  \frac{\partial H(x_0,u_1^0)}{\partial u_1}
  =\frac1{\dfrac{\partial F_2(x_0,u_1^0,u_2^0)}{\partial u_2}}\cdot
  \left.
  \frac{\partial(F_1,F_2)}{\partial(u_1,u_2)}
  \right|_{(x_0,u_1^0,u_2^0)}
  \ne0.
\]
因此 $\exists\delta_0(0<\delta_0<\delta_1)$, 使得 $H(x,u_1)=0$ 在 $U(x_0,\delta_0)$ 内唯一确定一个函数 $u_1=f_1(x)$, 满足 $u_1^0=f_1(x_0)$, 且
\[
  F_1(x,f_1(x),g(x,f_1(x)))=0.
\]
现记
\[
  u_1=f_1(x),\qquad u_2=f_2(x)=g(x,f_1(x)),
\]
则在 $U(x_0,\delta_0)$ 内, 有
\[
  \begin{cases}
    F_1(x,f_1(x),f_2(x))=0,\\
    F_2(x,f_1(x),f_2(x))=0
  \end{cases}
\]
和 $u_0=(f_1(x_0),f_2(x_0))$. 再由 $F(x,u)$ 具有各个连续偏导数, 我们可以推出 $f_1(x),f_2(x)$ 在 $U(x_0,\delta_0)$ 内具有各个连续偏导数 (从而 $f_1(x)$ 与 $f_2(x)$ 在该邻域内连续)。

对于 $\forall i(1\le i\le n)$, 利用复合函数的求导法则有
\[
  \begin{cases}
    \dfrac{\partial F_1}{\partial x_i}
    +\dfrac{\partial F_1}{\partial u_1}\cdot\dfrac{\partial f_1}{\partial x_i}
    +\dfrac{\partial F_1}{\partial u_2}\cdot\dfrac{\partial f_2}{\partial x_i}=0,\\[1.2ex]
    \dfrac{\partial F_2}{\partial x_i}
    +\dfrac{\partial F_2}{\partial u_1}\cdot\dfrac{\partial f_1}{\partial x_i}
    +\dfrac{\partial F_2}{\partial u_2}\cdot\dfrac{\partial f_2}{\partial x_i}=0.
  \end{cases}
\]
将上式写成向量形式有
\[
  \begin{pmatrix}
    \dfrac{\partial F_1}{\partial u_1} &
    \dfrac{\partial F_1}{\partial u_2}\\
    \dfrac{\partial F_2}{\partial u_1} &
    \dfrac{\partial F_2}{\partial u_2}
  \end{pmatrix}
  \begin{pmatrix}
    \dfrac{\partial f_1}{\partial x_i}\\
    \dfrac{\partial f_2}{\partial x_i}
  \end{pmatrix}
  =-
  \begin{pmatrix}
    \dfrac{\partial F_1}{\partial x_i}\\
    \dfrac{\partial F_2}{\partial x_i}
  \end{pmatrix},
\]
因此有

\[
  \begin{pmatrix}
    \dfrac{\partial F_1}{\partial u_1} &
    \dfrac{\partial F_1}{\partial u_2}\\
    \dfrac{\partial F_2}{\partial u_1} &
    \dfrac{\partial F_2}{\partial u_2}
  \end{pmatrix}
  \begin{pmatrix}
    \dfrac{\partial f_1}{\partial x_1} &
    \dfrac{\partial f_1}{\partial x_2} &
    \cdots &
    \dfrac{\partial f_1}{\partial x_n}\\
    \dfrac{\partial f_2}{\partial x_1} &
    \dfrac{\partial f_2}{\partial x_2} &
    \cdots &
    \dfrac{\partial f_2}{\partial x_n}
  \end{pmatrix}
  =-
  \begin{pmatrix}
    \dfrac{\partial F_1}{\partial x_1} &
    \dfrac{\partial F_1}{\partial x_2} &
    \cdots &
    \dfrac{\partial F_1}{\partial x_n}\\
    \dfrac{\partial F_2}{\partial x_1} &
    \dfrac{\partial F_2}{\partial x_2} &
    \cdots &
    \dfrac{\partial F_2}{\partial x_n}
  \end{pmatrix}.
\]
由此得
\[
  \begin{pmatrix}
    \dfrac{\partial f_1}{\partial x_1} &
    \dfrac{\partial f_1}{\partial x_2} &
    \cdots &
    \dfrac{\partial f_1}{\partial x_n}\\
    \dfrac{\partial f_2}{\partial x_1} &
    \dfrac{\partial f_2}{\partial x_2} &
    \cdots &
    \dfrac{\partial f_2}{\partial x_n}
  \end{pmatrix}
  =-
  \begin{pmatrix}
    \dfrac{\partial F_1}{\partial u_1} &
    \dfrac{\partial F_1}{\partial u_2}\\
    \dfrac{\partial F_2}{\partial u_1} &
    \dfrac{\partial F_2}{\partial u_2}
  \end{pmatrix}^{-1}
  \begin{pmatrix}
    \dfrac{\partial F_1}{\partial x_1} &
    \dfrac{\partial F_1}{\partial x_2} &
    \cdots &
    \dfrac{\partial F_1}{\partial x_n}\\
    \dfrac{\partial F_2}{\partial x_1} &
    \dfrac{\partial F_2}{\partial x_2} &
    \cdots &
    \dfrac{\partial F_2}{\partial x_n}
  \end{pmatrix}.
\]
现证所确定的隐函数组是唯一的. 事实上, 倘若由
\[
  F(x,u_1,u_2)=0
\]
在 $U(x_0,\delta_0)$ 内确定另一组函数
\[
  \begin{cases}
    u_1=\tilde f_1(x),\\
    u_2=\tilde f_2(x),
  \end{cases}
\]
并且它们仍然满足定理结论 (a), (b) 和 (c), 则容易看出 $\tilde f_1$ 与 $f_1$, $\tilde f_2$ 与 $f_2$ 在该邻域内具有相同的各个偏导数, 从而 $f_1(x)-\tilde f_1(x)$, $f_2(x)-\tilde f_2(x)$ 在 $U(x_0,\delta_0)$ 内均为常数函数. 但它们在 $x_0$ 处有 $f_1(x_0)=\tilde f_1(x_0)$, $f_2(x_0)=\tilde f_2(x_0)$, 因此在 $U(x_0,\delta_0)$ 内有 $f_1(x)=\tilde f_1(x)$, $f_2(x)=\tilde f_2(x)$. 证毕。

\noindent\textbf{例 14.4.3}\quad 证明方程组
\begin{equation}
  \begin{cases}
    x^2+y^2+u^2-v^2=2,\\
    u+v+x+y=0
  \end{cases}
  \tag{14.4.4}
\end{equation}
在 $(x_0,y_0,u_0,v_0)=(1,1,-1,-1)$ 的某个邻域 $U((1,1,-1,-1),\delta_0)(\delta_0>0)$ 内确定隐函数组
\[
  \begin{cases}
    u=u(x,y),\\
    v=v(x,y),
  \end{cases}
\]
并在该邻域内求 $\dfrac{\partial u}{\partial x}$, $\dfrac{\partial v}{\partial y}$。

\noindent\textbf{解}\quad 显然点 $x_0=(1,1,-1,-1)$ 满足方程组 (14.4.4). 记
\[
  F(x,y,u,v)=x^2+y^2+u^2-v^2-2,
  \qquad
  G(x,y,u,v)=u+v+x+y,
\]
则 $F,G$ 都具有连续偏导数, 且
\[
  \left.
  \frac{\partial(F,G)}{\partial(u,v)}
  \right|_{x_0}
  =
  \left.
  \begin{vmatrix}
    2u & -2v\\
    1 & 1
  \end{vmatrix}
  \right|_{x_0}
  =-2-2=-4\ne0.
\]
因此, $\exists\delta_0>0$, 使得方程组 (14.4.4) 在 $U((1,1,-1,-1),\delta_0)$ 内唯一确定隐函数组
\[
  \begin{cases}
    u=u(x,y),\\
    v=v(x,y),
  \end{cases}
\]
并且 $u(x,y)$ 与 $v(x,y)$ 在 $U((1,1),\delta_0)$ 内存在连续偏导数. 下面我们来求偏导数. 对方程组 (14.4.4) 微分得
\begin{equation}
  \begin{cases}
    2x\,dx+2y\,dy+2u\,du-2v\,dv=0,\\
    du+dv+dx+dy=0.
  \end{cases}
  \tag{14.4.5}
\end{equation}
令 $dy=0$, 得
\[
  \begin{cases}
    u\,du-v\,dv=-x\,dx,\\
    du+dv=-dx.
  \end{cases}
\]
从中解出 $du=-\dfrac{x+v}{u+v}\,dx$, 从而
\[
  \frac{\partial u}{\partial x}=-\frac{x+v}{u+v}.
\]
在方程组 (14.4.5) 中令 $dx=0$, 即得
\[
  \begin{cases}
    u\,du-v\,dv=-y\,dy,\\
    du+dv=-dy,
  \end{cases}
\]
从而可求出
\[
  \frac{\partial v}{\partial y}=\frac{y-u}{u+v}.
\]

\subsection{逆映射存在定理}

对于一个可微映射, 如何确定它的逆映射的存在性呢? 隐函数组存在定理可以给出一个连续可微映射局部存在逆映射的充分条件。

\noindent\textbf{定理 14.4.4(逆映射存在定理)}\quad 设
\[
  y=(y_1,y_2,\cdots,y_n)=(f_1(x),f_2(x),\cdots,f_n(x))=f(x)
\]
是区域 $D\subset\mathbb R^n$ 到区域 $\Omega\subset\mathbb R^n$ 的一个 $C^1$ 映射, 并且在 $x_0\in D$ 处有
\[
  \left.
  \frac{\partial(f_1,f_2,\cdots,f_n)}
       {\partial(x_1,x_2,\cdots,x_n)}
  \right|_{x_0}
  \ne0.
\]
记 $y_0=f(x_0)$, 则存在 $x_0$ 的邻域 $U(x_0,\delta_0)\subset D$, 使得映射 $y=f(x)$ 是 $U(x_0,\delta_0)$ 到 $f(U(x_0,\delta_0))$ 的 $C^1$ 同胚映射 (变换), 其中 $f(U(x_0,\delta_0))$ 是包含 $y_0$ 的一个区域。

\noindent\textbf{证明}\quad 对于 $j=1,2,\cdots,n$, 记 $F_j(x,y)=y_j-f_j(x)$, 并考虑下面 $n$ 个方程组成的方程组
\[
  \begin{cases}
    F_1(x,y)=0,\\
    F_2(x,y)=0,\\
    \cdots\\
    F_n(x,y)=0.
  \end{cases}
\]
由定理条件知它们满足

\noindent (1) 对任意的 $j(1\le j\le n)$, $F_j(x_0,y_0)=0$;

\noindent (2) $\exists\delta'>0$, 对任意的 $j(1\le j\le n)$, 使得 $F_j(x,y)$ 在 $U((x_0,y_0),\delta')$ 内具有连续偏导数 (因此 $F_j(x,y)$ 在该邻域连续), 并且
\[
  \left.
  \frac{\partial(F_1,F_2,\cdots,F_n)}
       {\partial(x_1,x_2,\cdots,x_n)}
  \right|_{(x_0,y_0)}
  =
  (-1)^n
  \left.
  \frac{\partial(f_1,f_2,\cdots,f_n)}
       {\partial(x_1,x_2,\cdots,x_n)}
  \right|_{x_0}
  \ne0,
\]
因此由定理 14.4.3 知, 在 $y_0$ 的某个邻域 $U(y_0,\delta_1)(0<\delta_1<\delta')$ 内存在一个 $n$ 维向量函数
\[
  x=g(y)=(g_1(y),\cdots,g_n(y)),
\]
满足 $x_0=(g_1(y_0),\cdots,g_n(y_0))$ 和
\[
  \begin{cases}
    y_1-f_1(g_1(y),\cdots,g_n(y))=0,\\
    \cdots\\
    y_n-f_n(g_1(y),\cdots,g_n(y))=0,
  \end{cases}
\]
并且 $(g_1(y),\cdots,g_n(y))$ 在 $U(y_0,\delta_1)$ 内具有各个连续偏导数. 上面的 $n$ 个等式同时表明, 在 $U(y_0,\delta_1)$ 内 $f(g(y))\equiv y$, 即 $x=g(y)$ 是 $y=f(x)$ 的逆映射. 我们从定理 14.4.3 推知 $f'(x)g'(y)=E$ ($n$ 阶单位矩阵), 因此由 $|f'(x_0)|\ne0$ 有 $g'(y_0)=[f'(x_0)]^{-1}$, 从而 $g(y)$ 在 $y_0$ 处的雅可比行列式也不为零. 如果将 $f(x)$ 与 $g(y)$ 的地位对换, 则在 $x_0$ 的某个邻域 $U(x_0,\delta_1')(0<\delta_1'<\delta')$ 内存在 $g(y)$ 的逆映射, 显然该逆映射是 $f(x)$。

现证存在 $x_0$ 的邻域 $U(x_0,\delta_0)$, 使得 $f(x)$ 是 $U(x_0,\delta_0)$ 到区域 $f(U(x_0,\delta_0))$ 的 $C^1$ 同胚映射. 首先取 $\delta_2(0<\delta_2<\delta_1)$ 充分小, 使得对于 $\forall x\in U(x_0,\delta_2)$, $f(x)$ 在 $x$ 处的雅可比行列式不等于零. 下面我们证明, 对于 $\forall x'\in U(x_0,\delta_2)$ 及任何包含 $x'$ 的开集 $U_{x'}$, $f(x')$ 是 $f(U_{x'})$ 的内点. 由于我们证明该事实只是利用 $f(x)$ 是 $C^1$ 映射, 并且它在 $x'$ 处的雅可比行列式不为零, 容易看出, 我们只要对 $x_0$ 证明上述结论即可. 任取包含 $x_0$ 的开集 $U_{x_0}$, 则由 $g(y)$ 在 $y_0$ 处的连续性, 存在 $y_0$ 的邻域 $U(y_0,\delta_3)$ (当然是开集), 使得 $g(U(y_0,\delta_3))\subset U_{x_0}$. 因此有
\[
  U(y_0,\delta_3)=f(g(U(y_0,\delta_3)))\subset f(U_{x_0}).
\]
再注意到连续映射将连通集映成连通集, 因此 $f(U(x_0,\delta_2))$ 是一个包含 $y_0$ 的区域. 同理, $\exists\delta_2'(0<\delta_2'<\delta_1')$, 使得 $g(U(y_0,\delta_2'))$ 是包含 $x_0$ 的一个区域. 因此取 $\delta_0(0<\delta_0<\delta_2')$ 充分小, 使得 $f(U(x_0,\delta_0))\subset U(y_0,\delta_2')$, 则 $f(U(x_0,\delta_0))$ 是 $\mathbb R^n$ 中包含 $y_0$ 的区域. 由于在 $U(x_0,\delta_0)$ 满足 $g(f(x))=x$, 并且 $f$ 在 $U(x_0,\delta_0)$ 内是 $C^1$ 映射, $g(y)=f^{-1}(y)$ 在 $f(U(x_0,\delta_0))$ 内是 $C^1$ 映射, 因此 $f(x)$ 是 $U(x_0,\delta_0)$ 到 $f(U(x_0,\delta_0))$ 的 $C^1$ 同胚映射 (变换). 证毕。

从逆映射存在定理的证明可知, 如果一个 $n$ 维 $n$ 元向量函数
\[
  y=f(x)=(f_1(x),f_2(x),\cdots,f_n(x))
\]
在区域 $D\subset\mathbb R^n$ 内具有连续偏导数, 并且对于 $\forall x\in D$, $f'(x)$ 非退化, 即它的雅可比行列式
\[
  \left.
  \frac{\partial(f_1,f_2,\cdots,f_n)}
       {\partial(x_1,x_2,\cdots,x_n)}
  \right|_x
  \ne0,
\]
则局部总存在逆映射, 因此 $f$ 必将开集映成开集. 另外, 由于此时 $y=f(x)$ 是连续函数, 因此推出 $f(D)$ 为一个区域。

\noindent\textbf{注 1}\quad 当一个 $n$ 维 $n$ 元向量函数 $y=f(x)$ 在区域 $D\subset\mathbb R^n$ 内具有连续偏导数, 并且对于 $\forall x\in D$ 有 $|f'(x)|\ne0$ 时, 我们仅仅能推出 $y=f(x)$ 是局部同胚映射, 但不能推出它是 $D$ 到 $f(D)$ 的整体同胚映射. 例如:
\[
  \begin{pmatrix}
    u\\
    v
  \end{pmatrix}
  =
  \begin{pmatrix}
    e^x\cos y\\
    e^x\sin y
  \end{pmatrix},
  \qquad (x,y)\in\mathbb R^2
\]
是 $\mathbb R^2$ 到 $\mathbb R^2\setminus\{(0,0)\}$ 的映射, 对于 $\forall(x,y)\in\mathbb R^2$, 有
\[
  \frac{\partial(u,v)}{\partial(x,y)}
  =
  \begin{vmatrix}
    e^x\cos y & -e^x\sin y\\
    e^x\sin y & e^x\cos y
  \end{vmatrix}
  =e^x\ne0,
\]
因此该映射局部总是同胚映射. 由 $\cos y$ 及 $\sin y$ 的周期性, 上述映射显然不是 $\mathbb R^2\to\mathbb R^2\setminus\{(0,0)\}$ 的同胚映射。

\noindent\textbf{注 2}\quad 对于逆映射存在定理, 我们容易举例说明, 如果仅仅要求存在逆映射, 而不要求逆映射具有连续偏导数, 雅可比行列式不为零的条件不是必要条件. 例如,
\[
  \begin{pmatrix}
    u\\
    v
  \end{pmatrix}
  =
  \begin{pmatrix}
    x^3\\
    y
  \end{pmatrix}
\]
是 $\mathbb R^2\to\mathbb R^2$ 的同胚映射. 事实上, 它的逆映射为
\[
  \begin{pmatrix}
    x\\
    y
  \end{pmatrix}
  =
  \begin{pmatrix}
    u^{1/3}\\
    v
  \end{pmatrix},
\]
但在 $(0,0)$ 处有
\[
  \left.
  \frac{\partial(u,v)}{\partial(x,y)}
  \right|_{(0,0)}
  =0.
\]
此例同时说明了当 $\dfrac{\partial x}{\partial u}$ (或 $\dfrac{\partial y}{\partial u}$) 不是处处存在时, 仍然具有连续的逆映射。

\section{多元函数的极值}

\subsection{通常极值问题}

在一元函数中, 我们利用导数讨论了函数的极值. 在本节中, 我们来讨论多元函数的极值问题. 对于多元函数来说, 极值的定义与一元函数是相似的。

\noindent\textbf{定义 14.5.1}\quad 设函数 $u=f(x)$ 在区域 $D$ 内有定义, $x_0\in D$, 若存在 $x_0$ 的邻域 $U(x_0,\delta_0)\subset D(\delta_0>0)$, 使得当 $x\in U(x_0,\delta_0)$ 时, 有 $f(x)\le f(x_0)$ (或 $f(x)\ge f(x_0)$), 则称 $f(x)$ 在 $x_0$ 处取极大值 (或极小值), $x_0$ 称为 $f(x)$ 的极大值点 (或极小值点). 如果上述不等式为严格不等式, 则称 $x_0$ 为 $f(x)$ 的严格极大值点 (或严格极小值点)。

对于一个二元可微函数 $z=f(x,y)((x,y)\in D)$, 一般来说它的图像是一块曲面. 从几何上来看, 若 $(x_0,y_0)\in D$ 是 $z=f(x,y)$ 的极值点, 则 $(x_0,y_0,f(x_0,y_0))$ 是曲面上局部的最高点或最低点. 显然, 对于曲面上的任何一条过 $(x_0,y_0,f(x_0,y_0))$ 的曲线 $\Gamma$, 该点也是 $\Gamma$ 在它附近的最高点或最低点. 当 $\Gamma$ 在该点具有切线时, 它的切线应该平行于 $Oxy$ 平面. 所以, 当 $z=f(x,y)$ 在 $(x_0,y_0)$ 处可微时, 有
\[
  \frac{\partial f(x_0,y_0)}{\partial x}
  =
  \frac{\partial f(x_0,y_0)}{\partial y}
  =0,
\]
即 $f'(x_0,y_0)=\operatorname{grad} f(x_0)=0$。

\noindent\textbf{定理 14.5.1}\quad 设函数 $f(x)$ 在区域 $D\subset\mathbb R^n$ 内有定义, $x_0=(x_1^0,x_2^0,\cdots,x_n^0)\in D$, 再设 $f(x)$ 在 $x_0$ 处取极值并且 $f(x)$ 在该点关于 $x_i(1\le i\le n)$ 可偏导, 则有
\[
  \frac{\partial f(x_0)}{\partial x_i}=0.
\]
特别地, 若 $f(x)$ 在该点可微, 则 $f'(x_0)=0$。

\noindent\textbf{证明}\quad 在定理的条件下, 一元函数
\[
  F_i(x_i)=f(x_1^0,\cdots,x_{i-1}^0,x_i,x_{i+1}^0,\cdots,x_n^0)
  \qquad (1\le i\le n)
\]
在 $x_i^0$ 处取极值, 因此
\[
  F_i'(x_i^0)=\frac{\partial f(x_0)}{\partial x_i}=0
  \qquad (1\le i\le n).
\]
证毕。

如同一元函数, 对于一个 $n$ 元可微函数 $f(x)$, 若 $f'(x_0)=0$, 则称 $x_0$ 为 $f(x)$ 的一个驻点或临界点. 因此, 对可微函数 $f(x)$ 而言, $f(x)$ 的极值点一定是驻点, 但反命题一般不真。

例如, $f(x,y)=y^2-x^2$ 在 $(0,0)$ 处显然不取极值. (如图 14.5.1 所示, 该函数的图像在原点附近有点像一个马鞍, 因此人们称之为马鞍面). 当一个可微函数 $f(x)$ 的驻点不是极值点时, 该驻点也称为 $f(x)$ 的一个鞍点。

在研究一元函数的极值问题中, 我们曾经利用它的高阶导数来判别驻点是否为极值点. 对于多元函数, 我们也可以利用高阶偏导数来讨论其极值问题. 由于多元函数高阶偏导数的复杂性, 我们只考虑涉及二阶偏导数的情形。

\noindent\textbf{定理 14.5.2}\quad 设函数 $f(x)=f(x_1,x_2,\cdots,x_n)$ 在区域 $D\subset\mathbb R^n$ 内具有二阶连续偏导数, 且 $f'(x_0)=0(x_0\in D)$, 再设 $f(x)$ 在 $x_0$ 处的海色矩阵 $H_f(x_0)$ 为满秩矩阵, 则

\noindent (1) 当 $H_f(x_0)$ 正定时, $f(x)$ 在 $x_0$ 处取极小值;

\noindent (2) 当 $H_f(x_0)$ 负定时, $f(x)$ 在 $x_0$ 处取极大值;

\noindent (3) 当 $H_f(x_0)$ 不定时, $f(x)$ 在 $x_0$ 处不取极值。

\noindent\textbf{证明}\quad 在 $x_0$ 的充分小邻域 $U(x_0,\delta_0)(\delta_0>0)$ 内, $f(x)$ 有泰勒公式
\[
\begin{aligned}
  f(x)
  &=f(x_0)+f'(x_0)(x-x_0)^{\mathrm T}
    +\frac12(x-x_0)H_f(x_0)(x-x_0)^{\mathrm T}
    +o(|x-x_0|^2)\\
  &=f(x_0)+\frac12
    \frac{x-x_0}{|x-x_0|}
    H_f(x_0)
    \frac{(x-x_0)^{\mathrm T}}{|x-x_0|}
    |x-x_0|^2
    +o(|x-x_0|^2)
    \qquad (|x-x_0|\to0).
\end{aligned}
\]
对于 $\forall x\in U_0(x_0,\delta_0)$, 记 $x'=\dfrac{x-x_0}{|x-x_0|}$, 则 $|x'|=1$. 由于 $x'H_f(x_0)x'^{\mathrm T}$ 是 $S(0;1)=\{x':|x'|=1\}$ 上的一个连续函数, 注意到 $S(0;1)$ 是紧集, $x'H_f(x_0)x'^{\mathrm T}$ 在 $S(0;1)$ 上取到最大值 $M$ 和最小值 $m$。

\noindent (1) 当 $H_f(x_0)$ 正定时, 对任意非零向量 $x-x_0$ 均有
\[
  (x-x_0)H_f(x_0)(x-x_0)^{\mathrm T}>0,
\]
因此有 $m>0$, 从而存在 $0<\delta_1<\delta_0$, 使得当 $x\in U_0(x_0,\delta_1)$ 时, 有
\[
\begin{aligned}
  f(x)
  &=f(x_0)+\frac12
    \frac{x-x_0}{|x-x_0|}
    H_f(x_0)
    \frac{(x-x_0)^{\mathrm T}}{|x-x_0|}
    |x-x_0|^2+o(|x-x_0|^2)\\
  &>f(x_0)+\frac m4|x-x_0|^2>f(x_0).
\end{aligned}
\]
这就证明了 $f(x)$ 在 $x_0$ 处取极小值。

\noindent (2) 当 $H_f(x_0)$ 负定时, 有 $M<0$. 因此, 存在 $0<\delta_2<\delta_0$, 当 $x\in U_0(x_0,\delta_2)$ 时, 有
\[
\begin{aligned}
  f(x)
  &=f(x_0)+\frac12
    \frac{x-x_0}{|x-x_0|}
    H_f(x_0)
    \frac{(x-x_0)^{\mathrm T}}{|x-x_0|}
    |x-x_0|^2+o(|x-x_0|^2)\\
  &<f(x_0)+\frac M4|x-x_0|^2<f(x_0).
\end{aligned}
\]
这就证明了 $f(x)$ 在 $x_0$ 取极大值。

\noindent (3) 当 $H_f(x_0)$ 不定时, 有 $m<0$, $M>0$. 任取 $x_1\in\mathbb R^n$, 使得
\[
  \frac{x_1-x_0}{|x_1-x_0|}
  H_f(x_0)
  \frac{(x_1-x_0)^{\mathrm T}}{|x_1-x_0|}
  =m,
\]
则对充分小的 $t>0$, 有
\[
  f\left(x_0+t\frac{x_1-x_0}{|x_1-x_0|}\right)
  =f(x_0)+\frac12mt^2+o(t^2)
  \le f(x_0)+\frac14mt^2<f(x_0).
\]
再取 $x_2\in\mathbb R^n$, 使得
\[
  \frac{x_2-x_0}{|x_2-x_0|}
  H_f(x_0)
  \frac{(x_2-x_0)^{\mathrm T}}{|x_2-x_0|}
  =M,
\]
则对充分小的 $t>0$, 有
\[
  f\left(x_0+t\frac{x_2-x_0}{|x_2-x_0|}\right)
  =f(x_0)+\frac12Mt^2+o(t^2)
  \ge f(x_0)+\frac14Mt^2>f(x_0).
\]
这说明 $f(x)$ 在 $x_0$ 处不取极值. 证毕。

在代数课程中, 有关于 $H_f(x_0)$ 正定、负定和不定的判定法则, 在这里我们就不再赘述了. 我们经常用到 $n=2$ 的情况, 下面给出定理 14.5.2 在 $\mathbb R^2$ 情形下的特殊形式。

\noindent\textbf{推论}\quad 设函数 $f(x,y)$ 在区域 $D\subset\mathbb R^2$ 内具有二阶连续偏导数, $(x_0,y_0)\in D$, 再设 $f'(x_0,y_0)=0$, 且记
\[
  H_f(x_0,y_0)
  =
  \left.
  \begin{pmatrix}
    \dfrac{\partial^2 f}{\partial x^2} &
    \dfrac{\partial^2 f}{\partial x\partial y}\\
    \dfrac{\partial^2 f}{\partial x\partial y} &
    \dfrac{\partial^2 f}{\partial y^2}
  \end{pmatrix}
  \right|_{(x_0,y_0)}
  \triangleq
  \begin{pmatrix}
    A & B\\
    B & C
  \end{pmatrix},
\]
则

\noindent (1) 当 $A>0$,
\[
  \begin{vmatrix}
    A & B\\
    B & C
  \end{vmatrix}
  =AC-B^2>0
\]
时, $f(x,y)$ 在 $(x_0,y_0)$ 处取极小值;

\noindent (2) 当 $A<0$,
\[
  \begin{vmatrix}
    A & B\\
    B & C
  \end{vmatrix}
  =AC-B^2>0
\]
时, $f(x,y)$ 在 $(x_0,y_0)$ 处取极大值;

\noindent (3) 当
\[
  \begin{vmatrix}
    A & B\\
    B & C
  \end{vmatrix}
  =AC-B^2<0
\]
时, $f(x,y)$ 在 $(x_0,y_0)$ 处不取极值。

\noindent\textbf{例 14.5.1}\quad 试求函数 $f(x,y)=x^2+3xy+3y^2-6x-3y$ 的极值。

\noindent\textbf{解}\quad 由
\[
  f_x'(x,y)=2x+3y-6=0,\qquad
  f_y'(x,y)=3x+6y-3=0
\]
得唯一驻点 $(9,-4)$, 再由于
\[
  f_{xx}''(x,y)=2,\qquad f_{xy}''(x,y)=3,\qquad f_{yy}''(x,y)=6,
\]
从而
\[
  f_{xx}''(9,-4)=2>0,
\]
\[
  \begin{vmatrix}
    f_{xx}''(9,-4) & f_{xy}''(9,-4)\\
    f_{xy}''(9,-4) & f_{yy}''(9,-4)
  \end{vmatrix}
  =
  \begin{vmatrix}
    2 & 3\\
    3 & 6
  \end{vmatrix}
  =3>0,
\]
因此 $f(x,y)$ 在 $(9,-4)$ 处取极小值 $-21$. 因为 $f(x,y)$ 在 $\mathbb R^2$ 处处可微, 所以 $(9,-4)$ 是 $f(x,y)$ 的唯一极值点。

\noindent\textbf{例 14.5.2}\quad 某工厂生产一批长方体无盖盒子, 要求其体积为 $1\mathrm{m}^3$, 盒子的底的厚度是侧面厚度的三倍. 问: 如何设定盒子的长、宽、高才能使得用料最省?

\noindent\textbf{解}\quad 设盒子的长、宽、高分别为 $x,y,z$, 并令
\[
  F(x,y,z)=3xy+2yz+2xz.
\]
由于盒子的体积为 $1\mathrm{m}^3$, 即 $xyz=1$, 因此, 若令
\[
  S(x,y)=3xy+2\left(\frac1x+\frac1y\right),
\]
则本题的最值问题转化成求 $S(x,y)$ 的最小值问题. 由
\[
  \frac{\partial S(x,y)}{\partial x}=3y-\frac2{x^2}=0,\qquad
  \frac{\partial S(x,y)}{\partial y}=3x-\frac2{y^2}=0
\]
解得
\[
  x=y=\sqrt[3]{\frac23}.
\]
由于该实际问题总是存在最小值, 且 $\left(\sqrt[3]{\dfrac23},\sqrt[3]{\dfrac23}\right)$ 是 $S(x,y)$ 唯一的驻点, 因此它是 $S(x,y)$ 的极小值点, 且 $S(x,y)$ 必在该点取最小值. 所以当
\[
  x=y=\sqrt[3]{\frac23},\qquad z=\sqrt[3]{\frac94}
\]
时可使用料最省。

我们也可以由 $S(x,y)$ 的海色矩阵的正定性来判断该点为极小值点. 事实上, 在 $\left(\sqrt[3]{\dfrac23},\sqrt[3]{\dfrac23}\right)$ 处, $S(x,y)$ 的海色矩阵为
\[
  H_S\left(\sqrt[3]{\frac23},\sqrt[3]{\frac23}\right)
  =
  \begin{pmatrix}
    6 & 3\\
    3 & 6
  \end{pmatrix}.
\]
容易看出它是正定的, 从而 $\left(\sqrt[3]{\dfrac23},\sqrt[3]{\dfrac23}\right)$ 为 $S(x,y)$ 的极小值点。

\noindent\textbf{例 14.5.3}\quad 设 $(x_k,y_k)(k=1,2,\cdots,K)$ 为 $\mathbb R^2$ 内 $K$ 个两两不同的点, 且它们不同在 $Oxy$ 平面内的一条垂直于 $x$ 轴的直线上. 证明: 存在唯一直线 $L_0:y=ax+b$, 使得函数
\[
  f(s,t)=\sum_{k=1}^K(sx_k+t-y_k)^2
\]
在 $(a,b)$ 处达到最小。

\noindent\textbf{证明}\quad 显然 $f(s,t)=\sum\limits_{k=1}^K(sx_k+t-y_k)^2$ 在 $\mathbb R^2$ 中处处可微. 由
\[
  \begin{cases}
    \dfrac{\partial f}{\partial s}
    =2\displaystyle\sum_{k=1}^K x_k(sx_k+t-y_k)=0,\\[1.2ex]
    \dfrac{\partial f}{\partial t}
    =2\displaystyle\sum_{k=1}^K (sx_k+t-y_k)=0
  \end{cases}
\]
得
\[
  \begin{cases}
    s\displaystyle\sum_{k=1}^K x_k^2
    +t\displaystyle\sum_{k=1}^K x_k
    =\displaystyle\sum_{k=1}^K x_ky_k,\\[1.2ex]
    s\displaystyle\sum_{k=1}^K x_k
    +Kt
    =\displaystyle\sum_{k=1}^K y_k.
  \end{cases}
\]
由此求出上述方程组的解为
\[
  s=
  \frac{
    K\displaystyle\sum_{k=1}^K x_ky_k
    -\left(\displaystyle\sum_{k=1}^K x_k\right)
     \left(\displaystyle\sum_{k=1}^K y_k\right)}
  {
    K\displaystyle\sum_{k=1}^K x_k^2
    -\left(\displaystyle\sum_{k=1}^K x_k\right)^2}
  \triangleq a,
\]
\[
  t=
  \frac{
    \left(\displaystyle\sum_{k=1}^K x_k^2\right)
    \left(\displaystyle\sum_{k=1}^K y_k\right)
    -\left(\displaystyle\sum_{k=1}^K x_k\right)
     \left(\displaystyle\sum_{k=1}^K x_ky_k\right)}
  {
    K\displaystyle\sum_{k=1}^K x_k^2
    -\left(\displaystyle\sum_{k=1}^K x_k\right)^2}
  \triangleq b.
\]
由于该问题中的最小值总是存在的, 且 $f(s,t)$ 只有唯一的驻点 $(a,b)$, 因此 $(a,b)$ 必为 $f(s,t)$ 的最小值点。

在上述例子中, 已知平面 $\mathbb R^2$ 内 $K$ 个不落在一条垂直于 $x$ 轴的直线上的点, 求满足定理条件的唯一直线 $y=ax+b$ 的方法称为是最小二乘法。

\subsection{条件极值问题}

为了导出条件极值问题, 我们先来看一个例子. 设 $\Gamma\subset\mathbb R^3$ 是一条曲线, 现在要求它到原点的最短距离. 为了求解此问题, 我们可令 $f(x,y,z)=\sqrt{x^2+y^2+z^2}$, 该函数表示了 $\mathbb R^3$ 中任一点到原点的距离. 于是该问题转化为: 当 $(x,y,z)$ 在 $\Gamma$ 上变化时, 求 $f(x,y,z)$ 的最小值. 实际上, 这是 $(x,y,z)$ 在 $\Gamma$ 上变化的条件下, 求 $f(x,y,z)$ 的极小值问题. 像这样具有外加约束条件的极值问题称为条件极值问题. 为了求解这一条件极值问题, 一个自然的想法是: 求出 $\Gamma$ 的参数方程然后代入 $f(x,y,z)$, 从而将此问题转化成上一小节讨论的极值问题. 但在很多情形下, 这绝非易事. 因此我们需用其他方法来求解条件极值问题。

首先, 设 $\Gamma$ 是一条平面曲线, 它由方程 $\varphi(x,y)=0$ 所确定. 我们假定 $\varphi(x,y)$ 在区域 $D\subset\mathbb R^2$ 内具有连续偏导数, 且对于 $\forall(x,y)\in D$,
\[
  \left.
  \left(\frac{\partial\varphi}{\partial x},\frac{\partial\varphi}{\partial y}\right)
  \right|_{(x,y)}
\]
的秩为 $1$, 再设 $(x_0,y_0)\in\Gamma$ 且是函数 $f(x,y)=\sqrt{x^2+y^2}$ 的一个极值点. 在理论上 (由隐函数存在定理) 我们可以从 $\varphi(x,y)=0$ 在 $(x_0,y_0)$ 附近解出 $y=y(x)$ 或 $x=x(y)$. 现在不妨设
\[
  \frac{\partial\varphi(x_0,y_0)}{\partial y}\ne0,
\]
从而 $\exists\delta>0$, 使得在 $U(x_0,\delta)$ 内存在 $y=y(x)$, 满足 $y_0=y(x_0)$ 且 $\varphi(x,y(x))=0$. 将 $y=y(x)$ 代入 $f(x,y)$ 得 $f(x,y(x))$, 则 $x_0$ 是 $f(x,y(x))$ 的一个通常极值点, 从而有
\[
  \frac{\partial f(x_0,y_0)}{\partial x}
  +
  \frac{\partial f(x_0,y_0)}{\partial y}y'(x_0)=0.
\]
由
\[
  y'(x_0)=-\frac{\varphi_x'(x_0,y_0)}{\varphi_y'(x_0,y_0)}
\]
知
\[
  \frac{\dfrac{\partial f(x_0,y_0)}{\partial x}}
       {\dfrac{\partial\varphi(x_0,y_0)}{\partial x}}
  =
  \frac{\dfrac{\partial f(x_0,y_0)}{\partial y}}
       {\dfrac{\partial\varphi(x_0,y_0)}{\partial y}},
\]
即 $\exists\lambda\in\mathbb R$, 使得
\[
  \begin{cases}
    \dfrac{\partial f(x_0,y_0)}{\partial x}
    +\lambda\dfrac{\partial\varphi(x_0,y_0)}{\partial x}=0,\\[1.2ex]
    \dfrac{\partial f(x_0,y_0)}{\partial y}
    +\lambda\dfrac{\partial\varphi(x_0,y_0)}{\partial y}=0.
  \end{cases}
\]
再设 $\Gamma$ 是由下述方程组确定的一条空间曲线:
\[
  \begin{cases}
    \varphi_1(x,y,z)=0,\\
    \varphi_2(x,y,z)=0,
  \end{cases}
\]
并假定 $\varphi_1,\varphi_2$ 在区域 $D\subset\mathbb R^3$ 内具有连续偏导数, 且对于 $\forall(x,y,z)\in D$,
\[
  \left.
  \begin{pmatrix}
    \dfrac{\partial\varphi_1}{\partial x} &
    \dfrac{\partial\varphi_1}{\partial y} &
    \dfrac{\partial\varphi_1}{\partial z}\\
    \dfrac{\partial\varphi_2}{\partial x} &
    \dfrac{\partial\varphi_2}{\partial y} &
    \dfrac{\partial\varphi_2}{\partial z}
  \end{pmatrix}
  \right|_{(x,y,z)}
\]
的秩为 $2$. 由此条件, 理论上从方程组我们仍然可解出曲线的参数方程, 我们将其代入 $f(x,y,z)$ 就可以推出在 $f(x,y,z)$ 的极值点 $(x_0,y_0,z_0)$ 处, 存在常数 $\lambda_1,\lambda_2$ 成立下述方程组:
\[
  \begin{cases}
    \dfrac{\partial f(x_0,y_0,z_0)}{\partial x}
    +\lambda_1\dfrac{\partial\varphi_1(x_0,y_0,z_0)}{\partial x}
    +\lambda_2\dfrac{\partial\varphi_2(x_0,y_0,z_0)}{\partial x}=0,\\[1.2ex]
    \dfrac{\partial f(x_0,y_0,z_0)}{\partial y}
    +\lambda_1\dfrac{\partial\varphi_1(x_0,y_0,z_0)}{\partial y}
    +\lambda_2\dfrac{\partial\varphi_2(x_0,y_0,z_0)}{\partial y}=0,\\[1.2ex]
    \dfrac{\partial f(x_0,y_0,z_0)}{\partial z}
    +\lambda_1\dfrac{\partial\varphi_1(x_0,y_0,z_0)}{\partial z}
    +\lambda_2\dfrac{\partial\varphi_2(x_0,y_0,z_0)}{\partial z}=0.
  \end{cases}
\]
对于一般情形, 我们可以证明下面的定理。

\noindent\textbf{定理 14.5.3}\quad 设函数 $f(x)$, $\varphi(x)=(\varphi_1(x),\varphi_2(x),\cdots,\varphi_m(x))$ 在区域 $D\subset\mathbb R^n(m<n)$ 内具有各个连续偏导数, 再设 $x_0=(x_1^0,x_2^0,\cdots,x_n^0)\in D$ 为 $f(x)$ 在约束条件
\begin{equation}
  \begin{cases}
    \varphi_1(x)=0,\\
    \varphi_2(x)=0,\\
    \cdots\\
    \varphi_m(x)=0
  \end{cases}
  \tag{14.5.1}
\end{equation}
下的极值点, 并且 $\varphi'(x_0)$ 的秩为 $m$, 则存在常数 $\lambda_1,\lambda_2,\cdots,\lambda_m\in\mathbb R$, 使得在 $x_0$ 处成立下述等式:
\begin{equation}
  \begin{cases}
    \dfrac{\partial f(x_0)}{\partial x_i}
    +\displaystyle\sum_{j=1}^m\lambda_j
    \dfrac{\partial\varphi_j(x_0)}{\partial x_i}=0,
    \qquad i=1,2,\cdots,n,\\[1.2ex]
    \varphi_j(x_0)=0,\qquad j=1,2,\cdots,m.
  \end{cases}
  \tag{14.5.2}
\end{equation}

\noindent\textbf{证明}\quad 由于 $\varphi'(x_0)$ 的秩为 $m$, 我们不妨设行列式
\[
  \left.
  \frac{\partial(\varphi_1,\varphi_2,\cdots,\varphi_m)}
       {\partial(x_{n-m+1},x_{n-m+2},\cdots,x_n)}
  \right|_{x_0}
  \ne0.
\]
因此, 在 $x_0$ 的某个邻域内唯一确定一组具有各个连续偏导数的隐函数
\begin{equation}
  \begin{cases}
    x_{n-m+1}=g_1(x_1,x_2,\cdots,x_{n-m}),\\
    x_{n-m+2}=g_2(x_1,x_2,\cdots,x_{n-m}),\\
    \cdots\\
    x_n=g_m(x_1,x_2,\cdots,x_{n-m}),
  \end{cases}
  \tag{14.5.3}
\end{equation}
满足
\[
  x_j^0=g_j(x_1^0,x_2^0,\cdots,x_{n-m}^0)
  \qquad (j=n-m+1,n-m+2,\cdots,n)
\]
及对 $k=1,2,\cdots,m$, 有
\[
  \varphi_k(x_1,\cdots,x_{n-m},g_1(x_1,x_2,\cdots,x_{n-m}),\cdots,g_m(x_1,x_2,\cdots,x_{n-m}))=0.
\]
将隐函数组 (14.5.3) 代入 $f(x)$, 得
\[
  f(x_1,\cdots,x_{n-m},g_1(x_1,x_2,\cdots,x_{n-m}),\cdots,g_m(x_1,x_2,\cdots,x_{n-m})).
\]
因此, $x_0$ 是条件极值点转化为 $(x_1^0,x_2^0,\cdots,x_{n-m}^0)$ 为上述函数的通常极值点。

令 $x_0'=(x_1^0,x_2^0,\cdots,x_{n-m}^0)$, 则对 $i=1,2,\cdots,n-m$, 有
\[
  \frac{\partial f(x_0)}{\partial x_i}
  +\frac{\partial f(x_0)}{\partial x_{n-m+1}}\cdot
  \frac{\partial g_1(x_0')}{\partial x_i}
  +\cdots+
  \frac{\partial f(x_0)}{\partial x_n}\cdot
  \frac{\partial g_m(x_0')}{\partial x_i}
  =0.
\]
令 $g(x')=(g_1(x'),g_2(x'),\cdots,g_m(x'))^{\mathrm T}$, 其中 $x'=(x_1,x_2,\cdots,x_{n-m})$. 将上述 $n-m$ 个等式写成向量形式, 有
\begin{equation}
  \left(
    \frac{\partial f(x_0)}{\partial x_1},
    \cdots,
    \frac{\partial f(x_0)}{\partial x_{n-m}}
  \right)
  +
  \left(
    \frac{\partial f(x_0)}{\partial x_{n-m+1}},
    \cdots,
    \frac{\partial f(x_0)}{\partial x_n}
  \right)
  g'(x_0')=0.
  \tag{14.5.4}
\end{equation}
由于
\begin{equation}
\begin{aligned}
  g'(x_0')
  &=
  -
  \begin{pmatrix}
    \dfrac{\partial\varphi_1(x_0)}{\partial x_{n-m+1}} &
    \dfrac{\partial\varphi_1(x_0)}{\partial x_{n-m+2}} &
    \cdots &
    \dfrac{\partial\varphi_1(x_0)}{\partial x_n}\\
    \dfrac{\partial\varphi_2(x_0)}{\partial x_{n-m+1}} &
    \dfrac{\partial\varphi_2(x_0)}{\partial x_{n-m+2}} &
    \cdots &
    \dfrac{\partial\varphi_2(x_0)}{\partial x_n}\\
    \vdots & \vdots & & \vdots\\
    \dfrac{\partial\varphi_m(x_0)}{\partial x_{n-m+1}} &
    \dfrac{\partial\varphi_m(x_0)}{\partial x_{n-m+2}} &
    \cdots &
    \dfrac{\partial\varphi_m(x_0)}{\partial x_n}
  \end{pmatrix}^{-1}\\
  &\quad\cdot
  \begin{pmatrix}
    \dfrac{\partial\varphi_1(x_0)}{\partial x_1} &
    \dfrac{\partial\varphi_1(x_0)}{\partial x_2} &
    \cdots &
    \dfrac{\partial\varphi_1(x_0)}{\partial x_{n-m}}\\
    \dfrac{\partial\varphi_2(x_0)}{\partial x_1} &
    \dfrac{\partial\varphi_2(x_0)}{\partial x_2} &
    \cdots &
    \dfrac{\partial\varphi_2(x_0)}{\partial x_{n-m}}\\
    \vdots & \vdots & & \vdots\\
    \dfrac{\partial\varphi_m(x_0)}{\partial x_1} &
    \dfrac{\partial\varphi_m(x_0)}{\partial x_2} &
    \cdots &
    \dfrac{\partial\varphi_m(x_0)}{\partial x_{n-m}}
  \end{pmatrix}
  \triangleq -A^{-1}B.
\end{aligned}
\tag{14.5.5}
\end{equation}
注意到
\[
  -
  \left(
    \frac{\partial f(x_0)}{\partial x_{n-m+1}},
    \frac{\partial f(x_0)}{\partial x_{n-m+2}},
    \cdots,
    \frac{\partial f(x_0)}{\partial x_n}
  \right)\cdot A^{-1}
\]
是一个 $m$ 维行向量, 我们可以将其记为
\begin{equation}
  -
  \left(
    \frac{\partial f(x_0)}{\partial x_{n-m+1}},
    \frac{\partial f(x_0)}{\partial x_{n-m+2}},
    \cdots,
    \frac{\partial f(x_0)}{\partial x_n}
  \right)\cdot A^{-1}
  =(\lambda_1,\lambda_2,\cdots,\lambda_m).
  \tag{14.5.6}
\end{equation}
将式 (14.5.5) 与 (14.5.6) 代入式 (14.5.4) 得
\begin{equation}
\begin{aligned}
  &\left(
    \frac{\partial f(x_0)}{\partial x_1},
    \frac{\partial f(x_0)}{\partial x_2},
    \cdots,
    \frac{\partial f(x_0)}{\partial x_{n-m}}
  \right)\\
  &\quad+
  (\lambda_1,\lambda_2,\cdots,\lambda_m)
  \begin{pmatrix}
    \dfrac{\partial\varphi_1(x_0)}{\partial x_1} &
    \dfrac{\partial\varphi_1(x_0)}{\partial x_2} &
    \cdots &
    \dfrac{\partial\varphi_1(x_0)}{\partial x_{n-m}}\\
    \dfrac{\partial\varphi_2(x_0)}{\partial x_1} &
    \dfrac{\partial\varphi_2(x_0)}{\partial x_2} &
    \cdots &
    \dfrac{\partial\varphi_2(x_0)}{\partial x_{n-m}}\\
    \vdots & \vdots & & \vdots\\
    \dfrac{\partial\varphi_m(x_0)}{\partial x_1} &
    \dfrac{\partial\varphi_m(x_0)}{\partial x_2} &
    \cdots &
    \dfrac{\partial\varphi_m(x_0)}{\partial x_{n-m}}
  \end{pmatrix}
  =0.
\end{aligned}
\tag{14.5.7}
\end{equation}
另外, 我们可以将式 (14.5.6) 改写成下述形式
\begin{equation}
  \left(
    \frac{\partial f(x_0)}{\partial x_{n-m+1}},
    \frac{\partial f(x_0)}{\partial x_{n-m+2}},
    \cdots,
    \frac{\partial f(x_0)}{\partial x_n}
  \right)
  +
  (\lambda_1,\lambda_2,\cdots,\lambda_m)
  A=0.
  \tag{14.5.8}
\end{equation}
最后将向量方程 (14.5.7) 与 (14.5.8) 写成分量形式的方程, 再加上约束条件 (14.5.1), 即得式 (14.5.2). 证毕。

\noindent\textbf{注}\quad 若构造函数
\[
  F(x_1,\cdots,x_n,\lambda_1,\cdots,\lambda_m)
  =
  f(x)+\sum_{j=1}^m\lambda_j\varphi_j(x),
\]
则上述条件极值的必要条件形式上化为 $F$ 的通常极值的必要条件
\begin{equation}
  \begin{cases}
    \dfrac{\partial F(x_0)}{\partial x_i}=0
    \qquad (i=1,2,\cdots,n),\\[1.2ex]
    \dfrac{\partial F(x_0)}{\partial\lambda_j}=0
    \qquad (j=1,2,\cdots,m).
  \end{cases}
  \tag{14.5.9}
\end{equation}
上述求条件极值点的必要条件的方法称为拉格朗日乘数法。

由于在拉格朗日乘数法中, 条件极值问题转化为
\[
  F(x)=f(x)+\sum_{j=1}^m\lambda_j\varphi_j(x)
\]
的通常极值问题, 而当 $x_0$ 是驻点时, $\lambda_j(j=1,2,\cdots,m)$ 是 $m$ 个常数, 因此我们利用 $F(x)$ 在 $x_0$ 处的海色矩阵 $H_F(x_0)$ 来判定 $x_0$ 是否为极值点. 可以证明: 当 $H_F(x_0)$ 正定时, 条件极值点为极小值点; 当 $H_F(x_0)$ 负定时, 条件极值为极大值点。

值得提出的是, 当 $H_F(x_0)$ 不定时, 条件极值仍可以取到. 例如, 设函数 $f(x,y)=x^2-y^2$, $\varphi(x,y)=y=0$, 则 $\varphi(x,y)=0$ 确定直线方程 $y=0$, 显然 $f(x,0)=x^2$ 在 $(0,0)$ 处取极小值 $0$. 容易算出
\[
  F(x,y)=x^2-y^2+\lambda y
\]
的海色矩阵在 $(0,0)$ 处为
\[
  \begin{pmatrix}
    2 & 0\\
    0 & -2
  \end{pmatrix},
\]
它是一个不定矩阵。

\noindent\textbf{例 14.5.4}\quad 求椭球面 $16x^2+4y^2+9z^2=144$ 内接且各面均平行于坐标平面的长方体的最大体积。

\noindent\textbf{解}\quad 设 $(x,y,z)$ 是该椭球面内接长方体在第一卦限内的顶点, 由对称性, 该长方体的体积为
\[
  V=f(x,y,z)=8xyz.
\]
因此所求的最大体积转化为求 $f(x,y,z)$ 在约束条件
\[
  g(x,y,z)=16x^2+4y^2+9z^2-144=0
\]
下的最大值. 由拉格朗日乘数法, 我们作函数
\[
  F(x,y,z,\lambda)=f(x,y,z)+\lambda g(x,y,z).
\]
求 $F(x,y,z,\lambda)$ 的各个偏导数得方程组
\[
  \begin{cases}
    8yz+32\lambda x=0,\\
    8xz+8\lambda y=0,\\
    8xy+18\lambda z=0,\\
    16x^2+4y^2+9z^2-144=0.
  \end{cases}
  \tag{14.5.10--14.5.13}
\]
由上面方程组易推出
\begin{equation}
  xyz=-12\lambda.
  \tag{14.5.14}
\end{equation}
将式 (14.5.14) 代入式 (14.5.10) 得
\[
  -32\lambda(3-x^2)=0.
\]
因此得 $\lambda=0$ 或 $x=\sqrt3$. 当 $\lambda=0$ 时得 $xyz=0$. 显然不合题意. 因此必有 $\lambda\ne0$, 从而有 $x=\sqrt3$。

同理将式 (14.5.14) 分别代入式 (14.5.11) 和 (14.5.12) 可得 $y=2\sqrt3$ 和 $z=4\sqrt3/3$。

由于在该问题中最大值总是存在的, 因此所求最大值为
\[
  V_{\max}=8\times\sqrt3\times2\sqrt3\times\frac{4\sqrt3}{3}=64\sqrt3.
\]

\noindent\textbf{例 14.5.5}\quad 设 $n$ 个正数之和为 $c$, 试求它们乘积的最大值, 并证明对任何正数 $a_1,a_2,\cdots,a_n$, 有
\[
  \sqrt[n]{a_1a_2\cdots a_n}
  \le
  \frac{a_1+a_2+\cdots+a_n}{n}.
\]

\noindent\textbf{解}\quad 在 $E=\{(x_1,x_2,\cdots,x_n):x_i\ge0,\ i=1,2,\cdots,n\}$ 上定义函数
\[
  f(x_1,x_2,\cdots,x_n)=\prod_{i=1}^n x_i.
\]
依题意, 我们要求 $f(x_1,x_2,\cdots,x_n)$ 在约束条件
\[
  \sum_{i=1}^n x_i=c
\]
下的最大值。

作辅助函数
\[
  F(x_1,x_2,\cdots,x_n,\lambda)
  =
  \prod_{i=1}^n x_i
  +\lambda\left(\sum_{i=1}^n x_i-c\right).
\]
求 $F$ 的 $n$ 个偏导数及 $\dfrac{\partial F}{\partial\lambda}$, 并令其为零, 得方程组
\[
  \begin{cases}
    x_2x_3\cdots x_n+\lambda=0,\\
    x_1x_3\cdots x_n+\lambda=0,\\
    \cdots\\
    x_1x_2\cdots x_{n-1}+\lambda=0,\\
    \displaystyle\sum_{i=1}^n x_i=c,
  \end{cases}
\]
解出
\[
  x_1=x_2=\cdots=x_n=\frac cn,\qquad
  \lambda=-\left(\frac cn\right)^{n-1}.
\]
我们知道 $f(x_1,x_2,\cdots,x_n)$ 在紧集 $E$ 上取到最大值. 另外, 容易看出, 若 $(x_1^0,x_2^0,\cdots,x_n^0)$ 是 $f(x_1,x_2,\cdots,x_n)$ 的最大值点, 必有 $x_i^0\ne0(i=1,2,\cdots,n)$. 由于
\[
  \left(\frac cn,\frac cn,\cdots,\frac cn,-\left(\frac cn\right)^{n-1}\right)
\]
是 $F(x_1,x_2,\cdots,x_n,\lambda)$ 的唯一驻点, 因此它为 $F(x_1,x_2,\cdots,x_n,\lambda)$ 的极大值点, 从而 $f(x_1,x_2,\cdots,x_n)$ 在该点取最大值
\[
  \frac{c^n}{n^n}.
\]
现对 $n$ 个正数 $a_1,a_2,\cdots,a_n$, 令 $c=\displaystyle\sum_{i=1}^n a_i$, 并利用上述函数 $f(x)$ 所证结论, 则对任何满足 $\displaystyle\sum_{i=1}^n x_i=c$ 的正数 $x_i$ 均有
\[
  \prod_{i=1}^n x_i\le\frac{c^n}{n^n}.
\]
取 $x_i=a_i(i=1,2,\cdots,n)$, 则有
\[
  \sqrt[n]{a_1a_2\cdots a_n}
  \le
  \frac{a_1+a_2+\cdots+a_n}{n}.
\]

\section{多元微分学的几何应用}

在本章的最后一节, 我们来讨论多元微分学在几何上的一些应用。

\subsection{空间曲线的切线与法平面}

我们回忆一下, 在上一章曾经给出的曲线定义: $\mathbb R^n$ 中的一条曲线是 $[\alpha,\beta]\to\mathbb R^n$ 的一个连续映射
\[
  h(t)=(x_1(t),x_2(t),\cdots,x_n(t)).
\]
在 $n=1,2,3$ 时, 曲线是我们经常遇到的. 但是值得指出的是, 曲线有时可能是非常复杂的. 在第二册中, 我们曾经遇到过不可求长的曲线. 皮亚诺甚至构造了一条连续曲线 $\to D=[0,1]\times[0,1]$:
\[
  \begin{cases}
    x=x(t),\\
    y=y(t),
  \end{cases}
  \quad (0\le t\le1);
\]
使得它充满了整个 $D$。

在以后各章节讨论的曲线中, 我们作如下约定: 设曲线 $\Gamma$ 由连续映射
\[
  h(t)=(x_1(t),x_2(t),\cdots,x_n(t)),\quad t\in[\alpha,\beta]
\]
所确定. 若对于 $\forall t_1,t_2\in[\alpha,\beta]$, $h(t_1)\ne h(t_2)$, 则称 $\Gamma$ 是一条简单曲线. 若对于 $\forall t_1,t_2\in[\alpha,\beta)$, 有 $h(t_1)\ne h(t_2)$, 但 $h(\alpha)=h(\beta)$, 则 $\Gamma$ 是一条封闭曲线, 则称 $\Gamma$ 是一条简单闭曲线. 简单闭曲线也称为约当 (Jordan) 曲线。

我们首先讨论比较简单的情形, 即空间曲线是由参数方程
\[
  \Gamma:
  \begin{cases}
    x=x(t),\\
    y=y(t),\\
    z=z(t),
  \end{cases}
  \quad t\in(\alpha,\beta)
\]
给出, 并假定 $x(t),y(t)$ 与 $z(t)$ 都是 $t$ 的可微函数, 且 $x'^2(t)+y'^2(t)+z'^2(t)\ne0$. 对于这类曲线, 它局部总是简单曲线. 在上述条件下, 我们现在来求过 $t_0\in(\alpha,\beta)$ 对应的曲线上的点 $x(t_0)=(x(t_0),y(t_0),z(t_0))$ 处的切线方程。

由切线的定义, 它是曲线 $\Gamma$ 上过点 $x(t)=(x(t),y(t),z(t))$ 与 $x(t_0)$ 的割线当 $t\to t_0$ 时的极限位置. 由于过点 $x(t_0)$ 与 $x(t)$ 的割线方程为
\[
  \frac{x-x(t_0)}{x(t)-x(t_0)}
  =
  \frac{y-y(t_0)}{y(t)-y(t_0)}
  =
  \frac{z-z(t_0)}{z(t)-z(t_0)};
\]
在上式分母中同除以 $t-t_0$, 并令 $t\to t_0$ 即得所求切线方程:
\[
  \frac{x-x(t_0)}{x'(t_0)}
  =
  \frac{y-y(t_0)}{y'(t_0)}
  =
  \frac{z-z(t_0)}{z'(t_0)}.
\]
从上述方程可以看出, 若记曲线 $h(t)=(x(t),y(t),z(t))(t\in(\alpha,\beta))$, 则
\[
  h'(t_0)=(x'(t_0),y'(t_0),z'(t_0))
\]
即为曲线 $\Gamma$ 在 $x(t_0)=(x(t_0),y(t_0),z(t_0))$ 处的切向量。

我们同时也可得到曲线 $\Gamma$ 在 $x(t_0)=(x(t_0),y(t_0),z(t_0))$ 处的法平面方程, 即过点 $x(t_0)$ 且以切向量 $h'(t_0)$ 为法向的平面方程:
\[
  x'(t_0)(x-x(t_0))+y'(t_0)(y-y(t_0))+z'(t_0)(z-z(t_0))=0.
\]
上式用向量内积记之为
\[
  h'(t_0)\cdot(x-x(t_0))=0.
\]

其次, 我们来讨论由两个方程所确定的曲线的切线方程. 给定两个方程所组成的方程组
\[
  \begin{cases}
    F_1(x,y,z)=0,\\
    F_2(x,y,z)=0,
  \end{cases}
\]
其中 $(x,y,z)\in D\subset\mathbb R^3$, 这里 $D$ 是一个区域. 假定 $F_1$ 与 $F_2$ 均是 $D$ 上的可微函数, 并设 $x_0\in D$. 当
\[
  \begin{pmatrix}
    F_1'(x_0)\\
    F_2'(x_0)
  \end{pmatrix}
  =
  \left.
  \begin{pmatrix}
    \dfrac{\partial F_1}{\partial x} & \dfrac{\partial F_1}{\partial y} & \dfrac{\partial F_1}{\partial z}\\
    \dfrac{\partial F_2}{\partial x} & \dfrac{\partial F_2}{\partial y} & \dfrac{\partial F_2}{\partial z}
  \end{pmatrix}
  \right|_{x_0}
\]
的秩为 $2$ 时, 由定理 14.4.3, 在 $x_0=(x_0,y_0,z_0)$ 的附近, 其中有两个变量可确定为第三个变量的函数, 从而该方程组能确定一条过点 $x_0$ 的曲线 $\Gamma$. 从定理 14.4.3 我们还可以知道, 在 $x_0$ 附近 $\Gamma$ 还可以具有参数形式: $(x(t),y(t),z(t)),t\in(\alpha,\beta)$, 其中
\[
  (x(t_0),y(t_0),z(t_0))=(x_0,y_0,z_0)=x_0.
\]
由上面的讨论知, 曲线在点 $x_0$ 处的切向量为 $(x'(t_0),y'(t_0),z'(t_0))$. 由于曲线 $\Gamma$ 由方程组确定, 从而有
\[
  \begin{cases}
    F_1(x(t),y(t),z(t))=0,\\
    F_2(x(t),y(t),z(t))=0.
  \end{cases}
\]
对此方程组的第一个方程两边关于 $t$ 求导数, 得
\[
  \frac{\partial F_1(x_0)}{\partial x}x'(t_0)
  +
  \frac{\partial F_1(x_0)}{\partial y}y'(t_0)
  +
  \frac{\partial F_1(x_0)}{\partial z}z'(t_0)
  =0.
\]
将上式写成内积形式有
\[
  (x'(t_0),y'(t_0),z'(t_0))\cdot
  \left(
    \frac{\partial F_1(x_0)}{\partial x},
    \frac{\partial F_1(x_0)}{\partial y},
    \frac{\partial F_1(x_0)}{\partial z}
  \right)
  =0.
\]
同理,
\[
  (x'(t_0),y'(t_0),z'(t_0))\cdot
  \left(
    \frac{\partial F_2(x_0)}{\partial x},
    \frac{\partial F_2(x_0)}{\partial y},
    \frac{\partial F_2(x_0)}{\partial z}
  \right)
  =0.
\]
因此, 曲线 $\Gamma$ 在 $x_0$ 处的切向量与
\[
  F_1'(x_0)\times F_2'(x_0)
  =
  \left.
  \begin{vmatrix}
    \mathbf i & \mathbf j & \mathbf k\\
    \dfrac{\partial F_1}{\partial x} & \dfrac{\partial F_1}{\partial y} & \dfrac{\partial F_1}{\partial z}\\
    \dfrac{\partial F_2}{\partial x} & \dfrac{\partial F_2}{\partial y} & \dfrac{\partial F_2}{\partial z}
  \end{vmatrix}
  \right|_{x_0}
\]
平行. 记
\[
  A=\left.\frac{\partial(F_1,F_2)}{\partial(y,z)}\right|_{x_0},\quad
  B=\left.\frac{\partial(F_1,F_2)}{\partial(z,x)}\right|_{x_0},\quad
  C=\left.\frac{\partial(F_1,F_2)}{\partial(x,y)}\right|_{x_0},
\]
则 $F_1'(x_0)\times F_2'(x_0)=A\mathbf i+B\mathbf j+C\mathbf k$. 因此曲线 $\Gamma$ 在 $x_0$ 处的切线方程为
\[
  \frac{x-x_0}{A}=\frac{y-y_0}{B}=\frac{z-z_0}{C},
\]
而法平面方程为
\[
  A(x-x_0)+B(y-y_0)+C(z-z_0)=0.
\]

\subsection{曲面的切平面与法线}

设曲面 $S$ 由方程
\[
  F(x,y,z)=0
\]
给出, 其中 $F(x,y,z)\in C^1(D)$, $D\subset\mathbb R^3$ 为一个区域, 再设点 $(x_0,y_0,z_0)\in D$, 使得 $F(x_0,y_0,z_0)=0$. 现在我们来求曲面 $S$ 在点 $(x_0,y_0,z_0)$ 处的切平面及法线的方程。

在曲面 $S$ 上任取一条过点 $(x_0,y_0,z_0)$ 的光滑曲线 $(x(t),y(t),z(t))$, $t\in(\alpha,\beta)$. 因此有
\[
  F(x(t),y(t),z(t))=0,\quad \forall t\in(\alpha,\beta).
\]
设 $t_0\in(\alpha,\beta)$, 使得 $(x(t_0),y(t_0),z(t_0))=(x_0,y_0,z_0)$. 将上述方程两边在 $t_0$ 处求导数得
\[
  F_x'(x_0,y_0,z_0)x'(t_0)+F_y'(x_0,y_0,z_0)y'(t_0)+F_z'(x_0,y_0,z_0)z'(t_0)=0.
\]
注意到 $(x'(t_0),y'(t_0),z'(t_0))$ 是该曲线在 $(x_0,y_0,z_0)$ 处的切向量, 因此上式说明了该切向量与向量 $F'(x_0,y_0,z_0)$ 正交, 从而当 $F'(x_0,y_0,z_0)$ 不是零向量时, 上式表明曲面 $S$ 过点 $(x_0,y_0,z_0)$ 的任何光滑曲线的切向量与固定向量 $F'(x_0,y_0,z_0)$ 正交. 由此我们推知曲面 $S$ 过点 $(x_0,y_0,z_0)$ 的任何光滑曲线在 $(x_0,y_0,z_0)$ 处的切线都在平面
\[
  F_x'(x_0,y_0,z_0)(x-x_0)+F_y'(x_0,y_0,z_0)(y-y_0)+F_z'(x_0,y_0,z_0)(z-z_0)=0
\]
上. 我们自然地称上述平面为曲面 $S$ 在 $(x_0,y_0,z_0)$ 处的切平面。

另外, 我们称直线
\[
  \frac{x-x_0}{F_x'(x_0,y_0,z_0)}
  =
  \frac{y-y_0}{F_y'(x_0,y_0,z_0)}
  =
  \frac{z-z_0}{F_z'(x_0,y_0,z_0)}
\]
为曲面 $S$ 在 $(x_0,y_0,z_0)$ 处的法线。

当 $F'(x_0,y_0,z_0)\ne0$ 时, 我们可以证明: 曲面 $S$ 在 $(x_0,y_0,z_0)$ 处的切平面上任何过 $(x_0,y_0,z_0)$ 的直线都是该曲面上某光滑曲线的切线 (见本章习题)。

现在设曲面 $S$ 由参数方程
\[
  \begin{cases}
    x=x(u,v),\\
    y=y(u,v),\\
    z=z(u,v),
  \end{cases}
  \quad (u,v)\in D
\]
给出, 其中 $D$ 是 $\mathbb R^2$ 中的区域, 而上述三个函数均具有连续偏导数. 我们来求曲面 $S$ 在 $(x_0,y_0,z_0)=(x(u_0,v_0),y(u_0,v_0),z(u_0,v_0))$ 处的切平面与法线方程. 为此, 我们在曲面 $S$ 上取两条特殊的曲线:
\[
  \begin{cases}
    x=x(u,v_0),\\
    y=y(u,v_0),\\
    z=z(u,v_0),
  \end{cases}
  \quad\text{和}\quad
  \begin{cases}
    x=x(u_0,v),\\
    y=y(u_0,v),\\
    z=z(u_0,v).
  \end{cases}
\]
它们在 $(x_0,y_0,z_0)$ 处的切向量分别为
\[
  (x_u'(u_0,v_0),y_u'(u_0,v_0),z_u'(u_0,v_0)),
\]
\[
  (x_v'(u_0,v_0),y_v'(u_0,v_0),z_v'(u_0,v_0)).
\]
若这两个向量不平行, 它们对应的切线将确定曲面 $S$ 在 $(x_0,y_0,z_0)$ 处的法向量:
\[
  n=(x_u(u_0,v_0),y_u(u_0,v_0),z_u(u_0,v_0))
  \times(x_v(u_0,v_0),y_v(u_0,v_0),z_v(u_0,v_0))
\]
\[
  =
  \left.
  \begin{vmatrix}
    \mathbf i & \mathbf j & \mathbf k\\
    x_u' & y_u' & z_u'\\
    x_v' & y_v' & z_v'
  \end{vmatrix}
  \right|_{(u_0,v_0)}.
\]
因此当
\[
  \left.
  \begin{pmatrix}
    x_u' & y_u' & z_u'\\
    x_v' & y_v' & z_v'
  \end{pmatrix}
  \right|_{(u_0,v_0)}
\]
的秩为 $2$ 时, 记
\[
  A=\left.\frac{\partial(y,z)}{\partial(u,v)}\right|_{(u_0,v_0)},\quad
  B=\left.\frac{\partial(z,x)}{\partial(u,v)}\right|_{(u_0,v_0)},\quad
  C=\left.\frac{\partial(x,y)}{\partial(u,v)}\right|_{(u_0,v_0)},
\]
则该曲面在 $(x_0,y_0,z_0)$ 处的切平面方程与法线方程分别为
\[
  A(x-x_0)+B(y-y_0)+C(z-z_0)=0,
\]
\[
  \frac{x-x_0}{A}=\frac{y-y_0}{B}=\frac{z-z_0}{C}.
\]

到现在为止, 我们已经研究了 $\mathbb R^3$ 中的曲线切线与曲面的切平面等问题. 作为推广, 我们还可以考虑 $\mathbb R^n$ 中的曲线. 在 $\mathbb R^n$ 中, 我们曾称 $[0,1]$ 到 $\mathbb R^n$ 的一个连续映射 $h(t)=(x_1(t),x_2(t),\cdots,x_n(t))$ 为一条曲线. 当 $h(t)$ 具有连续导数且 $h'(t)\ne0$ 时, 称此曲线为光滑曲线, 并称 $h'(t_0)$ 是曲线在 $h(t_0)$ 处的切向量. 类似地, 我们称 $F(x)=F(x_1,x_2,\cdots,x_n)=0$ 为 $\mathbb R^n$ 中的一个曲面, 其中 $F(x)$ 是一个 $C^1$ 函数. 当 $F'(x_0)\ne0$ 时, 称 $F'(x_0)$ 为此曲面在 $x_0$ 处的法向量, 并称
\[
  F'(x_0)(x-x_0)
  =
  \sum_{i=1}^n \frac{\partial F(x)}{\partial x_i}(x_i-x_i^0)=0
\]
为该曲面在 $x_0$ 处的切平面方程. 作为练习, 请读者写出 $\mathbb R^n$ 中曲线的切线方程与法平面方程以及曲面的法线方程。

\noindent\textbf{例 14.6.1}\quad 求曲线
\[
  \begin{cases}
    x^2-y^2+2z^2=2,\\
    x+y+z=3
  \end{cases}
\]
在 $(1,1,1)$ 处的切线方程与法平面方程。

\noindent\textbf{解}\quad 令 $F(x,y,z)=x^2-y^2+2z^2-2$, $G(x,y,z)=x+y+z-3$, 则
\[
  \frac{\partial F}{\partial x}=2x,\quad
  \frac{\partial F}{\partial y}=-2y,\quad
  \frac{\partial F}{\partial z}=4z,\quad
  \frac{\partial G}{\partial x}
  =
  \frac{\partial G}{\partial y}
  =
  \frac{\partial G}{\partial z}
  =1,
\]
于是
\[
  \left.\frac{\partial(F,G)}{\partial(y,z)}\right|_{(1,1,1)}
  =
  \left.
  \begin{vmatrix}
    \dfrac{\partial F}{\partial y} & \dfrac{\partial F}{\partial z}\\
    \dfrac{\partial G}{\partial y} & \dfrac{\partial G}{\partial z}
  \end{vmatrix}
  \right|_{(1,1,1)}
  =
  \begin{vmatrix}
    -2 & 4\\
    1 & 1
  \end{vmatrix}
  =-6,
\]
\[
  \left.\frac{\partial(F,G)}{\partial(z,x)}\right|_{(1,1,1)}
  =
  \left.
  \begin{vmatrix}
    \dfrac{\partial F}{\partial z} & \dfrac{\partial F}{\partial x}\\
    \dfrac{\partial G}{\partial z} & \dfrac{\partial G}{\partial x}
  \end{vmatrix}
  \right|_{(1,1,1)}
  =
  \begin{vmatrix}
    4 & 2\\
    1 & 1
  \end{vmatrix}
  =2,
\]
\[
  \left.\frac{\partial(F,G)}{\partial(x,y)}\right|_{(1,1,1)}
  =
  \left.
  \begin{vmatrix}
    \dfrac{\partial F}{\partial x} & \dfrac{\partial F}{\partial y}\\
    \dfrac{\partial G}{\partial x} & \dfrac{\partial G}{\partial y}
  \end{vmatrix}
  \right|_{(1,1,1)}
  =
  \begin{vmatrix}
    2 & -2\\
    1 & 1
  \end{vmatrix}
  =4.
\]
因此, 所求切线方程为
\[
  \frac{x-1}{-3}=\frac{y-1}{1}=\frac{z-1}{2},
\]
法平面方程为
\[
  -3(x-1)+(y-1)+2(z-1)=0,\quad \text{即}\quad 3x-y-2z=0.
\]

对于上述切线方程, 还可用另一方法求解. 由于曲面 $F(x,y,z)=0$ 在 $(1,1,1)$ 处的法向量为 $(F_x',F_y',F_z')|_{(1,1,1)}=(2,-2,4)$, 因此在该点处的切平面方程为
\[
  2(x-1)-2(y-1)+4(z-1)=0,
\]
即 $x-y+2z=2$. 由于切线是两曲面在该点的切平面的交线, 而另一曲面 $G(x,y,z)=0$ 为一平面, 因此所求切线方程为
\[
  \begin{cases}
    x-y+2z=2,\\
    x+y+z=3.
  \end{cases}
\]

\noindent\textbf{例 14.6.2}\quad 证明曲面 $S:x^2+y^2+2xyz=4$ 与曲面族 $S_t:2x+ty^2-(1+t)z^2=1(t\in\mathbb R)$ 在 $(1,1,1)$ 处正交 (即它们的法向量相互垂直), 并求出它们的交线在该点的切线方程。

\noindent\textbf{解}\quad 令
\[
  F(x,y,z)=x^2+y^2+2xyz-4,
\]
\[
  G_t(x,y,z)=2x+ty^2-(1+t)z^2-1,
\]
则 $S$ 在 $(1,1,1)$ 处的法向量为
\[
  \left(
    \frac{\partial F(1,1,1)}{\partial x},
    \frac{\partial F(1,1,1)}{\partial y},
    \frac{\partial F(1,1,1)}{\partial z}
  \right)
  =(2,2,2),
\]
而 $S_t$ 在 $(1,1,1)$ 处的法向量为
\[
  \left(
    \frac{\partial G_t(1,1,1)}{\partial x},
    \frac{\partial G_t(1,1,1)}{\partial y},
    \frac{\partial G_t(1,1,1)}{\partial z}
  \right)
  =(2,2t,-2(1+t)).
\]
因为
\[
  (2,2,2)(2,2t,-2(1+t))=0,
\]
所以 $S$ 与 $S_t$ 在 $(1,1,1)$ 处总是正交的。

由于 $S$ 与 $S_t$ 的交线在 $(1,1,1)$ 处的切线方程即两曲面在 $(1,1,1)$ 处的切平面的交线, 因此切线方程为
\[
  \begin{cases}
    (x-1)+(y-1)+(z-1)=0,\\
    (x-1)+t(y-1)-(1+t)(z-1)=0,
  \end{cases}
\]
化简得
\[
  \begin{cases}
    x+y+z=3,\\
    x+ty-(1+t)z=0.
  \end{cases}
\]

\subsection{多元凸函数}

作为泰勒公式的应用, 我们来证明 $n$ 元函数凸性的一个结果. 由于研究凸性需要考虑任意两点之间线段上的函数值与端点函数值之间的关系, 因此我们假定 $D\subset\mathbb R^n$ 是一个凸域. 若 $f(x)$ 在凸域 $D$ 内具有 $k$ 阶连续偏导数, 则对于 $\forall x_0\in D$, $f(x)$ 在 $D$ 内总可以在 $x_0$ 处展成 $k$ 阶泰勒公式。

\noindent\textbf{定义 14.6.1}\quad 设 $D\subset\mathbb R^n$ 是一个凸域, $f(x)$ 在 $D$ 内有定义. 如果对于 $\forall x_0,x_1\in D$ 和 $\forall t\in(0,1)$, 有
\[
  f(tx_1+(1-t)x_2)\le tf(x_1)+(1-t)f(x_2),
\]
则称 $f(x)$ 在 $D$ 内是凸函数; 如果上述不等式总成立严格不等式, 则称 $f(x)$ 在 $D$ 内是严格凸函数。

在微分学中, 我们主要利用导数 (偏导数) 来研究凸函数. 回忆一下, 对一元函数 $f(x)$ 来说, 若 $f''(x)$ 存在, 且 $f''(x)>0$ 处处成立, 则 $f(x)$ 是凸的. 我们前面提到过, 对于 $n$ 元函数 $f(x)$, 其导数 $f'(x)$ 在很多情形下将起着一元函数导数相似的作用. 下面定理的结论说明, 在凸函数的研究中, $n$ 元二阶连续可微函数 $f(x)$ 的海色矩阵 $H_f(x)$ 起着类似于一元函数二阶导数的作用。

\noindent\textbf{定理 14.6.1}\quad 设 $D\subset\mathbb R^n$ 是一个凸区域, 函数 $f(x)$ 在 $D$ 内具有二阶连续偏导数, 则以下结论等价:

(1) $f(x)$ 在 $D$ 内是凸函数;

(2) 对于 $\forall x_0,x\in D$, 成立 $f(x)\ge f(x_0)+f'(x_0)(x-x_0)^{\mathrm T}$;

(3) 对于 $\forall x_0\in D$, $f(x)$ 在 $x_0$ 处的海色矩阵 $H_f(x_0)$ 半正定。

\noindent\textbf{证明}\quad 先证 (1) $\Rightarrow$ (2). 由凸函数的定义知, 对于 $\forall x_0,x\in D$ 及 $\forall t\in(0,1)$, 有
\begin{equation}
  f(tx+(1-t)x_0)\le tf(x)+(1-t)f(x_0).
  \tag{14.6.1}
\end{equation}
由于 $f(x)$ 具有二阶连续偏导数, 我们有
\begin{equation}
  \begin{aligned}
    f(tx+(1-t)x_0)
    &=f(x_0+t(x-x_0))\\
    &=f(x_0)+tf'(x_0)(x-x_0)^{\mathrm T}
      +o(t|x-x_0|)\quad (t|x-x_0|\to0).
  \end{aligned}
  \tag{14.6.2}
\end{equation}
将式 (14.6.2) 代入式 (14.6.1) 得
\[
  t(f'(x_0)(x-x_0)^{\mathrm T})+o(t|x-x_0|)\le t(f(x)-f(x_0)).
\]
在上式中两边除以 $t$, 并注意到 $|x-x_0|$ 在 $t$ 变化的过程中是常数, 我们有
\[
  \lim_{t\to0+0}\frac{o(t|x-x_0|)}{t}=0,
\]
因此
\[
  f(x)\ge f(x_0)+f'(x_0)(x-x_0)^{\mathrm T}.
\]

再证 (2) $\Rightarrow$ (3). 任取 $x_0\in D$ 及 $\Delta x\in\mathbb R^n$, $\Delta x\ne0$, 则当 $t$ 充分小时, $x_0+t\Delta x\in D$. 由泰勒公式有
\[
  f(x_0+t\Delta x)=f(x_0)+tf'(x_0)\Delta x^{\mathrm T}
  +\frac{t^2}{2}\Delta xH_f(x_0)\Delta x^{\mathrm T}
  +o(t^2|\Delta x|^2)\quad (t|\Delta x|\to0).
\]
在 (2) 中令 $x=x_0+t\Delta x$, 我们有
\[
  \frac{t^2}{2}\Delta xH_f(x_0)\Delta x^{\mathrm T}
  +o(t^2|\Delta x|^2)\ge0,
\]
即
\[
  \Delta xH_f(x_0)\Delta x^{\mathrm T}
  +\frac{o(t^2|\Delta x|^2)}{t^2}\ge0.
\]
令 $t\to0+0$, 即得
\[
  \Delta xH_f(x_0)\Delta x^{\mathrm T}\ge0,
\]
这说明 $H_f(x_0)$ 是半正定的。

最后证 (3) $\Rightarrow$ (1). 对于 $\forall x_1,x_2\in D$, 令 $x_0=tx_1+(1-t)x_2(t\in(0,1))$, 则有 $x_0\in D$, 从而有
\[
  f(x_1)=f(x_0)+f'(x_0)(x_1-x_0)^{\mathrm T}
  +\frac12(x_1-x_0)H_f(x_0+\theta(x_1-x_0))(x_1-x_0)^{\mathrm T},
\]
其中 $0<\theta<1$. 由 (3) 知
\[
  f(x_1)\ge f(x_0)+f'(x_0)(x_1-x_0)^{\mathrm T}.
\]
同理,
\[
  f(x_2)\ge f(x_0)+f'(x_0)(x_2-x_0)^{\mathrm T}.
\]
因此有
\[
  \begin{aligned}
    tf(x_1)+(1-t)f(x_2)
    &\ge t[f(x_0)+f'(x_0)(x_1-x_0)^{\mathrm T}]
      +(1-t)[f(x_0)+f'(x_0)(x_2-x_0)^{\mathrm T}]\\
    &=f(x_0)=f(tx_1+(1-t)x_2).
  \end{aligned}
\]
证毕。

\noindent\textbf{注}\quad 条件 (2) 说明 $u=f(x)$ 表示的 ``曲面'' 在其任一点处的切平面的上方。

\noindent\textbf{例 14.6.3}\quad 证明: 若一元凸函数 $y=g(x)$ 在 $(a,b)$ 内具有二阶导数, 则 $f(x)=\displaystyle\sum_{i=1}^n g(x_i)$ 是
\[
  D=\underbrace{(a,b)\times(a,b)\times\cdots\times(a,b)}_{n\text{ 个}}
\]
内的凸函数。

\noindent\textbf{证明}\quad 因为 $f'(x)=(g'(x_1),g'(x_2),\cdots,g'(x_n))$, 所以
\[
  H_f(x)=
  \begin{pmatrix}
    g''(x_1) & 0 & \cdots & 0\\
    0 & g''(x_2) & \cdots & 0\\
    \vdots & \vdots & & \vdots\\
    0 & 0 & \cdots & g''(x_n)
  \end{pmatrix}.
\]
由于 $g''(x_i)\ge0(i=1,2,\cdots,n)$, 从代数知识知, $H_f(x)$ 是半正定矩阵. 由定理 14.6.1 知 $f(x)$ 是 $D$ 内的凸函数。

\section*{习题十四}

1. 设函数 $u=f(x)$ 在 $U(x_0,\delta_0)\subset\mathbb R^n(\delta_0>0)$ 内存在各个偏导数, 并且所有的偏导数在该邻域内有界, 证明 $f(x)$ 在 $x_0$ 处连续; 举例说明存在函数 $u=g(x)$, 它在 $x_0$ 的某个邻域内存在无界的各个偏导数, 但它在 $x_0$ 处连续。

2. 举例说明在 $\mathbb R^2$ 内存在函数 $z=f(x,y)$, 使得 $f(x,y)$ 在 $\mathbb R^2$ 内处处不连续, 但它在原点处存在两个偏导数。

3. 求下列函数在指定点处的偏导数:

(1) $f(x,y)=xy\ln[x^2+\sin(xy^2)+\sin(xy)]$, 求 $f_x'(1,-1),f_y'(1,-1)$;

\end{document}
